\input{preamble.sty}%available at github.com/danimalabares/config
\input{thms-por.sty}%available at github.com/danimalabares/config

\renewcommand{\curvearrowright}{\mathbb{y}}
\renewcommand{\curvearrowleft}{\mathbb{x}}

%This preamble is in github.com/danimalabares/config

\begin{document}

{\Huge Geometria simpl\'etica}

\tableofcontents

\section{Aula 1}
Além do material do curso, uso bastante Lee, Intro. to Smooth Manifolds, e \href{https://www.damtp.cam.ac.uk/user/tong/dynamics.html}{Tong, Lectures on Classical Mechanics}.

\subsection{Origem da geometria simpl\'etica}
\begin{itemize}
	\item Formula\c c\~ao da geom\'etrica da mec\^anica (s\'ec XIX).
	\item Vers\~ao moderna, 1960-70.
	\item Diferentes descrip\c c\~oes da mec\^anica cl\'asica:
		\begin{itemize}
			\item Newtoniano: $F=ma$, ecua\c c\~ao diferencial ordin\'aria de segunda ordem.
			\item Lagrangiano: princ\'ipio gravitacional (Eq. E-L). Following Tong, these equations are:

			\item Hamiltoniano.
		\end{itemize}
\end{itemize}

\subsection{Formalismo hamiltoniano (simplificado)}

This happened in the 1880's (according to Tong).


\begin{itemize}
	\item Espa\c co de base $\mathbb{R}^{2}=\{(p,q)\} $ (conjunto de estados)
	\item Fun\c c\~ao Hamiltoniana $H\in C^{\infty}(\mathbb{R}^{2m} )$.
	\item Campo Hamiltoniano: $X_{H}\in \mathfrak{X}(\mathbb{R}^{2n})$.
		\begin{align*}
			X_{H}=\begin{pmatrix}  \frac{\partial H}{\partial p_{i}}\\-\frac{\partial H}{\partial q_{i}} \end{pmatrix} = \left(\begin{tabular}{c|c}0&$\operatorname{Id}_{n}$\\ \hline
		$-\operatorname{Id}_{n}$&0
\end{tabular}\right)
\end{align*}

	Which coincides with Lee's formula
\begin{align*}
	\dot x^{i}(t)&=\frac{\partial H}{\partial y^{i}}(x(t),y(t)),\\
	\dot y^{i}(t)&=-\frac{\partial H}{\partial x^{i}}(x(t),y(t))
\end{align*}
where Lee defined the \textit{\textbf{Hamiltonian vector field}} as the  \textit{analogue of the gradient with respect to the symplectic form}, that is, satisfying $\omega(X_{H},Y)=dH(Y)$ for any vector field $Y$.

Also look at Tong's formulation:
\begin{align*}
	\dot p_{i}&=-\frac{\partial H}{\partial q_{i}}\\
	\dot q_{i}&=\frac{\partial H}{\partial p_{i}}\\
	-\frac{\partial L}{\partial t}&=\frac{\partial H}{\partial t}
\end{align*}
where $L$ is the Lagrangian and the Hamiltonian function $H$ is obtained as the Legendre transform of the Langrangian. Tong shows how the Hamiltonian formalism allows to replace the $n$ $2^{\operatorname{nd}}$ order differential equations by $2n$ $1^{\operatorname{st}}$ order differential equations for $q_{i}$ and $p_{i}$.

\begin{quotation}
	In practice, for solving problems, this isn't particularly helful. But, as we shall see, conceptually it's very useful!
\end{quotation}
At least for me, it looks like a first insight on why symplectic geometry lives on even-dimensional spaces.
\end{itemize}

\subsection{Evolu\c c\~ao temporal (equa\c c\~oes de Hamilton)}
Curvas integrais
\[c(t)=(q_{i}(t),p_{i}(t))\]
de $X_{H}$, ie.
\[c'(t)=X_{H}(c(t))\iff\begin{cases}
	\dot q_{i}&=\frac{\partial H}{\partial p_{i}} \\
	\dot p_{i} &=\frac{\partial H}{\partial q_{i}}
\end{cases}\]
que s\~ao as \textit{\textbf{Equa\c c\~oes de Hamilton}} (de novo).

\begin{example}
	Part\'icula de massa $m$ em $\mathbb{R}^{3} =\{q_{1},q_{2},q_{3}\} $ sujeita a campo de for\c ca conservativa
	\[F=-\nabla V,\quad V\in C^{\infty}(\mathbb{R}^{3}\]
	\[q(t)=(q_{1},q_{2},q_{3})\]
	Equa\c c\~ao de Newton:
	\[m\ddot q=\partial V(q) \iff m\ddot q_{i}=\frac{\partial V}{\partial q_{i}}(q),\qquad i=1,2,3. \]

Ponto de vista Hamiltoniano:
\begin{itemize}
	\item Espa\c code fase $\mathbb{R}^{5}=\{(q_{i},p_{i})\} $.
	\item Hamiltoniano: $H(p,q)=\frac{1}{2m}\sum_{i}p_{i}^{2} +V(q)$ 
	\item Equa\c c\~oes de Hamilton
	 \[\begin{cases}
	 	\dot q_{i}=p_{i}/m\iff p_{i}=m\dot q_{i}\\
	 	\dot p_{i}=-\frac{\partial V}{\partial q_{i}}\qquad &
	 \end{cases}\]
\end{itemize}

\[\begin{tikzcd}
	H\in C^{\infty}(\mathbb{R}^{2n}) \arrow[r,rightsquigarrow]&\nabla H\arrow[r,rightsquigarrow ,"-J_{0}\nabla H"]&X_{H}
\end{tikzcd}\]
where $J_{0}=\begin{pmatrix}0&-\operatorname{I}\\\operatorname{I}&0\end{pmatrix} $. So it looks like another way of obtaining (defining?) the Hamiltonian vector field is to take the gradient of $H$ and then applying $J_{0}$. So it would be nice to see eventually that this is the same as Lee's definition of "symplectic gradient" so to say.
\end{example}

Compondo $\nabla H$ e $X_{H}$ : taxa de varia\c c\~ao de $H$ ao longo dos fluxos. {\color{persimmon}Mas: o que \'e a composi\c c\~ao de dois campos vetoriais? Tal vez \'e a derivada exterior de $H$, $dH$ em lugar do gradiente de  $H$.}

\begin{itemize}
	\item \textit{\textbf{Fluxo gradiente}}
		\begin{align*}c'(t)&=\nabla H(c(t))\\
			\frac{d}{dt}H(c(t))&=\left<\nabla H(c(t)),c'(t)\right> =\|\nabla H(c(t))\|^{2}
\end{align*}
$\nabla H$ aponta na dire\c c\~ao que $H$ varia\c c\~ao.
\item \textit{\textbf{Fluxo hamiltoniano}} 
	\begin{align*}
		c'(t)& =X_{H}(c(t))\\
		\frac{d}{dt}H(c(t))&=\left<\nabla H(c(t)),c'(t)\right> \\
		&=\left<\nabla H(c(t)),-J_{0}\nabla H(c(t))\right>\\
		&=0
	\end{align*}
	{\color{persimmon}?}, $H\in C^{\infty}(\mathbb{R}^{2n} )$, $H\rightsquigarrow dH\in \Omega^{1}(\mathbb{R}^{2n} )$.
	\item \textit{\textbf{Gradiente.}} $\nabla H(x)\in T_{x}\mathbb{R}^{2n} =\mathbb{R}^{2n}$ \'e \'unico.
\begin{align*}
	g_{0}(\nabla H(x),\cdot )=\left<\nabla H(x),\cdot \right> =dH(x)
\end{align*}
onde $g_{0}$ \'e a m\'etrica Euclidiana. De outra forma,
\begin{align*}
	g_{0}^{\flat}:\mathbb{R}^{2n} &\overset{\sim}{\to  } (\mathbb{R}^{2n})^{*}\\
u&\mapsto g_0(u,\cdot )
\end{align*}
assim,
\[\nabla H(x)\overset{\sim}{\to }dH(x).\]
Analogamente, $X_{H}(x)\in \mathbb{R}^{2n}$ \'e \'unico {\color{persimmon}tal que?}
\[\Omega_0(X_{H}(x),\cdot )=dH(x),\qquad \Omega_0(u,v)=-dJ_0V,\]
ou:
\begin{align*}
	\Omega_0^{\flat}: \mathbb{R}^{2n} &\overset{\sim}{\longrightarrow}(\mathbb{R}^{2n})^{*}\\
	X_{H}(x) &\longleftrightarrow dH(x) 
\end{align*}
\end{itemize}

\begin{remark}
	Note que $\Omega_{q}$ define uma 2-forma ({\color{persimmon}c…?)} em $\mathbb{R}^{2n}=\{(q_{i},p_{i})\} $.
	\[\omega_0=\sum_{i=1}^{n} dq_{i}\wedge dp_{i}\in \Omega_2(\mathbb{R}^{2n}),\]
	$X_{H}$ \'e \'unico tal que $i_{X_{H}}\omega_{0}=dH$. So this was Lee's definition $\ddot \smile$.
\end{remark}

\begin{defn}[tempor\'aria]
	Uma \textit{\textbf{variedade simpl\'etica}} \'e $(M,\omega)$, $\omega\in \Omega^{2}(M)$ localmente isomorfa a $(\mathbb{R}^{2n},\sum_{i}dq_{i}\wedge dp_{i})$.

	[Dessenho mostrando que o pullback da carta coordenada leva $\omega$ em $\sum_{i}dq_{i}\wedge dp_{i}$.

\begin{thm}[de Darboux, em Lee]\leavevmode
	Let $(M,\omega)$ be a $2n$-dimensional symplectic manifold. For any $p \in M$ there are smooth coordinates $(x^{1},\ldots,x^{n},y^{1},\ldots,y^{n})$ centered at $p$ in which $\omega$ has the coordinate representation $\omega=\sum_{i=1}^{n} dx^{i} \wedge dy^{i}$.
\end{thm}

And Lee does a proof using the \textit{theory of time-dependant flows}.

	\[\begin{tikzcd}
		\substack{\text{mec\^anica}  \\ \text{cl\'asica} }\arrow[r,"\text{qu\^antica} "] \arrow[dr]&\substack{\text{Teoria de rep.}  \\ \text{de Lie} }\arrow[d]&\text{geo. Riemanniana} \arrow[dl]\\
		&\text{geo. simpl\'etica}\arrow[d]\arrow[u]\arrow[dl,bend right]\\
		\text{geo. K\"ahler} &\substack{\text{din\'amica}  \\ \text{{\color{persimmon}Lagrangiana?}} } 
	\end{tikzcd}\]
	
\end{defn}

\subsection{\'Algebra linear simpl\'etica}

$V$ espa\c co vetorial real, $\Omega:V\times V\to \mathbb{R}$ forma bilinea ansim\'etrica, i.e. $\Omega\in \Lambda^{2} V^{*}$.

\begin{defn}
	$\Omega$ \'e n\~ao degenerada se $\Omega(u,v)=0\forall v\iff u=0$.

	Following Lee, this can also be stated as: for each nonzero $v\in V$ there exists $w\in V$ such that $\omega(v,w)\neq 0$; and it is equivalent to the linear map $v\mapsto \omega(v,\cdot )\in V^{*}$ being invertible, and also that in terms of some (hence every) basis, the matrix $(\omega_{ij})$ representing $\omega$ is nonsingular.

Ou seja, se
\[\ker \Omega:=\{u\in V|\Omega(u,v)=0\;\forall v\} \]
ent\~ao $\Omega$ \'e n\~ao degenerada se e somente se $\ker (\Omega)=\{0\} $.

$\Omega\in \Lambda^{2} V^{*}$ \'e n\~ao degenerada \'e chamada simpl\'etica. $(V,\Omega)$ \'e um \textit{\textbf{espa\c co vectorial simpl\'etico}}.
\end{defn}

\begin{remark}\leavevmode 
	\begin{enumerate}
		\item $\{e_1,..,e_{n}\} $ base de $V$, $\Omega$ \'e representado por uma matriz antisim\'etrica
	\[A=(A_{ij}),\qquad A_{ij}=\Omega(e_{i},e_{j}),\qquad \Omega(u,v)=u^{\operatorname{t}} A,v.\]

	\item $\Omega$ \'e n\~ao degenerada se e somente se $\det (A)\neq 0$.

		Note que
		 \begin{align*}\det A=\det A^{\operatorname{t}} =\det (-A)=(-1)^{\dim V}\det (A)\\
			 \text{ implica que }\quad  \det A\neq 0\implies m=\dim V=2n
			 \end{align*}
	
	\item $\Omega\in \Lambda^{2} V^{*}$. Defina
		\begin{align*}
			\Omega^{\flat}: V &\longrightarrow V^{*} \\
			u &\longmapsto \Omega(u,\cdot )
		\end{align*}
		note que $\ker \Omega=\ker (\Omega^{\flat} )$, assim $\Omega$ \'e n\~ao degenerada se e somente se  $\Omega^{\flat}$ \'e isomorfismo.
	\end{enumerate}
	\end{remark}

\section{Aula 2}

\subsection{Subespa\c cos de evs}

Sejam $(V,\Omega)$ evs e $V\subseteq V$ subespa\c co.

\begin{defn}
	Seja $(V,\Omega)$ um espaço vetorial simplético. Dado um subespaço $W\subseteq V$, definimos seu \textit{\textbf{ortogonal simplético}} como
\[W^{\Omega} :=\{v\in V|\Omega(v,w)=0\;\forall w\in W\}\]
Note que se
\begin{align*}
	\Omega^\flat: V &\longrightarrow V^* \\
	v &\longmapsto i_v\Omega=\Omega(v,\cdot )
\end{align*}
temos que
\[\ker \Omega:=\ker \Omega^\flat=W^\Omega.\]

Considere a restri\c c\~ao  de $\Omega$ \`a W:
	\[i:W\hookrightarrow V,\qquad i^{*} \Omega=\Omega|_{W}\in \Lambda^2W^{*},\]
	ent\~ao
	\[\ker (\Omega|_{W})=W\cap W^{\Omega}\]
Casos de interesse:
\begin{itemize}
	\item \textit{\textbf{Isotr\'opico}}: $W\subseteq W^{\Omega}$ ($\iff\Omega|_{W}\equiv 0$).
	\item \textit{\textbf{Coisotr\'opico}}: $W^{\Omega} \subseteq W$.
	\item  \textit{\textbf{Lagrangiano}}: $W=W^{\Omega}$.
	\item \textit{\textbf{Simpl\'etico}}: $W\cap W^{\Omega} =\{0\}$ ($\Omega|_{W}$ \'e n\~ao degenerado (=simpl\'etico)). 
\end{itemize}
\end{defn}

\begin{lemma}
	$\dim W+\dim W^{\Omega} =\dim V$.
\end{lemma}
\begin{proof}[Demostra\c c\~ao]
	\begin{align*}
		\Omega^{1}: V &\overset{\sim}{\to }V^{*} \\
	u &\longmapsto \Omega(u,\cdot )
	\end{align*}
	Note que $W^{\Omega} \mapsto \operatorname{Ann}(W)$, assim
	\[\dim W+\dim \operatorname{Ann}(W)'=\dim V\]
\end{proof}

\begin{remark}\leavevmode
	\begin{itemize}
		\item $W\subseteq V$ subespa\c co simpl\'etico  se e somente se  $V=W\oplus W^{\Omega}$.
		\item $W$ isotr\'opico $\implies \dim W\leq \frac{\dim V}{2}$.
		\item $W$ coisotr\'opico $\implies \dim W\geq \frac{\dim V}{2}$.
		\item $W$ Lagrangiano se $\dim W=\frac{\dim V}{2}$.
	\end{itemize}
\end{remark}

De fato, $W$ \'e Lagrangiano se e somente se  $W$ \'e isotr\'opico e $\dim W=\frac{\dim V}{2}$.

\begin{exercise}\leavevmode 
	\begin{itemize}
		\item $(W^{\Omega})^{\Omega}=\Omega$ ($W$ isotr\'opico se e somente se  $W^{\Omega}$).
		\item $(W_1\cap W_2)^{\Omega} =W_1^{\Omega} +W_2^{\Omega}$.
	\end{itemize}
\end{exercise}

\begin{example}\leavevmode 
	\begin{itemize}
		\item Subespa\c cos de dimens\~ao 1 s\~ao isotr\'opicos (subespa\c cos de codimens\~ao 1 s\~ao\\ coisotr\'opicos).
		\item $V=V\oplus W^{*}$, onde $V$ tem a forma $\Omega_{can?}$ e $W$ e $W^{*}$ s\~ao Lagrangianos.
		\item $\mathbb{R}^{2n}$, $\{e_1,\ldots,e_n,f_1,\ldots,f_n\} $ base simpl\'etica, ent\~ao $\operatorname{span} \{e_i,f_i\}$ \'e simpl\'etico, e \\$\operatorname{span} \{e_1,\ldots,e_k\} $ \'e isotr\'opico (se $k=n$ \'e Lagrangiano).
		\item $(V_1,\Omega_1)$ e $(V_2,\Omega_2)$ evs's, $T:V_1\to V_2$ isometr\'ia linear, $\operatorname{graf}(T):=\{(u,Tu):u\in V_1\} \subseteq V_1\times V_2$. $T$ \'e simplectomorfismo se e somente se $\operatorname{graf}(T)$ \'e um subespa\c co Lagrangiano em $V_1\times V_2$.
		\item $\dim \operatorname{graf}(T)=\dim V_1=\frac{1}{2}\dim (V_1\times V_2)$.
		\item $\Omega_{V_1\times \bar{V_2}}((u,Tu),(v,Tv))=\Omega(u,v)-\underbrace{\Omega_2(Tu,Tv)}_{=T^{*} \Omega_2(u,v)}$ ($=0\iff \Omega_1=T^{*} \Omega_2$).
	\end{itemize}
\end{example}

\begin{thm}[Exist\^encia das bases simpl\'eticas]\leavevmode
	Para cualquer $(V,\Omega)$ evs existe uma base simpl\'etica.
\end{thm}
\begin{proof}[Demostra\c c\~ao]
	Seja $e_1\in V\setminus \{0\} $. Como $\Omega$ \'e n\~ao degenerada, existe $f_1\in V$ tal que $\Omega(e_1,f_1)=1$. Considere $W_1=\operatorname{span}\{e_1,f_1\} $. Ent\~ao $\Omega|_{W_1}$ \'e n\~ao degenerado (ie. $W_1$ \'e simpl\'etico), o que acontece se e somente se $V=W_1\oplus W_1^{\Omega}$. Assim, existem $e_2\neq 0$ in $W_1^{\Omega}$ e $f_2\in W_1^{\Omega}$ tal que $\Omega(e_2,f_2)=1$, etc… ($V=W_1\oplus  W_2\oplus \ldots\oplus W_n$). O conjunto $\{e_1,\ldots,e_n,f_1,\ldots,f_n\} $ \'e uma base simpl\'etica.
\end{proof}

\begin{exercise}
	$V$ ev de dimens\~ao $2n$ e $\Omega\in \Lambda^{2} V^{*}$ \'e n\~ao degenerada se e somente se $\Omega^{n}=\Omega\wedge \ldots\wedge \Omega\in \Lambda^{2n} V^{*} \neq 0$.
\end{exercise}

\subsection{Equival\^encia entre ev's simpl\'eticos}

$(V,\Omega)$ e $(V',\Omega')$ s\~ao \textit{\textbf{equivalentes}} se existe um \textit{\textbf{simplectomorfismo}} linear $\varphi :V\overset{\sim}{\to }V'$ (isometr\'ia linear) tal que
	\[\varphi^{*}\Omega'=\Omega\in \Lambda^{2} V^{*}\] 
onde
\[\varphi^{*} \Omega'(u,v)=\Omega'(\varphi(u),\varphi(v).\]
Dado $(V,\Omega)$ evs, definimos
\[\operatorname{Sp}(V):=\{T\in \operatorname{GL}(V)|T^{*} \Omega=\Omega\} \]
\begin{example}\leavevmode 
	\begin{enumerate}
	\item $V=\mathbb{R}^{2n}$, $\Omega_0(u,v)=-u^{\mathbf{T}} J_0v$ onde $J_0=\begin{pmatrix}0&-\operatorname{I}\\\operatorname{I}&0\end{pmatrix}$, com base can\^onica $\{e_1,\ldots,e_n,$ $f_1,\ldots,f_n\} $. Temos

\begin{equation}\label{eq:1}
	\begin{cases}
				\Omega_0(e_i,e_{j})=0\\
				\Omega_0(e_i,f_j)=\delta_{ij} \\
				\Omega_0(f_i,f_j)=0
				\qquad &
			\end{cases}
\end{equation}

\begin{defn}
	Uma base de $(V,\Omega)$ satisfazendo \cref{eq:1} \'e chamada \textit{\textbf{base simpl\'etica}}.
\end{defn}

Following Lee, Example. 22.2, the condition may be that $\Omega=\sum_{i=1}^{n} \alpha^{i} \wedge \beta^{i}$ where $\alpha^{i}$ and $\beta^{i}$ are just the dual basis covectors of the base $\{A_{1},\ldots,A_{n},B_1,\ldots,B_{n}\} $ of $V$.

\begin{remark}\leavevmode 
		Escolher/Achar uma base simpl\'etica \'e equivalente \`a escolher/achar um simplectomorfismo
	\[(V,\Omega)\overset{\sim}{\to }(\mathbb{R}^{2n},\Omega_{0})\]
\end{remark}

	\item $W$ espa\c co vetorial sobre $\mathbb{R}$, sejam $V=W\oplus W^{*}$, $w,w\in W$ e $\alpha,\alpha \in W^{*}$
		\[\Omega_{\operatorname{?}}( (w,\alpha),(w',\alpha')):=\alpha'(w)-\alpha(w')\]
		\'e n\~ao degenerada e anti-sim\'etrica. Assim,
		\[(W\oplus W^{*},\Omega_{?})\]
	\'e um espa\c co vetorial simpl\'etico.

\begin{remark}
	Se $\{e_1,\ldots,e_n\} $ \'e uma base simpl\'etica de $W$ e $\{f_1,\ldots f_n\} $ \'e a base dual de $W^{*}$, ent\~ao
\[(W\oplus W^{*},\Omega_{?}\cong (\mathbb{R}^{2n},\Omega_0).\]
\end{remark}

Note que ainda que dado 
\[A:W\overset{\sim}{\to}W\]
automorfismo ?,
\[T_{A}:=\begin{pmatrix}A&0\\0&(A^{*} )^{-1}\end{pmatrix}:W\oplus W^{*} \to W\oplus W^{*}\]
\'e simplectomorfismo, ($T_{A}=A\oplus (A^{*} )^{-1}$ ).

\paragraph{Moral:} $\operatorname{GL}(W)\hookrightarrow \operatorname{Sp}(W\oplus W^{*} )$ 
\begin{align*}
	EV&\overset{\text{funtor} }{\rightsquigarrow}EVS\\
	A\circlearrowleft W&\longmapsto W\oplus W^{*} \circlearrowright T_{A}
\end{align*}

\item $V$ ev sobre $\mathbb{C}$, $\dim_{\mathbb{C}}=n$, com produto interno hermitiano
	\[h:V\times V\to \mathbb{C}\]
i.e. satisfazendo
\begin{enumerate}
	\item $h(u,\lambda v)=\lambda h(u,v)$  $\forall \lambda\in \mathbb{C}$,
	\item $h(u,v)=\overline{h(v,w)}$,
	\item $h(u,u)>0$ $\forall u\neq 0$,
\end{enumerate}
pode ser escrito como
\[h(u,v)=g(u,v)+i\Omega(u,v)\]

Agora considere $V$ como espa\c co vetorial sobre $\mathbb{R}$ (de dimens\~ao $2n$ ).

\begin{exercise}\leavevmode 
	\begin{itemize}
		\item $g$ \'e produto interno positivo definido.
		\item $\Omega$ \'e antisim\'etrica, n\~ao degenerada (simpl\'etica).
		\item Ache uma base de $V$ (dica: extens\~ao de base ortonormal de $h$…)
		\item $\operatorname{U}(n)\subset \operatorname{SP}(V,\Omega)$.
	\end{itemize}
\end{exercise}

\item Produto direto: $(V_1,\Omega_1)$, $(V_2,\Omega_2)$ espa\c cos vetoriais.
	\[\begin{tikzcd}
		&V_1\times V_2\arrow[ld,"\pi_1",swap]\arrow[rd,"\pi_2"]\\
		V_1&&V_2
	\end{tikzcd}\]
Tem a forma simpl\'etica \'e o pullback:
	\[\Omega:=\pi_1^{*} \Omega_1+\pi_2^{*} \Omega_2\]
	ou seja,
	\[\Omega((u_1,u_2),(v_1,v_2)):=\Omega_1(u_1,v_1)+\Omega_2(u_2,v_2),\]
	que \'e n\~ao degenerado e antsim\'etrico tamb\'em.
\end{enumerate}
\end{example}

\paragraph{Nota\c c\~ao:} se $(V,\Omega)$ \'e um espa\c co vetorial simpl\'etico, denotamos por $(V,-\Omega):=\bar{V}$, que tamb\'em \'e um evs.

\section{Aula 3}

Se $V$ é ume spaço vetorial real, uma \textit{\textbf{estrutura complexa}} em $V$ é um endomorfismo linear $J:V\to V$ tal que $J^2=-\operatorname{id}$.

Seja $(V,\Omega)$ um espaço vetorial simplético. Uma estrutura complexa em $V$ é \textit{\textbf{compatível}} com $ \Omega$ se para todo $u, v\in V$,
\[g(u,v):=\Omega(u,Jv)\]
é um produto interno. (Acho que) isso implica que
\[\Omega(Ju,Jv)=\Omega(u,v),\qquad\text{e}  \qquad \Omega(u,Ju)>0\;\forall u \neq 0\]
\section{Aula 4}

\section{Aula 5}

Lembranza da \'ultima aula:
\begin{enumerate}
	\item Defini\c c\~ao de variedade simpl\'etica.
	\item Pelo menos dois exemplos.
	
	\item Forma de volume/orientabilidade.
	
	\item Campos simpl\'eticos/campos hamiltonianos.

	\item Obstru\c c\~ao cohomol\'ogica de para estrutura simpl\'etica.
\end{enumerate}

\paragraph{Hoje:} Fibrados cotangentes.

\subsection{Forma tautol\'ogica no fibrado cotangente}

Seja $Q$ uma variedade e $M :=T^* Q$ o fibrado cotangente.

\paragraph{Lembrando} Se $Q$ \'e uma variedade, $x\in Q$. O \textit{\textbf{espa\c co tangente}} em $ x$ s\~ao derivaç\~oes ou clases de equivalencia de curvas… base local do espa \c co tangente $\partial_{x_i}$…base dual disso \'e base do espa\c co cotangente nesse ponto… o fibrado cotangente $\bigsqcup_{x\in Q}T^*_{x}Q$ \'e variedade suave.

O fibrado cotangente possui uma 1-forma tautol\'ogica definida assim:

\begin{defn}
	$\alpha\in\Omega^{1}(M)$, onde $M:=T^*Q$, dada por
	\[\alpha_p(X)=p(\pi_*(X))\]
	ou seja, como $X$ \'e tangente ao fibrado cotangente, ele est\'a anclado a algum covetor, assim a gente pode evaluar ele no covector. Tamb\'em pode ser pensado como o pullback de um covector em $Q$ baixo a proje\c c\~ao cotangente usual.
\end{defn}

\begin{defn}[Monitoria]
	\[\begin{tikzcd}
		T^*M=\{(p,\xi)|\xi:T_pM\to \mathbb{R}\text{ linear} \} \arrow[d,"\pi"]\\
		M
	\end{tikzcd}\]
	A \textit{\textbf{forma tautologica}}  \'e $\lambda$ dada por
	\[\lambda_{(p,\xi)}(v)\in\mathbb{R},\qquad v\in T_{(p,\xi)(T^*M)}\]
	\'e igual a
	\[\xi(d\pi_{(q,\xi)(v)})\]
	usando o mapa
\[T_{(p,\xi)}(T^* M)\overset{d\pi_{(p,\xi)}}{\longrightarrow}T_{p}M\]
\end{defn}

Em coordenadas locais $(x_1,...,x_n,y_1,\ldots,y_n)$ do espaço cotangente, temos que
\[\lambda=\sum_{i=1}^nA_i dx_i+\sum_{i=1}^nB_i dy_i\]
Avaliando $\lambda$ nos vectores can\'onicos $\frac{\partial}{\partial x_j}\Big|_{(p,\xi)}$ e $\frac{\partial}{\partial y_j}$ notamos que $A_i=\xi\left( \frac{\partial}{\partial x_j} \right) $ pois a diferencial de $\pi$ faz as $B_j$ ser zero.

\begin{exercise}\leavevmode 
	\begin{enumerate}
		\item A 1-forma tautol\'ogica $\alpha\in\Omega^{1}(T^*Q)$ \'e a \'unica 1-forma satisfazendo
			\[\forall \mu\in\Omega^{1}(Q),\qquad \mu^*\alpha=\mu\]
			onde pensamos a $\mu$ do lado izquerdo como um mapa  $\mu:Q\to T^*Q$, ie. uma se\c c \~ao do fibrado cotangente, e do lado direito simplesmente como uma 1-corma em $Q$.

	\end{enumerate}
\end{exercise}

\begin{defn}
	$M=T^*Q$, $\alpha\in\Omega^{1}(M)$ ent\~ao a \textit{\textbf{forma simpl\'etica can\'onica}} de $T^*Q$ \'e
	\[\omega_{\operatorname{can}}=-d\alpha\]
\end{defn}

\begin{remark}\leavevmode
	\begin{itemize}
		\item $d\omega_{\operatorname{can}}=-d^{2}\alpha=0$.
		\item Formalmente $\omega=\sum_{i=1}^{n}dx_i\wedge d\xi_i$
	\end{itemize}
\end{remark}

Assim, temos uma variedade simpl\'etica can\'onica associada a toda variedade, $\\(T^* Q,\omega_{\operatorname{can}}$.

\begin{remark}\leavevmode 
	\begin{itemize}
		\item Dado $B\in\Omega^{2}(Q)$ com $dB=0$, a forma
			 \[\omega_{B}\omega_{\operatorname{can}}+\pi^*B\]
			 \'e simpl\'etica e o termo $\pi^*B$ se chama de \textit{\textbf{magn\'etico}}.

			\item Se $Q$ \'e Riemanniana com m\'etrica $g$ temos o mapa induzido
				\begin{align*}
					g^{\sharp}: TQ &\longrightarrow T^*Q \\
					u &\longmapsto g(u,\cdot )
				\end{align*}
				Assim, o pullback the $\omega_{\operatorname{can}}$ \'e uma forma simpl\'etica em $TQ$.

	Al \'em disso, a m\'etrica nos fornece de uma fun\c c\~ao Hamiltoniana dada por $H\in C^{\infty}(TQ)$, $H(v)=\frac{1}{2}g(v,v)=\frac{1}{2}\|v\|^{2}$.

	Veremos que o fluxo Hamiltoniano de $H$ em $(TQ,\omega)$ \'e fluxo geod\'esico em $Q$.

	Tem dois generaliza\c c\~oes naturais:
	\begin{itemize}
		\item  $\bar{H}(v)=\frac{1}{2}g(u,v)+V(x)$ com $V\in C^{\infty}(Q)$, mec\^anica cl\'asica.

		\item $H(v)=\frac{1}{2}g(v,v)$ com respeito a $\omega_{B}$.
	\end{itemize}
	\end{itemize}
\end{remark}

\begin{question}[Projeto?]
	Exist\^encia de \'orbitas peri\'odicas em n\'iveis de energia?
\end{question}

\begin{defn}
	O \textit{\textbf{levantamiento cotangente}} de um difeomorfismo (na mesma dire\c c\~ao do difeomorfismo) \'e $\varphi:Q_1\overset{\sim}{\to }Q_2$ \'e $\hat{\varphi}=((T\varphi)^{*})^{-1}$.
\end{defn}

\begin{question}
	Preserva a forma can\'onica?
\end{question}

\begin{prop}
	Sim. $\hat{\varphi}:T^*Q_1\to T^*Q_2$ satisfaz $\hat{\varphi}^*\alpha_2=\alpha_1$ onde $\alpha_i$ \'e a forma tautol\'ogica, para $i=1,2$. Isso implica que  $\hat{\varphi}^*\omega_2=\omega_1$.
\end{prop}

Isso implica que temos um funtor $Q\rightsquigarrow T^*Q$ que se chama de \textit{\textbf{funtor cotagente}} e permite levar problemas de geometria diferencial para a geometria simpl \'etica.

\begin{proof}[Demostra\c c\~ao]
	\[\begin{tikzcd}
		T^*Q_1\arrow[r,"\varphi"]\arrow[d,"\pi_1"]&T^*Q_2\arrow[d,"\pi_2"]\\
		Q_1\arrow[r,"\varphi"]&Q_2
	\end{tikzcd}\]
	A clave dessa prova \'e que o diagrama commuta, assim pode se-trocar um termo $\pi_2\circ \hat{\varphi}$ por $\varphi \circ \pi_1$.
\end{proof}

O funtor que produzimos $\operatorname{Dif}(Q)\hookrightarrow \operatorname{Simp }(T^*Q$ n\~ao e fiel (surjetivo), ie. existem simplectomorfismos no fibrado cotangente que n\~ao vem de difeomorfismos na variedade.

\begin{remark}
	Dada uma 1-forma $A\in\Omega^{1}$. Pode se-produzir um mapa no cotangente simplesmente trasladando por $A$:
	\begin{align*}
		T_A: T^*Q &\longrightarrow T^*Q \\
		(x,\xi) &\longmapsto (x,\xi+A_x)
	\end{align*}
	que n\~ao pode ser um levantamento porque se projecta na identidade!

	\begin{exercise}
		$T_A$ \'e um simplectomofrismo $\iff$ $dA=0$.
	\end{exercise}
\end{remark}

Mas, como sabemos quais simplectomorfismos no cotangente s\~ao sim levantamentos de difeomorfismos na variedade?

\begin{exercise}
	Seja $F:T^*Q\to T^*Q$ um simplectomorfismo. Quando $F=\hat{\varphi}$ \'e levantamento de algum $\varphi:Q\overset{\sim}{\to } Q $. Pois, isso acontece $\iff$ $F$ preserva a forma tautol\'ogica, ie. $F^*\alpha=\alpha$.
\end{exercise}

\begin{remark}
	Levantamento cotangente de campos de vetores. Come\c ca com um campo $X\in\mathfrak{X}(Q)$, integra para obter um fluxo $\varphi_{t}$, que \'e uma fam\'ilia de difeomorfismos na variedada, voc\^e sabe levantar isso com o funtor obtendo outro fluxo (porque levantamento de fluxo \'e fluxo) $\hat{\varphi}_{t}$, e diferenciando obt\'em $\hat{X}\in\mathfrak{X}(T^*Q)$.
\end{remark}

\begin{remark}
	Para cualquer fibrado vetorial $E \to  M$, podemos ver a se\c c\~oes $\Gamma(E)$ como um subconjunto das fun \c c\~oes suaves na variedade $C^{\infty}(E)$---s\~ao as fun\c c\~oes lineares nas fibras. A\'i tem um modo natural de definir para cualquer campo vetorial $X\in\Gamma(TQ)\subseteq C^{\infty}(T^*Q)$ uma fun\c c\~ao, $H_X(p)=p(X_{\pi(p)}=\alpha(\hat{X})$.
\end{remark}

\begin{prop}
	$\hat{X}=$ campo Hamiltoniano de $H_X$.
\end{prop}

\section{Aula 6}

Hoje: Colchete de Poisson, Darboux.

\subsection{Colchete de Poisson}


$M$ variedade, $\omega\in\Omega^2(M)$ n\~ao degenerada (quase-simpl\'etica). Podemos fazer

\begin{align*}
	w^\flat: TM &\longrightarrow T^*M \\
	x &\longmapsto i_X\omega
\end{align*}
So that
\[f\in C^\infty(M)\rightsquigarrow X_f\in\mathfrak{X}(M)\]
e
\[i_{X_f}\omega=df.\]

\begin{defn}
	$f,g\in C^\infty(M)$.
	\begin{align*}
		\{\cdot ,\cdot \}: C^\infty(M)\times C^\infty(M) &\longrightarrow C^\infty(M) \\
		\{f,g\} & =\omega(X_g,X_f)=dg(X_f)=\mathcal{L}_{X_f}g=-\mathcal{L}_{X_g}f 
	\end{align*}
	
\end{defn}

\begin{prop}[Exerc\'icio]
	$d\omega=0\iff \{\cdot ,\cdot \}$ satisfaz identidade de Jacobi.  $\implies (M,\omega)$ simpl\'etica, $\{\cdot ,\cdot \}$ \'e colchete de Lie em $C^\infty(M)$ e isso se chama de um \textit{\textbf{colchete de Poisson em $(M,\omega)$}}.
\end{prop}

\begin{exercise}
	$\{f,gh\} =\{f,g\} h+\{f,h\} g$.
\end{exercise}

\begin{example}
	$\mathbb{R}^{2n}$.
\end{example}

\begin{defn}
	$f,g\in C^\infty(M)$ est\~ao em \textit{\textbf{involu\c c\~ao}} se $\{f,g\} =0$. ie. $X_g$ \'e tangente aos n\'iveis $f=$const (e vice versa).
\end{defn}

\begin{remark}
	Nesse caso, a derivada de $g$ ao longo das curvas integrais de $X_f$ \'e zero.
\end{remark}

\paragraph{Motiva\c c\~ao} $(M,\omega)$ simpl\'etica, $H\in C^\infty(M)$ queremos integrar $X_H$ (ie. resolver $c'(t)=X_H(c(t))$). Suponha que existe $f\in C^\infty(M)$ com $\{f,H\} =0$, chamada \textit{\textbf{integral primeira}}. ie.  $f$ \'e constante ao longo do fluxo Hamiltoniano.

No s\'eculo XIX, quando Poisson vivia, a ideia era que se temos um n\'umero sufieiente de integrais primeiras "independentes", podemos "integrar" $X_H$. (Aqui "integrar" significa dar uma solu\c c\~ao a equa\c c\~ao diferencial do fluxo Hamiltoniano).

Em 1810, Poisson deu a f\'ormula
\[\{f,g\} =\sum_{i,j}\frac{\partial f}{\partial p_i}\frac{\partial g}{\partial q_i}-\frac{\partial f}{\partial q_i}\frac{\partial g}{\partial p_i}.\]

\begin{thm}[Poisson]
	$\{f,H\} =0=\{g,H\} \implies \{\{f,g\},H\} =0$.
\end{thm}

\begin{thm}[Jacobi]\leavevmode
	\[\{H,\{f,g\}\} +\{g,\{H,f\}\} +\{f,\{g,H\}\} =0\]
\end{thm}

\paragraph{1880} Lie usou essa identidade no seu trabalho de transforma\c c\~oes (\'algebras de  Lie).

\paragraph{Vers\~ao moderna (sec.  XX) de integrabilidade} Veremos adiante…

\begin{thm}[Arnold-Liouville]\leavevmode
	$(M,\omega)$ de dimens\~ao $2n$ e seu Hamiltoniano $H=f_1$ que \'e a primeira de uma sequencia de $n=\dim M/2$ fun\c c\~oes independentes (as derivadas s\~ao linearmente independentes) $f_2,\ldots,f_n\in C^{\infty}(M)$ tais que $\{f_i,f_j\} =0$ e que $(f_1,\ldots,f_n):M\to \mathbb{R}^{n}$ \'e uma submers\~ao. Ent\~ao
	\[N=\{(f_1,\ldots,f_n)=\text{cte} \} \cong \mathbb{T}^n\]
	se compacto e conexo. Al\'em disso, a din\^amica de $X_H$ em $\mathbb{T}^n$ \'e quase peri\'odica (=\'e um fluxo linear no toro, que pode ser racional ou irracional).
\end{thm}

\begin{remark}[Projeto?]\leavevmode
	Qu\'e acontece com essa din\^amica no toro se perturbamos o sistema? O problema de dois corpos \'e completamente integravel. Por exemplo, a din\^amica da Terra e o Sol pode se-resolver, mas o problema adicionando a Lua \'e o problema de 3 corpos, que ningu\'em sabe c\'omo resolver. Aqui a Lua \'e uma perturba\c c\~ao.

	Teorema KAM, quanto mais irracional \'e o fluxo, mais robusto \'e o toro, mais inest\'avel.
\end{remark}

Em fim, tudo isso para motivar os colchetes de Poisson.

\subsection{Teorema de Darboux}

$(M,\omega)$ variedade simpl\'etica com o colchete $\{\cdot ,\cdot \}$.

\begin{remark}\leavevmode
	\begin{enumerate}
		\item $\omega$ est\'a completamente determinada por $\{\cdot ,\cdot \}$, ie. se duas estruturas simpl\'eticas d\~ao lugar ao mesmo colchete de Poisson, elas s\~ao iguais.Por que?
			\begin{align*}
				\omega^\sharp: T^*M &\longrightarrow TM 
			\end{align*}
			est\'a dada em cada ponto por
			\[(\omega^\sharp )_{ij}=\{x_i,x_j\}\]
			por defini\c c\~ao.

			{\color{magenta}(My interpretation)} Especificamente, considere coordenadas de Darboux \\$(x^1,\ldots,x^n,y^1,\ldots,y^n)$ en $M$. Em \cite{lee}, eq. 22.9 vemos que para qualquer fun\c c\~ao $f \in\mathcal{C}^\infty(M)$, o seu campo Hamiltoniano est\'a dado por
			\[X_f=\sum_{i=1}^n\left( \frac{\partial f}{\partial y^i}\frac{\partial }{\partial x^i}-\frac{\partial f}{\partial x^i} \frac{\partial }{\partial y^i} \right) \]
			assim,
			\[X_{x^i}=-\frac{\partial }{\partial y^i},\qquad X_{y^i}=\frac{\partial }{\partial x^i}.\]
			Isso \'e uma base do espaço tangente. De fato, para qualquer base $v_1,\ldots,v_n$ de um espaço vetorial com base dual $v^*_1,\ldots,v_n^*$, se $w=\sum_{i}w^iv_i$, \'e super básico que
			\[\omega\left( v_i,w \right) =\omega\left( v_i,\sum_{j}w^jv_j \right) =\sum_{j}\omega(v_i,v_j)w^j=\sum_{j}\omega(v_i,v_j)v_j^*(w)\]
			ie.
			\[i_{v_i}\omega=\omega(v_i,\cdot )=\sum_{j}\omega(v_i,v_j)v_j^*\]
			Da\'i, em coordenadas,
			\[\omega^\flat(v_i)=\begin{pmatrix} \omega(v_i,v_1)\\\vdots \\\omega(v_i,v_{n} \end{pmatrix} \]
			ou seja
			\[\omega^\flat =\begin{pmatrix} \omega(v_1,v_1)&\cdots &\omega(v_n,v_1)\\
			\vdots &&\vdots \\
		\omega(v_1,v_n)&\cdots &\omega(v_n,v_n)\end{pmatrix} \]
		Agora note que os vetores Hamiltonianos associados as coordenadas de Darboux $(x^1,\ldots,x^n,y^1,\ldots,y^n)$ satisfazem as rela\c c\~oes do seguinte item nesta observa\c c\~ao (pode comprovar isso usando a fórmula do colchete de Poisson em coordenadas de Darboux). Da\'i, nessa base de vetores Hamiltonianos,

		\[\omega^\flat =\begin{pmatrix} 0&-I\\I&0 \end{pmatrix} \]
		Mas que n\~ao \'e que a gente tava buscando $\omega^\sharp$? Pois \'e, essa matriz elevada ao quadrado \'e $-\operatorname{id}$, da\'i a sua inversa \'e s\'o botar um signo menos…

		\item A estrutura simpl\'etica can\'onica $\omega_0=\sum_{i}dp_i\wedge dp_i$ em $\mathbb{R}^{2n}$ est\'a determinada (\'e a \'unica tal que) por
			\[\{q_i,q_j\}=0=\{p_i,p_j\},\qquad \{p_i,q_j\} =\delta_{ij}.\]
			\'E como se tivesse uma base simpl\'etica boa em todos os pontos…
	\end{enumerate}
\end{remark}

\begin{thm}[Darboux]\leavevmode
	$(M,\omega)$ simpl\'etica, ent…åo ao redor de todo ponto $x\in M$ existem coordenadas locais $(q_1,\ldots,q_n,p_1,\ldots,p_n)$ tais que $\omega=\sum_{i=1}^ndq_i\wedge dp_i$, ou, equivalentemente vale
		\[\{q_i,q_j\}=0=\{p_i,p_j\},\qquad \{p_i,q_j\} =\delta_{ij}.\]
\end{thm}

Tem um lema que va a provar essencialmente tudo.

\begin{lemma}[Primeiro paso da indu\c c\~ao ]
	Ao redor de qualquer ponto $x\in M$ existem coordenadas $(q,p,y_1,\ldots,y_{2n-2}$ tais que
	\[1=\{p,q\},\quad \{p,y_j\}=0=\{q,y_j\},\qquad \{y_i,y_j\} =\varphi_{ij}(y).\]
Ou seja, a matriz da forma \'e
\[\begin{pmatrix} \begin{pmatrix}0&-1\\1&0\end{pmatrix}& 0&0&\ldots&0\\
0&&&&\\
0&&A(y)&&\\
0&&&&\end{pmatrix} \]
ou seja, temos uma expres\~ao
\[\omega=dq\wedge dp+\omega_N\]
onde $\omega_N$ \'e dada por $A(y)$ e  \'e simpl\'etica.
\end{lemma}

\begin{proof}[Demostra\c c\~ao do Lema]
	\begin{enumerate}[label=\textbf{Paso \arabic*}]
		\item Seja $p$ uma fun\c c\~ao tal que $X_p(x)\neq 0$. Pelo teorema de fluxo tabular (retifica\c c\~ao ) existe uma fun\c c\~ao $q$ tal que $X_p=\frac{\partial}{\partial q}$, de modo que $\{p,q\}=dq(X_p)=1$ e $dp(X_q)=-1$.

		\item En\~ao $X_p$ e $X_q$ s\~ao linearmente independentes, pois $1=\{p,q\} =\omega(X_p,X_q)\neq 0$, o que aconteceria por antisimetria se s\~ao linearmente dependentes. Al\'em disso, comutam, pois
			\[ [X_p,X_q]\overset{\text{aula pasada?} }{=}X_{\{p,q\} =1}=0.\]

			Agora usamos a generaliza\c c\~ao do teorema do fluxo tabular: se $X_1,\ldots,X_k$ s\~ao campos  linearmente independentes e que comutam, ent\~ao existem coordenadas\\ $(x_1,\ldots,x_n)$ dais que $X_i=\frac{\partial}{\partial x_i}$. (Teo. fun\c c\~ao inversa.) Assim, existem coordenadas locais $y_1,\ldots,y_{2n}$ tais que
			\[X_q=\frac{\partial}{\partial y_{2n-1}},\qquad  X_p=\frac{\partial}{\partial y_{2n}}.\]
			Logo
			\[dy_j(X_q)=0=dy_j(X_p)\]
			para $j=1,\ldots,2n-2$.

		\item As diferenciais
		\[dq,dp,dy_1,\ldots,dy_{2n-2}\]
		s\~ao linearmente independentes, pois se
		\[adq+bdp+\sum_{i}c_{ij}y_i=0\]
		pois as $y_i$ j\'a s\~ao LI, e avaliando em $X_i$ obtemos $a=0$, e no $X_q$ que $b=0$.

		Temos um sistema de coordenadas  $(q,p,y_1,\ldots,y_{2n-2}$ ao redor de $x$ tal que as condi\c c\~oes  do teorema salvo a \'ultima se cumplem. Agora veamos que $\{y_i,y_j\}$ n\~ao depende de $p,q$.

		 \item S\'o lembrar que
			 \[X_q=-\frac{\partial}{\partial p},\qquad X_q=\frac{\partial}{\partial q}\]
			 assim
			 \[ \frac{\partial}{\partial p}\{y_i,y_j\} =-\{q,\{y_i,y_j\}\} =0\]
			 onde a segunda igualdade \'e jacobi. Fim.
	\end{enumerate}
\end{proof}

\begin{proof}[Demostra\c c\~ao do Teo. Darboux]
Segue do lema por indu\c c\~ao
\end{proof}

\begin{defn}
Uma \textit{\textbf{estrutura de Poisson}} em uma variedade $M$ \'e
\begin{align*}
	\{\cdot ,\cdot \}: C^\infty(M)\times C^\infty(M) &\longrightarrow C^\infty(M)
\end{align*}
$\mathbb{R}$-bilinear, antisim\'etrica, Jacobi e Leibniz, ie. $\{f,gh\} =\{f,g\} h+\{f,h\} g$.
\end{defn}

\begin{example}\leavevmode 
	\begin{itemize}
		\item $(M,\omega)$ simpl \'etica com $\{f,g\} =\omega(X_g,X_f)$.
	\end{itemize}
\end{example}

\section{Aula 7}

Na aula passada vimos:

\begin{itemize}
	\item Colchetes de Poisson.
	\item Teorema de Darboux. Prova: demostrar que tem relaç\~oes que caractetizam a forma de maneira \'unica.
	\item \'E poss\'ivel descrever estruturas cimpl\'eticas en termos de colchete de Poisson.: Variedades de Poisson. Issto \'e axiomatizar as propriedades b\'asicas do colchete de Poisson. Esses objetos podem ser entendidos como foleaç\~oes simpl\'eticas.
\end{itemize}

\subsection{Subvariedades}

Seja $(M,\omega)$ simpl\'etica e $N\overset{i}{\hookrightarrow }(M,\omega)$. Ent\~ao temos
\[\omega_N=i^* \omega\in\Omega^2(N)\]
que \'e fechada porque o pullback comuta com derivada exterior.

\begin{align*}
	 \ker(\omega_N)& =\{X\in T N:\omega(X,Y)=0\;\forall Y\in T N\} \\
	 &=T N\cap T N^\omega\subseteq T N
\end{align*}

\subsection{Pausa para distribuiç\~oes}
$P$ variedade.

\begin{defn}
	Uma \textit{\textbf{distribuiç\~ao (generalizada)}} em $P$ \'e
	\[P\ni x\longmapsto D_x\subseteq T_xP\text{ subespaço} \]
	e o posto da distribuiç\~ao em $x$ := $\dim D_x$.

	A distribuiç\~ao \'e \textit{\textbf{suave}} se para todo $x_0\in P$, $\forall v\in D_{x_0}$ existe um campo vetorial $X\in\mathfrak{X}(P)$ que extende a $v$ e est\'a contido na distribiuç\~ao no sentido de que $X_x\subseteq D_x\forall x$ e $X_{x_0}=v$.
\end{defn}

\begin{example}
	N\'uncleo de 2-formas \'e um exemplo de distribui\c c\~ao, mas n\~ao \'e suave em geral.
\end{example}

\begin{defn}
	Uma distribui\c c\~ao suave $D\subseteq TP$ \'e dita \textit{\textbf{integravel}} se $\forall x\in P$ existe uma subvariedade $S\ni x$,  $TS=D|_{S}$
\end{defn}

No caso de uma dsitribui\c c\~ao (suave) integr\'avel, por todo ponto passa uma subvariedade integral conexa maximal chamadas \textit{\textbf{folhas}}.

\begin{remark}\leavevmode 
	\begin{itemize}
		\item Distribui\c c\~ao suave, de posto constante \'e a mesma coisa que um subfibrado vetorial $D\subseteq TP$. Nesse caso,

		\begin{thm}[Frobenius]
			$D$ \'e integr\'avel se e somente se \'e \textit{\textbf{involutivo}}, ou seja  \[[\Gamma(D),\Gamma(D)]\subseteq \Gamma(D).\]
		\end{thm}

		\begin{proof}[Demostra\c c\~ao]
			Note que $\implies $ \'e trivial porque se tem uma  variedade que realiza a distribui\c c\~ao, o colchete de Lie sempre vai ser outro campo vetorial tangente.
		\end{proof}
		
	\item Suponha que $D=\ker(\omega)$ com $\omega\in\Omega^2(P)$ \'e suave  $\iff$ $D$ tem posto constante. Aqu\'i $\impliedby$ \'e f\'acil.

	\item Se $d\omega=0 \implies D=\ker \omega$ \'e involutivo.
	\end{itemize}
\end{remark}

\paragraph{Conclus\~ao}  Se $\omega$ \'e uma 2-forma fechada e $D=\ker \omega$ tem posto constante, da lugar a uma folhea\c c\~ao (regular=folhas de mesma dimens\~ao) em $P$.


\subsection{Voltando}

\begin{defn}
	$N$ \'e dita
	\begin{itemize}
		\item \textit{\textbf{isotr\'opica}} quando $T_x N\subseteq T_X N^\omega\iff \omega_N=0\iff \ker \omega_N=T N$.

		\item \textit{\textbf{coisotr\'opica}} quando $T_x N^\omega\subseteq T_x N$.

		\item \textit{\textbf{lagrangiana}} quando $T_x N=T_x N^\omega\iff i^*\omega=\omega_N=0$ e $\dim N=\dim M/2$.

		\item \textit{\textbf{simpl\'etica}} $T_x N\cap(T_x N)^\omega=\{0\}\;\forall x\in N\iff \omega_N$ \'e simpl\'etica.

		\item \textit{\textbf{posto constante}} $T_x N\cap T_x N^\omega\subseteq T_x N\;\forall N$ tem posto constante.
	\end{itemize}
\end{defn}

\begin{example}\leavevmode 
	\begin{itemize}
		\item curvas s\~ao isotr\'opicas.
		\item hipersuperficies s\~ao coisotr\'opicas.
		\item Veremos v\'arios exeplos de subespaços lagrangianos.
	\end{itemize}
\end{example}

\subsubsection{Sobre subvariedades coisotr\'opicas}

\textbf{Isto tamb\'em vale para subvariedades de posto constante.} 

Vamos ver uma vers\~ao geom\'etrica de um exer\'icio da lista 1, onde pegabamos o quociente de um espaço vetorial por el n\'ucleo de uma forma para obter um espaço vetorial simpl\'etico.

\begin{exercise}
	Suponha que as folhas da folhea\c c\~ao s\~ao fibras de uma sobmers\~ao

	\[\begin{tikzcd}
		N\arrow[r,hook]\arrow[d,two heads,swap,"q"]&(M,\omega)\\
		B=N/\sim
	\end{tikzcd}\]
	ent\~ao existe uma forma simpl\'etica $\bar{\omega} \in\Omega^2(B)$ tal que $q^*\bar{\omega} =\omega_N$.
\end{exercise}

\begin{example}
	O fluxo hamiltoniano do oscilador harm\'onico $H(p,q)=\frac{1}{2}\sum_{i}q_i^2+p_i^2$ com $c=1/2$ da  $\mathbb{C}P^{n-1}$
\end{example}

\begin{exercise}
	$\psi:M\to \mathbb{R}^{k}$, $\psi=(\psi_1,\ldots,\psi_k)$. $N=\psi^{-1}(c)$ para $c$ valor regular.
	\begin{itemize}
		\item $N$ coisotr\'opico $\iff \{\psi_i,\psi_j\}|_{N} =0$.

		\item $N$ simpl\'etico $\iff \left( \{\psi_i,\psi_j\} |_{N} \right)_{ij}$ \'e invert\'ivel.
	\end{itemize}
\end{exercise}

\section{Aula 8}

Lembre:

\begin{itemize}
	\item Subvariedades lagrangianas, (co-)isotr\'opicas, simpl\'eticas. Aprofundamos\\ nas coisotr\'opicas (posto constante), como as hipersuprficies ou conjuntos de n\'ivel, que tem uma folhea\c c\~ao, e com condi\c c\~oes  de regularidade pode passar para o espaço quociente, que \'e simpl\'etico, como $\mathbb{C}P^{n}$.
\end{itemize}

\subsection{Alguns exemplos de subvariedades lagrangianas}

\begin{example}
	Dois variedades simpl\'eticas e um difeomorfismo entre elas. Ent\~ao $\varphi$ \'e simplectomorfismo se e s\'o se seu gr\'afico \'e lagrangiano. Talvez isso pode ser ussado para pensar em simplectomorfismos em um objeto cuantico.
\end{example}

\begin{remark}
	Considere
	\begin{align*}
		\varepsilon: M_1 &\longrightarrow M_1\times M_2 \\
		x &\longmapsto (x,\varphi x)
	\end{align*}
	ent\~ao o grafo de $\varphi$ \'e lagrangiano $\iff \omega_1-\varphi^*\omega_2$.
\end{remark}

\begin{example}[no fibrado cotangente]\leavevmode 
	\begin{itemize}
		\item A se\c c\~ao zero $Q\hookrightarrow T^*Q$ \'e nos mostra que $Q$ \'e uma subvariedade lagrangiana.

		\item A fibra (cotangente) de um ponto tamb\'em \'e uma subvariedade lagrangiana de $T^*Q$.

		\item Logo, o espaço de fibras?

		\item Pegue uma subvariedade da base $S\subset Q$. Considera o \textit{\textbf{fibrado conormal}} $N^*S$, $\nu_S^*$. \'E o dual do fibrado tangente. \'E o anulador de $TS$,  $\{(x,\xi)\in T^*Q:x\in S, \xi|_{TxS}=0\}$. Note que \'e um subfibrado do fibrado cotangente.
	\end{itemize}
\end{example}

Os dois exemplos anteriores s\~ao $S=Q$ e  $S=\{x\}$ da seguinte prop:

\begin{prop}
	$N^*S\hookrightarrow T^*Q$ \'e (um subfibrado) uma subvariedade lagrangiana.
\end{prop}

\begin{proof}[Demostra\c c\~ao]
	Usando coordenadas adaptadas e a forma tautol\'ogica do $T^*Q$, damos coordenadas $N^*Q$ da forma $(x_1,\ldots,x_k,\xi_{k+1},\ldots,\xi_n)$ e assim o pullback da forma tautol\'ogica \'e zero porque ele evalua os covectores $\xi_{\text{grande} }$ em vectores $x_{\text{pequeno} }$.
\end{proof}

\begin{example}
	Uma forma $\mu$ vista como se\c c\~ao do fibrado cotangente pode ser pensada como um mergulho de $Q$ em $T^*Q$.
\end{example}

\begin{prop}
	Essa subvariedade \'e lagrangiana $\iff d\mu=0$.
\end{prop}

\subsection{M\'etodo de Moser}

\begin{upshot}
	Moser's trick is a thing that gives you a diffeomorphism that pulls back $\omega_2$ to $\omega_1$.
\end{upshot}

Dadas dois formas simpl\'eticas numa variedade, como podemos achar um simplectomorfismo entre elas? A ideia do m\'etodo \'e assim:

\begin{enumerate}[label=\textbf{Step \arabic*}]
	\item Interpolar as dois formas mediante uma familia cont\'inua $\omega_t$ de formas simpl\'eticas.

	\item Buscar uma (isotop\'ia) fam\'ilia de difeomorfismos $\varphi_t$ com $\varphi_0=\operatorname{id}$ e tal que $\varphi^*_t\omega_t=\omega_0$. Com isso a gente procura levar o problema para uma EDO.

	\item Os fluxos s\~ao isotop\'ias com uma rela\c c\~ao de comutatividade. Eles correspondem com campos vetoriais. As isotop\'ias em geral est\~ao em correspond\^encia com \textit{\textbf{campos de vetores n\~ao aut\'onomos}}.
\end{enumerate}

\begin{defn}
	Uma fam\'ilia suave de difeomorfismos $\{\phi_t\}$ com $\phi_0=\operatorname{id}$ \'e chamada \textit{\textbf{isotop\'ia}}. Suave significa que $(t,x)\mapsto \varphi_t$ \'e suave.
\end{defn}

\begin{example}
	Fluxos (complets) s\~ao isotop\'ias tq $\phi_s\circ \phi_t=\phi_{s+t}$.
\end{example}

\begin{defn}
	Um \textit{\textbf{campo de vetor $t$-dependente}} ou  \textit{\textbf{n\~ao aut\'onomo}} \'e fam\'ilia suave de campos $X_t\in\mathfrak{X}(M)$. De novo, suave \'e que $(t,x)\mapsto X_t(x)$ \'e suave.
\end{defn}

\[\text{isotop\'ia} \leftrightarrow \text{campos $t$-dependentes}  \]
A diferencia\c c\~ao sempre \'e simples n\'e? Fixa um ponto e varia o tempo, obt\'em uma curva.
\[\phi_t\mapsto X_t(\varphi_t(x)=\frac{d}{d\tau}|_{t=\tau}\varphi_\tau(x).\]

A rec\'iproca \'e mais dif\'icil. A ideia e extender a variedade  \'a $M\times \mathbb{R}$, e considerar $\overline{X}(x,t)=(X_t(x),\frac{d}{dt})$. Esse depende do tempo, assim podemos achar um fluxo $\phi_t$ de $\overline{X}_t$. Aqui se deve extender o fluxo usando bump functions, assim a gente tem que $\phi_t$ est\'a definido para toda $t$.

Note que $\phi_t(x,s)=(G_t,t+s)$ para alguma fun\c c\~ao $G$ na variedade. Podemos achar uma inversa dela assim:
\[(x,s)=\phi_{-t}(\phi_t(x,s))=G_{-t}(G_t(x,s),t+s),s)\]
ie. a inversa de
\[x\mapsto G_t(x,s)\]
\'e
\[y\mapsto G_{-t}(y,s+t)\]
Logo, 
\[\phi_t(x)=G_t(x,0)\]
\'e uma isotop\'ia e como a derivada do fluxo
\[\frac{d}{dt}\phi_t(x,0)=\overline{X}(G_t(x,0),t)\implies \frac{d}{dt}G_t(x,0)=X_t(x,0)).\]
E \'e isso. Temos a correspondencia.

Voltando ao m\'etodo de Moser, para achar $\varphi^*\omega_1=\omega_0$, pegamos uma isotop\'ia que puxa $\omega_t$ em $\omega_0$, e queremos diferenciar a isotop\'ia. No caso de um fluxo, trata-se da derivada de Lie por defini\c c\~ao.

\begin{lemma}
	$\{\varphi_t\}$ isotop\'ia em $M$, $\{X_t\}$ campo aut\'onomo. Sejam $\eta\in\Omega^k(M)$, $\beta_t\in\Omega^k(M)$. Ent\~ao vale:
	\[\frac{d}{dt}(\varphi^*_t\epsilon)=\varphi^*_t(\mathcal{L}_{X_t}\eta\]
	onde estamos pegando a derivada num tempo $t$ fixo. Da\'i veremos que pela regra da cadeia segue que
	\[\frac{d}{dt}(\varphi^*_t\beta_t)=\varphi^* (\mathcal{L}_{X_t}\beta_t+\frac{d}{dt}\beta_t\]
\end{lemma}

\begin{proof}[Demostra\c c\~ao]
	\begin{enumerate}[label=\alph*.]
		\item Considere os seguintes operadores em $\Omega^\bullet$ :
			\[Q_1(\eta)=\frac{d}{dt}\varphi^*_t\eta,\qquad Q_2(\eta)=\varphi^*_t\mathcal{L}_{X_t}\eta\]
			Da\'i note que esses operadores comutam com a derivada exterior, s\~ao Leibniz respeito ao producto cunha e coincidem em fun\c c\~oes . Da\'i segue que $Q_1=Q_2$.

		\item A regra da cadeia diz que para uma fun\c c\~ao $F(a,b)$,
\[\frac{d}{dt}F(t,t)=\frac{\partial}{\partial a}F(t,t)+\frac{\partial}{\partial b}F(t,t)\]
e olha para $\varphi^*_a\beta_b$ como a $F$. Sustiuindo e usando  a., o resultado segue.
	\end{enumerate}
\end{proof}

Uma aplica\c c\~ao disso \'e

\begin{thm}[de estabilidade de Moser]\leavevmode
	$M$ compacta, $\{\omega_t\}$ formas simpl\'eticas, $t\in [0,1]$. Se as formas s\~ao todas cohomologas ent\~ao elas s\~ao simplectomorfas, i.e. $[\omega_t]=[\omega_0]\implies \exists \phi_t$ tq $\phi^*_t\omega_t=\omega_0$. Ou, de outra forma, se existe uma fam\'ilia suave de formas $\beta_t$ tais que
	 \[\omega_t=\omega_0+d\beta_t\]
	 ent\~ao existe uma isotop\'ia $\{\varphi_t\}$ tal que $\varphi^*_t\omega_t=\omega_0$.
\end{thm}

\begin{proof}[Demostra\c c\~ao]
	Note que n\~ao \'e imediato que as clases de cohomologia nos dem uma famil\'ia suave, mas \'e equivalente sim (usando decomposi\c c\~ao de Hodge? Tem algo mais simples?). O m\'etodo \'e achar um campo de vetores aut\'onomo resolvendo
	\[i_{X_t}\omega_t=-\frac{d}{dt}\beta_t\]
	pois dela segue que
	\[\mathcal{L}_{X_t}\omega_t=-d\left( \frac{d}{dt}\beta_t \right)\]
	E da\'i a segunda afirma\c c\~ao do lema.
\end{proof}

\section{Aula 9}

Lembre: M\'etodo de Moser.
 
A prova foi:
\begin{proof}[Demostra\c c\~ao]
	Calcule
	\[\frac{d}{dt}\varphi^*_t\omega_t=0\]
	isso implica que
	\[\mathcal{L}_{X_t}\omega_t=-d\left( \frac{d}{dt}\beta_t \right)\]
	e isso que
	\[i_{X_t}\omega_t=-\frac{d}{dt}\beta_t\]
\end{proof}

Com isso conseguimos associar uma isotop\'ia a um campo $t$-dependente (integra\c c\~ao).

\subsection{Aplica \c c\~ao ao teorema de Darboux}

\begin{lemma}
	$X_t$ campo de vetores $t$-dependente, $t\in\mathbb{R}$. Suponha que
	\[X_t|_{x_0}=0\quad \forall t.\]
	Ent\~ao existe uma vizinhança $U\ni x_0$ e uma familia $\varphi_t:U\to M$ de 
	\begin{itemize}
		\item (Inclus\~ao ) $\varphi_0=\operatorname{id}$.
		\item $\frac{d}{dt}\varphi_t(x)=X_t(\varphi_t(x))$
		\item $\varphi_t(x_0)=x_0$
		\item $\varphi _t:U\overset{\operatorname{difeo}}{\longrightarrow}\varphi_t(U)$.
	\end{itemize}
\end{lemma}

\begin{proof}[Demostra\c c\~ao]
	Varia\c c\~ao do caso $M$ compacto
	\[\bar{X}(x,t):=\left( X_t(x),\frac{d}{dt}\right)\quad \text{em $M\times \mathbb{R}$} \]
	\[\bar{X}(x_0,t)=\left( 0,\frac{d}{dt} \right) \]
	assim existe uma curva integral $\gamma(t)=(x_0,t)$ de $ \bar{X}$ por $(x_0,0)$ est\'a definida para toda $t\in\mathbb{R}$.

	Por EDO, existe uma vizinhança $W$ de $(x_0,0)$ em $M\times \mathbb{R}$ onde o fluxo de $\bar{X}$ existe $\forall t\in[0,1]$.

	Tome $U=\bigcap_{w\{M\times \{0\}\}} $.
\end{proof}

Valem a f\'ormula para $\frac{d}{dt}(\varphi_t^*\omega_t)$…

\begin{thm}[Darboux]\leavevmode
	$(M,\omega)$ simpl\'etica, $\dim M=2n$. Para todo $x\in M$ existe uma vizinhança $U\ni x$, aberto  $0\in V\subseteq \mathbb{R}^{2n}$ e um difeomorfismo
	\begin{align*}
		\phi: V\subseteq \mathbb{R}^{2n} &\longrightarrow U\subseteq M \\
		0 &\longmapsto x
	\end{align*}
	tal que 
	\[\phi^* \omega=\sum_{i}dq_i\wedge dp_i.\]

	[Desenho de carta coordenada]
\end{thm}

\begin{proof}[Demostra\c c\~ao]
	Podemos assumir que $M$ \'e bola aberta de $\mathbb{R}^{2n}$ com estrutura sumpl\'etica $\omega$ aribtr\'aria.

	Para usar o m\'etodo de Moser, definamos
	\begin{align*}
		\omega_1& =\omega\\
		\omega_0=\sum_{i}dq_i\wedge dp_i
	\end{align*}
	Podemos assumir que na origem
	\[\omega_1|_{x=0}=\omega_0|_{x=0}\qquad T_0\mathbb{R}^{2n}=\mathbb{R}^{2n}\]
	simplesmente porque qualquer dois formas simpl\'eticas s\~ao equivalentes num espaço vetorial simpletico usando uma mudança de coordenadas.

\begin{itemize}
	\item	Podemos assumir pelo Lema de Poincar\'e que
	\[\omega_1-\omega_0=d\beta,\quad \beta|_{0}=0\]
supondo pela mesma raz\~ao que antes que $\beta|_{0}=0$.
	\item $\omega_t=(1-t)\omega_0+t\omega_1\iff \omega_t=\omega_0+td\beta$ 
\end{itemize}
Precisamos checar que $\omega_t$ s\~ao n\~ao degeneradas numa vizinhança de 0. 

Note que em $x=0$, $\omega_t|_{x=0}=\omega_0|_{x=0}=\omega_1|_{x=0}$, assim $\omega_t |_{x=0}$ \'e n\~ao degenerada para toda $t$, mas precisamos de uma vizinhança, n\~ao s\'o um ponto.

\begin{lemma}
	Se tem uma fam\'ilia $\omega_t|_{x_0}$ \'e simpl\'etica $\forall t$, $t\in [0,1]$, ent\~ao existe uma vizinhança de $x_0$ onde $\omega_t$ \'e n\~ao degenerada $\forall t\in[0,1]$.
\end{lemma}

\begin{proof}[Demostra\c c\~ao]
	Considere
	\[(x,s)\to \det (\omega_s(x))=\text{determinante da matriz que representa a forma} \]
	essa fun\c c\~ao \'e n\~ao zero em zero, assim para cada $t$ existe uma vizinhança onde ela n\~ao \'e zero. Logo, pela compacidade de $[0,1]$, $ \exists $ uma vizinhança $B\ni x_0$ onde $\det (\omega_s(x))$ n\~ao se anula $\forall s\in [0,1]$.
\end{proof}

Ent\~ao j\'a temos essa vizinhança que precisavamos.

Defina $X_t$ como a solu\c c\~ao da equa\c c\~ao de Moser:
\[i_{X_t}\omega_t=-\beta.\]
Como $\beta|_{0}=0\implies X_t|_{x=0}=0\implies \exists \varphi_t$, $t\in [0,1]$.

Pelo lema 1, existe uma vizinhança $V\ni 0$ e 
\begin{align*}
	\varphi_t: V &\longrightarrow B \\
	\varphi_t^*\omega_t&=\omega_0 
\end{align*}
tome $t=1$, $0\in U=\varphi_1(V)$.
\end{proof}

Com esse mesmo m\'etodo a gente consegue provar uma generaliza\c c\~ao do teorema de Darboux.

\subsection{Teorema de Darboux generalizado (Weinstein)}

\begin{thm}
	$Q\overset{i}{\hookrightarrow}M$ subvariedade (mergulhada) e $\omega_0,\omega_1$ em $M$ simpl\'eticas. Suponha que
	\[\omega_0|_{x}=\omega_1|_{x}\qquad \forall x\in Q\]
	ent\~ao existem vizinhanças $U_0$ e $U_1$ de $Q$ em M e um difeomorfismo
	\[\varphi:U_0\overset{\sim}{\longrightarrow}U_1  \]
	tal que 
	\[\varphi^*\omega_1=\omega_0\]
	e que $\varphi(x)=x\;\forall x\in Q$
\end{thm}

\begin{remark}
	O teorema de Darboux \'e quando $Q$ \'e um ponto s\'o!
\end{remark}

\begin{remark}
	A condi\c c\~ao $\omega_0|_{x}=\omega_1|_{x}$ significa que $\omega_0$ e $\omega_1$ coincidem em todo o espaço tangente a $M$ nos pontos de $Q$, n\~ao \'e que o pullback em $Q$ coincide. Tem mais vetores no espaço tangente a $M$.
\end{remark}

	Vamos precisar de um Lema de Poincar\'e relativo.

	\begin{lemma}
		$Q\hookrightarrow M$ subvariedade. Seja $\eta \in\Omega^{k}(M)$, $d\eta=0$,  $i^*\eta=0$. Ent\~ao existe uma vizinhança $U$ de $Q$ em $M$, $\beta\in\Omega^{k}(U)$ tal que
		\begin{align*}
			\eta & =d\beta\\
			\beta|_{x}&=0,\quad \forall x\in Q
		\end{align*}
		($\beta|_{T_xM}=0$ $\forall x\in Q$).

		A ideia aqui \'e simplesmente que podemos achar uma vizinhança de $Q$  que se contrae a $Q$ (retrato por deforma\c c\~ao?)
	\end{lemma}

\begin{proof}[Demostra\c c\~ao]
	Em fim, pelo lema, para $\eta=\omega_1-\omega_0$, $i^*\eta=0$. Compare com a demostra\c c\~ao anterior, $\beta$ se anulava no 0, agora $\eta$ se anula em toda $Q$ (\'e uma vers\~ao param\'etrica disso).

	$Q\hookrightarrow M$ tem vizinhança $U$ onde $\exists \beta\in\Omega^{1}(U)$,
	\[\omega_1-\omega_0=d\beta,\quad \beta |_{x}=0\]
	\begin{itemize}
		\item Seja $\omega_t=(1-t)\omega_0+t\omega_1=\omega_0+t d\beta$.
		\item $\forall t\in [0,1]$, $x\in Q$, $\omega_ t |_{x}=\omega_0|_{x}=\omega_1|_{x}$.
	\end{itemize}
	Pelo lema 2, $x$ tem vizinhança em M onde $\omega_t$ \'e simpl\'etica $\forall t\in [0,1]$.

	Tomando a uni\~ao das vizinhanças, temos vizinhança de $Q$ onde $\omega_t$ simpl\'etico $\forall t\in[0,1]$.

	\textbf{M\'etodo}
	\begin{itemize}
		\item Define $X_t$ por $i_{X_t}\omega=-\beta$. Isso implica que $\frac{d}{dt}\varphi^*_t\omega_t=0$.

		\item Como $\beta|_{x}=0$, ent\~ao $\forall x\in Q$, $X_t |_{x}=0\;\forall x\in Q$.

		\item Pelo lema 1, $\exists U_0$ onde $\varphi_t$ est\'a definido $\forall t\in [0,1]$.

		\item E mais $X_t |_{Q}=0\implies \varphi_t |_{Q}=\operatorname{id}_Q$.

		\item Tome $\phi=\varphi_1$ e $U_1=\varphi_1(U_0)$.
	\end{itemize}
\end{proof}

\subsubsection{Sobre o Lema de Poincar\'e relativo}

O principal ingrediente \'e teorema da vizinhança tubular.

Lembre:

\begin{thm}[Vizinhança tubular]
	$Q\hookrightarrow M$ subvariedade mergulhada. Existe uma vizinhança $Q\subseteq U\subseteq M$ para qual existe $\pi:U\to Q$ tal que 
	\begin{align*}
		\pi\circ i&=\operatorname{id}_Q\\
		i\circ \pi&\simeq \operatorname{id}_U, \quad \text{(homotop\'ia suave)} 
	\end{align*}
\end{thm}

Da\'i, o lema de Poincar\'e segue a existencia de um \textit{\textbf{operador de homotop\'ia}}.

Em geral, quando temos uma homotop\'ia
\begin{align*}
	F: M\times [0,1] &\longrightarrow N \\
	F_0 :M&\to N\\
	F_1:M&\to N
\end{align*}
exsite um operador 
\[H:\Omega^{k}(M) \to \Omega^{k-1}(M)\]
tal que
\[F_1^*\eta-F^*_0\eta=d(H\eta)-Hd\eta\]

Note que no caso de formas fechadas, o termo da direita se anula e a gente prova a invariança homot\'opica da cohomologia. No nosso caso, o operador de homotop\'ia nos da $\eta=dH\eta$j \'a que $d\eta$ se anula em  $Q$.

\subsubsection{Vizinhança tubular}

\begin{thm}
	Existe uma vizinhança $U_0$ de $Q$ em $NQ$ e uma vizinhança $U_1$ de $Q$ em $M$ tais que
	\begin{enumerate}[label=\alph*.]
		\item $U_0\cap (NQ)_x$ \'e convexo $\forall x\in Q$.

		\item Existe um difeomorfismo  $\phi:U_0\overset{\sim}{\longrightarrow} U_1$ tal que $\phi(x)=x$, e $d\phi|_{x}:T_x(NQ)\overset{\operatorname{id}}{\longrightarrow}TM|_{x}$
	\end{enumerate}
\end{thm}

\begin{proof}[Demostra\c c\~ao]
	Idea: aplica\c c\~ao exponencial.
\end{proof}

\subsection{Monitoria 2}

\begin{prop}
	$\phi:M^{2n}\to \mathbb{R}^{k}$ suave, $c \in\mathbb{R}^{k}$ valor regular.
	 \[N:=\phi^{-1}(c) \text{ coisotr\'opica }\iff \{\phi_i,\phi_j\} |_{N}=0 \]
\end{prop}


\section{Aula 10}

Lembre

\begin{itemize}
	\item Darboux generalizado: duas formas numa subvariedade que coinciden nos pontos da subvariedade, existem vizinhanças da subvariedade simplectomorfas.

	\item A prova disso: usa m\'etodo de Moser.  Para usar o m\'etodo de Moser:
		\begin{lemma}[Poincar\'e relativo]
			$Q \overset{i}{\longrightarrow}M$. $\eta\in\Omega^{k}(M)$ fechada e tal que $i^*\eta=0$. Ent\~ao existe  vizinhança $U\supset Q$ e $\beta\in\Omega^{k-1}(U)$ tal que $\eta=d\beta$ e $\beta|_{Q}=0$.

			(O lema de Poincar\'e usual \'e quando $Q$ \'e um ponto)
		\end{lemma}
		\begin{proof}[Demostra\c c\~ao do lema de Poincar\'e]
			A\'i tem que mergulhar $Q$ no fibrado tangente $NQ$ que \'e um fibrado que n\~ao precisa de m\'etrica para ser definido. Por\'em, na prova a gente intruduiz uma m\'etrica em $Q$ e identifica $NQ$ com  $T^\perp Q$. Da\'i usando a aplica\c c\~ao exponencial conseguimos ver que $Q$ \'e um retrato por deforma\c c\~ao de uma vizinhança dele no $M$---a exponencial \'e a ponte de $NQ$ [a $Q$.

			Isso da uma homotop\'ia
			\begin{align*}
				F_t: U_0 &\longrightarrow U_0 \\
				(x,v) &\longmapsto (x,tv)\\
				F_0&=i\circ \pi\\
				F_1&=\operatorname{id}_{U_0}
			\end{align*}
			Da\'i \'e s\'o pegar o operador de homotop\'ia
			\[\mathcal{H}:\Omega^{k}(U_0)\to \Omega^{k-1}(U_0)\]
			que \'e tal que
			\[F_1^*\eta=F_0^* \eta=\mathcal{H}(d\eta)+d(\mathcal{H}\eta)\]
			\begin{claim}
				O operador de homotop\'ia \'e
				\[H(\eta)=\int_{0}^1I^*_ti_{ \frac{\partial }{\partial t}}(F^*\eta)dt\]
				onde
				\begin{align*}
					F: [0,1]\times U_0 &\longrightarrow U_0 \\
					(t,y) &\longmapsto F_t(y)
				\end{align*}
				e
				\begin{align*}
					I_t: U_0 &\longrightarrow [0,1]\times U_0 \\
					y &\longmapsto (t,y)
				\end{align*}
				de forma que
				\[F_t=F\circ I_t\]
				\textbf{Nota\c c\~ao} Seja
				\begin{align*}
					\tau_t: \mathbb{R}\times U_0 &\longrightarrow \mathbb{R}\times U- \\
					(x,y) &\longmapsto (s+t,y)
				\end{align*}
				de forma que
				\[I_t=\tau_t\circ I_0,\quad F_t=F\circ I_t=F\circ \tau_t\circ I_0\]
				e a conta que a gente faiz \'e
				\begin{align*}
					\frac{d}{dt}F_t^*\eta&=I_0^*\frac{d}{dt}\tau_t^* (F^* \eta)\\
					&=I_0^* \tau_t^* (\mathcal{L}_{\frac{\partial }{\partial }t}F^*_{\eta}\\
					&\overset{\text{Cartan}}{=}I_0^* \tau_t^*\left( di_{ \frac{\partial }{\partial t}}F^* _\eta+i_{\frac{\partial }{\partial t}}d(F^*\eta \right) \\
					&=d\left( I_{t}^* i_{\frac{\partial }{\partial t}}F^*_{\eta}+I_t^ki_{\frac{\partial }{\partial t}}F^*(d\eta) \right) 
				\end{align*}
				e a\'i integramos para obter
				\[F_1^*\eta-F^*_0\eta=d(H\eta)+H(d\eta)\]
				Se $d\eta=0$,  $i^*\eta=0$, $\implies \eta=d(H\eta)$. Defina $\beta=H\eta$. Como  $F_t(x,0)=(x,0)\;\forall x\in Q$, assim
				\[\frac{\partial }{\partial t}F_t(x,0)=0\implies i_{\frac{\partial }{\partial t}}dF_t |_{x\in Q}=0\]
				e por fim
				\[\beta|_{x}=0.\]
			\end{claim}
		\end{proof}
\end{itemize}


\subsection{Darboux generalizado vers\~ao 2.0}
No teorema anterior pode imaginar que cada vizinhança de $Q$ \'e uma variedade diferente. Mas ent\~ao a condi\c c\~ao de que as duas formas sejam iguais acima de $Q$ j\'a n\~ao faz sentido. Precisamos de um isomorfismo simpl\'etico entre esses espaços tangentes.

\begin{thm}[Teorema de Darboux generalizado Vers\~ao 2.0]\leavevmode
	\[\begin{tikzcd}
		(M_0,\omega_0)&&(M_1,\omega_1)\\
	&Q\arrow[ul,"i_0",hook]\arrow[ur,"i_1",hook,swap]
	\end{tikzcd}\]
	Suponha que temos um isomorfismo de fibrados simpl\'ecticos
	\[\begin{tikzcd}
	TM_0|_{Q}\arrow[rr,"\phi"]\arrow[dr,swap]&&TM_1|_{Q}\arrow[dl]\\
	&Q
	\end{tikzcd}\]
tal que $\phi|_{TQ}:TQ\to TQ$ é $\operatorname{id}_{TQ}$.

	Ent\~ao $\phi$ estende à derivada de um simplectomorfismo
	\[\begin{tikzcd}
	U_0\subset M_0\arrow[rr,"\varphi"]&&U_1\subset M_1\\
	&Q\arrow[ul]\arrow[ur,swap]
	\end{tikzcd}\]
	i.e.,
	\[d\varphi|_{Q}=\phi:TM_0|_{Q}\to TM_1|_{Q}\]
	Em palavras: a derivada do simplectomofismo (entre as vizinhanças de $M_1$ e $M_2$) que obtemos é estendida pelo isomorfismo simpl\'etico dos fibrados tangentes dado.
\end{thm}

\begin{proof}[Demostra\c c\~ao]
	Podemos reduzir ao caso anterior! Basta achar um difeomorfismo $\psi:U_0\to U_1$ tal que $\psi|_{Q}=\operatorname{id}_Q$ e que $d\psi|_{Q}=\phi$. Nesse caso, $\omega_0$ e $\psi^*\omega_1$ s\~ao dois formas em $U_0$ que coincidem sobre $TM_1|_{Q}$. Vamo l\'a

	Tome dois complementos
	 \begin{align*}
		E_0,\quad TM_0|_{Q}&=TQ\oplus E_0\\
		E_1,\quad TM_1|_{Q}&=T_Q\oplus E_1
	\end{align*}
	Ent\~ao como $\phi$ preserva $T_Q$, ele tamb\'em preserva os complementos, \'e s\'o algebra linear. Isto \'e, $\phi$ se restringe a um isomorfismo
	\[\bar{\phi} :E_0\to E_1\]
	Note que
	\[ \bar{\phi}|_{Q} :TE_0\cong TQ\oplus E_0\to TE_1\cong TQ\oplus E_1\]
	Aqui estamos pegando a derivada do isomorfismo nos fibrados. O importante e que como ele \'e linear, sua derivada \'e ele mesmo (s\'o que a\'i aparecem muitas identifica\c c\~oes):
	\[d\bar{\phi}|_{Q}=\operatorname{id}\oplus \bar{\phi} =\phi\]
Agora pegamos vizinhanças tubulares de $Q$, $V_0\subset E_0$ e $V_1\subset E_1$ e usando a exponancial como antes podemos contraer
\[\begin{tikzcd}
	V_0\arrow[r,"\phi"]\arrow[d,"\phi_0{\color{magenta}= \operatorname{exp}}"]&V_1\arrow[d,"\phi_1"]\\
	U_0\subset M_0\arrow[r,"\psi",dashed]&U_1\subseteq M_1
\end{tikzcd}\]
e todo comuta:
\[\psi=\phi_1\circ \bar{\phi} \circ \phi^{-1}:U_0\overset{\cong }{\longrightarrow} U_1\]
e por fim
\[d\psi|_{Q}=\operatorname{id}\circ d \bar{\phi} \operatorname{id}=\phi\]
\end{proof}

Agora um caso particular:
\[\begin{tikzcd}
	(M,\omega)&&T^*\mathcal{L}\\
&\mathcal{L}\arrow[ul,hook,"i"]\arrow[ur,"i",swap,hook]
\end{tikzcd}\]
onde $Q$ est\'a dentro do fibrado cotangente como a se\c c\~ao zero.

\subsection{Teorema das vizinhanças Lagrangianas de Weinstein}


\begin{thm}[das vizinhanças Lagrangianas de Weinstein]\leavevmode
	(As subvariedades Lagrangianas est\~ao definidas "intrinsecamente", pois existe uma vizinhança delas que \'e simplectomorfa a ela como subvariedade no cotangente dela)

	Existem vizinhanças $U_0\supseteq \mathcal{L}$ em $T^*\mathcal{L}$ e $U_1\supseteq \mathcal{L}$ em $M$ e um simplectomorfismo
	\[\varphi:U_0\to U_1\]
\end{thm}

\begin{proof}[Demostra\c c\~ao]
	S\'o precisamos de um $\phi$ como no Darboux 2, i.e.,
	\begin{align*}
		\phi: TM|_{\mathcal{L}} &\longrightarrow T(T^*\mathcal{L})|_{\mathcal{L}}
	\end{align*}
	tal que
\[\phi|_{T\mathcal{L}}:T\mathcal{L}\to T\mathcal{L}\quad =\operatorname{id}_{T\mathcal{L}}\]

\begin{lemma}
	Suponha que $\mathcal{L}\hookrightarrow (M,\omega)$ \'e Lagrangiana. Considere $TM|_{\mathcal{L}}$ um fibrado vectorial simpl\'etico. Ent\~ao
	\begin{enumerate}
		\item Existe um subfibrado lagrangiano $E\subseteq TM|_{\mathcal{L}}$ tal que $TM|_{\mathcal{L}}=T\mathcal{L}\oplus E$.

		\item Existe um isomorfismo
			\[\begin{tikzcd}
			TM|_{\mathcal{L}}\arrow[rr,"\cong "]\arrow[dr,swap]&&T\mathcal{L}\oplus (T\mathcal{L})^*\arrow[dl]\\
			&\mathcal{L}
			\end{tikzcd}\]
			onde no espaço $T\mathcal{L}\oplus (T\mathcal{L})^*$ a forma simplética \'e,
			\[\nu((X,\alpha),(Y,\beta))=\beta(X)-\alpha(Y),\]
			lembrando um exerc\'icio da lista 1 (de \'alebra linear) que diz que um subespaço Lagrangiano nos da uma descomposi\c c\~ao em soma direta usando o seu dual.

			\begin{proof}[Demostra\c c\~ao do Lema]\leavevmode 

				\begin{enumerate}[label=\textbf{Step \arabic*}]
					\item Todo espaço simpl\'etico induiz uma estrutura complexa compat\'ivel.  Se $L$ \'e lagrangiano, $JL$ tamb\'em e o espaço vetorial (acho que isso coincide com o complemento ortogonal na m\'etrica compat\'ivel). Isso vale para fibrados vetorias.

					\item Note que 
						\begin{align*}
							E &\longrightarrow (T\mathcal{L})^* \\
							u &\longmapsto \omega(\cdot ,u)
						\end{align*}
						\'e um isomorfismo. Isso \'e super elementar de algebra linear.

Tome 
\begin{align*}
	(TM|_{\mathcal{L}},\omega)  &\longrightarrow (T\mathcal{L}\oplus T\mathcal{L}^*,\nu) \\
	(x,u) &\longmapsto (X,\omega(\cdot ,u)
\end{align*}
que \'e que acontece? Ent\~ao,
\begin{align*}
	\nu(T(X,u),T(Y,v)&=\nu((X,\omega(\cdot ,u)),(Y,\omega(\cdot ,v))\\
	&=\omega(X,u)-\omega(Y,u)\\
	&=\omega((X,u),(Y,u))
\end{align*}
				\end{enumerate}
			\end{proof}
	\end{enumerate}
	Da\'i, o teorema queda provado simplesmente aplicando o lema em $TM|_{\mathcal{L}}$ e em $T(T^*\mathcal{L})|_{\mathcal{L}}$ para construir o isomorfismo de fibrados que buscavamos:
	\[\begin{tikzcd}
	&T\mathcal{L}\oplus (T\mathcal{L})^*\arrow[dr,"\cong "]\\
	TM|_{\mathcal{L}}\arrow[rr,"\phi"]\arrow[ur,"\cong "]&&T(T^*\mathcal{L})|_{\mathcal{L}}
	\end{tikzcd}\]
\end{lemma}
\end{proof}



\section{Aula 11}

\subsection{Aplica\c c\~ao (de Weinstein): pontos fixos de simplectomorifsmos}

Generaliza o estudo (Poincar\'e-Birkoff) cl\'assico de pontos fixos de aplica\c c\~oes que preservan \'area:

\begin{thm}[\'Ultimo teorema de Poincar\'e]\leavevmode
	Um automorfismo de um anelo que preserva orianta\c c\~ao, \'area e rota a fronteira do anelo em dire\c c\~oes opostas tem um ponto fixo.
\end{thm}

Isso apareceo quando Poincar\'e estudava fluxos em $\mathbb{R}^{3}$.

Consideremos $(M,\omega)$ simpl\'ectica e $M \overset{f}{\longrightarrow}M$ simplectomorfismo. Nos interessa o caso em que $f$ \'e um fluxo hamiltoniano no tempo 1, ie. $f=\varphi^{t=1}_{X_{H_t}}$. Sabemos que
\[\Gamma_f=\{(x,f(x)):x\in M\} \subseteq M\times \bar{M} \]
\'e uma subvariedade lagrangiana, e tamb\'em
$\Delta =\Gamma_{\operatorname{id}_M}=\{(x,x):x\in M\} \subseteq M\times \bar{M}$ 
De forma que os pontos fixos de $f$ s\~ao os pontos de interse\c c\~ao entre $\Gamma_f$ e $\Delta$.

\begin{prop}
	Seja $M$ compacta, $H^{1}_{\operatorname{dR}}(M)=0$. Se $f$ \'e  \textit{\textbf{$C^1$-pr\'oximo}} (convergencia uniforme, \href{https://en.wikipedia.org/wiki/Compact-open_topology#Fr%C3%A9chet_differentiable_functions}{Fr\'echet differentiable}?) da $\operatorname{id}_M$, ent\~ao $f$ tem pelomenos 2 pontos fixos.
\end{prop}

\begin{proof}[Demostra\c c\~ao]
	 Note que $\Delta \cong M$ pelo teorema da vizinhança lagrangiana, como $\Delta$ \'e lagrangiana existe uma vizinhança $U\supseteq \Delta$ simplectomorfa a uma vizinhança $U'$ de  $M\hookrightarrow (T^* M,\omega_{\operatorname{can}}$.

	 \begin{itemize}
	 \item Se $f\in\operatorname{Simp}$ est\'a "perto" da $\operatorname{id}_M$, ent\~ao $\Gamma_f\subseteq U$.

	\item $f$ \'e $C^1$-pr\'oximo da $\operatorname{id}_M$, ent\~ao $\Gamma_f$ corresponde a 1-forma $\mu$ em $T^*M$ (a uma subvariedade $N_\mu$ de  $T^*M$?). (\'E uma gr\'afica de $M$ no fibrado cotangente!)

	\item $\Gamma_f$ lagrangiana  $\implies $ $d\mu=0$ (Lista 2)

	 \item $H^{1}(M)=0\implies \mu=dh$
	
	\item $M$ compacta $\implies $ $h$ tem pelo menos 2 pontos cr\'iticos.
	 \end{itemize}
\end{proof}

\begin{remark}[Monitoria]
	Se pedimos s\'o $C^0$-pr\'oximo, \'e poss\'ivel que a se\c c\~ao $\mu$ n\~ao esteja bem definida porque um ponto de $M$ pode n\~ao estar associado a um covector ancorado em outro ponto, ou algum outro problema assim. A condi\c c\~ao $C^1$ controla isso. 
\end{remark}

\begin{remark}\leavevmode 
	\begin{itemize}
	\item N\~ao podemos abrir m\~ao de $H^{1}(M)=0$. Eg. rota\c c\~ao no toro.

	\item Podemos substituir $H^{1}(M)=0$ por $f$ ser simplectomorfismo Hamiltoniano (ver McDuff-Salomon).
	\end{itemize}
\end{remark}

\begin{question}
	Remover $C^1$-proximidade da identidade? (Pelo menos no caso $f$ hamiltoniano.
\end{question}

\begin{conjecture}[Arnold]
	$M$ simpl\'etica compacta, $f$ simplectomorfismo Hamiltoniano. O n\'umero de pontos fixos de $f$ e maior o igual que o n\'umero m\'inimo de pontos cr\'iticos que uma fun\c c\~ao em $M$ deve ter:
	\begin{align*}
		\text{\# pontos fixos de $f$}&\geq \operatorname{Crit}(M)\\
		\geq  \text{LS category (Lusternik Schninelmann} 
	\end{align*}
Isso est\'a relacionado com o fato de tirar a hip\'otese de que a fun\c c\~ao est\'e pr\'oxima da identidade.
\end{conjecture}

\begin{conjecture}[Outra vers\~ao]
	Para pontos fixos n\~ao degenerados (s\~ao os pontos onde $N_\mu$ e  $M$ se intersectan transversalmente em $T^*M$).
	\begin{align*}
		\text{\# pontos cr\'iticos} &\geq \text{\# m\'inimo de pontos cr\'iticos que fun\c c\~oes de Morse devem ter.}\\
		&\underbrace{\geq }_{\text{desig. Morse} } \sum_{k}\operatorname{Betti}_{k}
	\end{align*}
\end{conjecture}

\paragraph{Projetos} 
\begin{itemize}
\item Conjetura de Arnold (Eliashbag (superf\'icies de Riemann), Hofer-Achander)
\item Homologia de Floer (\'e uma vers\~ao de Homologia de Morse em dimens\~ao infinita)
\end{itemize}

Professor Leonardo vai falar com mais detalhe desses temas.

\subsection{Outra classe de exemplos de variedades simpl\'eticas: \'orbitas coadjuntas}

S\~ao exeplos de redu\c c\~ao simpl\'etica. Isso est\'a relacionado com teor\'ia de Lie. $G \curvearrowright (M,\omega)$ simetr\'ias hamiltonianas. Da\'i vamos produzir uma nova variedade simpl\'etica.

\subsubsection{Revis\~ao de grupos e \'algebras de Lie}

Cada grupo de Lie age na sua \'algebra de Lie de maneira can\'onica. Da\'i podemos pegar a \'algebra dual. O fato importante \'e que as \'orbitas l\'a tem uma estrutura simpl\'etica.

\begin{defn}
	Um \textit{\textbf{grupo de Lie}} \'e uma variedade $C^\infty$ $G$ munida de estrutura de grupo tal que o produto e a invers\~ao s\~ao fun\c c\~oes  suaves. Os \textit{\textbf{morfismos}} s\~ao homomorfismos de grupos $C^\infty$. \textit{\textbf{Subgrupos de Lie}} s\~ao subvariedades imersas que s\~ao subgrupos.
\end{defn}

\begin{example}\leavevmode
	\begin{itemize}
	\item $\operatorname{GL}(n,\mathbb{R})=\{A\in M_{n\times n}(\mathbb{R}):\det A \neq 0\}$ 
	\item $V$ espaço vetorial, $(V,+)$  \'e grupo de Lie abeliano.
	\item  $S^1$, $S^1\times S^1\times \ldots\times S^1$ s\~ao grupos de Lie abelianos.
	\item Grupos finitos/enumer\'aveis: $\mathbb{Z},\mathbb{Z}_m,\ldots$ 
	\item 
		\begin{exercise}
			$G$ grupo de Lie conexo, ent\~ao o seu recobrimento universal $\tilde{G}$ \'e grupo de Lie.
		\end{exercise}
	\item Subgrupos de $\operatorname{GL}(n,\mathbb{R})$:
		\begin{itemize}
		\item Ortogonal $\operatorname{O}(n)=\{A\in\operatorname{GL}(n,\mathbb{R}):A A^{\mathbf{T}}=\operatorname{id}\}$.

			Mais generalmente, $(V,\left<\cdot,\cdot\right> $ espaço vetorial de produto interno, $\operatorname{O}(V) =\{T:V\to V|\left<Tv,Tv\right> =\left<u,v\right> \}$.

			Considerando
			\begin{align*}
				\psi:\operatorname{GL}(n,\mathbb{R}) &\longrightarrow \operatorname{Sim}(n) \\
				A &\longmapsto A A^{\mathbf{T}}
			\end{align*}
			temos que $\operatorname{id}$ \'e um valor regular, e assim $\operatorname{O}(n)=\psi^{-1}(\operatorname{id})$ \'e uma subvariedade (compacta \'e n\~ao conexa por $\det A=\pm 1$
	
	\item $\operatorname{SL}(n) =\{A\in\operatorname{GL}(n,\mathbb{R}):\det A=1\}$, conexo n\~ao compacto.

	\item $\operatorname{SO}(n) =\operatorname{O}(n) \cap \operatorname{SL}(n)$ compacto conexo

	\item $\operatorname{Sp}(2n) =\{A\in\operatorname{GL}(2,\mathbb{R}):A^{\mathbf{T}}J_0A=J_0\}$ com $J_0=\begin{pmatrix} 0&-I\\I&0 \end{pmatrix} $. \textit{\textbf{Grupo simpl\'etico}}.

	\item $\operatorname{GL}(n,\mathbb{C})=\{A\in M_{n\times n}(\mathbb{C}):\text{invertiveis} \} \overset{\text{aberto} }{\subseteq}M_n(\mathbb{C})=\mathbb{C}^{n^2}$.

	\item $\operatorname{U}(n) =\{A\in\operatorname{GL}(n,\mathbb{C}):A A^*=\operatorname{id}\}$, issto \'e $A^*=\overline{A}^{\mathbf{T}}$, temos $|\det A| =1$ e o mapa
		\[\det :\operatorname{U}(n)\to S^1\]
		e de fato $\operatorname{U}(1) \cong S^1$.

	\item $\operatorname{SU}(n) = \{A\in\operatorname{U}(n) :\det A=1\}$ \textit{\textbf{grupo unit\'ario especial}}
		
		\end{itemize}
	\end{itemize}	
\end{example}

\begin{remark}\leavevmode 
	\begin{itemize}
	\item 	\begin{thm}[de Cartan]\leavevmode
		(F. Warner) Subgrupo fechado de grupo de Lie \'e subgrupo de Lie! (mergulhado)
	\end{thm}

	\item Nem todo grupo de Lie \'e grupo de Lie de matrices. O espaço recobridor de $\operatorname{SL}(2,\mathbb{R})$, por exemplo.
	\end{itemize}
\end{remark}

\subsubsection{Sobre $\operatorname{SU}(2)$}

Sabemos que podemos escrever 
\[\mathbb{R}^{4}=\mathbb{H}=\{a+ib+jc+kd:i^2=j^2=k^2=-1\} \]
e como
\[S^3\hookrightarrow \mathbb{H}\]
$S^3$ herda uma estrutura de grupo de Lie, e de fato
\[S^3\overset{\cong }{\longrightarrow}\operatorname{SU}(2) = \left\{ (\frac{\alpha}{-\beta}\frac{\beta}{\alpha}:\alpha,\beta\in\mathbb{C},|\alpha|^2+|\beta|^2=1\right\} \]
Da\'i,
\[\begin{tikzcd}
	S^3\arrow[r,"\cong "]\arrow[d,"2:1",swap]&\operatorname{SU}(2)\arrow[d,"2:1"]\\
	\mathbb{R}P^{3}\arrow[r,"\cong "]&\operatorname{SO}(3)
\end{tikzcd}\]
Em geral recobrimentos duplos de $\operatorname{SO}(n)$ s\~ao grupos $\operatorname{Spin}(n)$. Para $n\geq 2$ s\~ao recobrimentos universais. S\~ao grupos de simetr\'ias de particulas que se chaman fermiones. A ideia \'e que a gente precisa dois voltas para virar a flecha que t\'a parada na particula.

Por \'ultimo vamos ver por qu\'e e que $\mathbb{R}P^{3}\cong \operatorname{SO}(3)$. Do mesmo jeito que $\mathbb{R}P^{2}$ \'e o hemisferio norte da esfera com os pontos no bordo identificados, $\mathbb{R}P^{3}$ \'e uma bola fechada em $\mathbb{R}^3$ com pontos antipodais no bordo identificados. Podemos pensar que os pontos de $\mathbb{R}P^{3}$ s\~ao rota\c c\~oes  de \'angulo $\pi\sqrt{x^2+y^2+z^2}$.


\section{Aula 12}

\subsection{\'Algebras de Lie}

\begin{defn}
	Uma \textit{\textbf{\'algebra de Lie}} \'e um espaço vetorial (sobre $\mathbb{R}$) munido de uma forma bilinear 
	\begin{align*}
		[\cdot,\cdot ]: \mathfrak{g} \times \mathfrak{g} &\longrightarrow \mathfrak{g}
	\end{align*}
tal que
\begin{itemize}
	\item $[u,v]=-[v,u]$.
	 \item Jacobi.
\end{itemize}

Um \textit{\textbf{morfismo de \'Algebras de Lie}} \'e um mapa linear $T:\mathfrak{g} \to \mathfrak{g}$ tal que $[Tu,Tv]=T[u,v]$.

Uma  \textit{\textbf{sub\'algebra de lie}} \'e $\mathfrak{h} \subset \mathfrak{g}$ subespaço tal que $[\mathfrak{h},\mathfrak{h} ]\subseteq\mathfrak{h}$, ie. $[\mathfrak{g},\mathfrak{h} ]\subseteq \mathfrak{h}$ ideal.
\end{defn}

\begin{example}\leavevmode 
	\begin{itemize}
		\item $M_n(\mathbb{R})$ matrizes de $n\times n$ com $[A,B]=AB-BA$. Se denota  $\mathfrak{gl}(n)$.

		\item $V$ espaço vetorial com $[\cdot,\cdot ]\equiv_0$ (abelianos).

		\item $\mathfrak{o}(n) =\{A\in M_n(\mathbb{R}):A=-A^{\mathbf{T}}\}$.

		\item $\mathfrak{sl}(n) =\{A\in M_n(\mathbb{R}):\operatorname{tr}A=0\}$.

		\item Em dimens\~ao $\infty$, $(\mathfrak{X}(M),[ \cdot,\cdot ])$ campos vetoriais, e tamb\'em $(\mathcal{C}^\infty(M),\{\cdot,\cdot\} )$ simpl\'etica.
	\end{itemize}
\end{example}

\subsection{Rela\c c\~ao entre grupos de Lie e álgebras de Lie}
\c c\~ao
\[\begin{tikzcd}
	\text{Grupos de lie}\arrow[rr,"\text{diferenciação}" ]&&\text{\'Algebras de Lie}  
\end{tikzcd}\]

\begin{defn}
	Dado um grupo de Lie $G$ e um elemento $g\in G$, a \textit{\textbf{multiplica\c c\~ao à esquerda}} \'e
	\begin{align*}
		L_g: G &\longrightarrow G \\
		a &\longmapsto ga
	\end{align*}
que \'e um difeomorfismo. Tamb\'em est\'a a \textit{\textbf{multiplica\c c\~ao à direita}} $R_g$. Temos que  $(L_g)^{-1}=L_{g^{-1}}$.
\end{defn}

\begin{defn}
	Um campo de vetores $X\in\mathfrak{X}(G)$ \'e \textit{\textbf{invariante à esquerda}} se  \[(L_{g})_*X=X \qquad \text{ para toda }g \in G\]
	ou
	\[(dL_g)|_{h}(X_h)=X_{gh}\qquad \forall g,h\]
\end{defn}

\begin{remark}\leavevmode 
	\begin{itemize}
	\item O conjunto de campos invariantes à esquerda $\mathfrak{X}^L(G) \subseteq \mathfrak{X}(G)$ \'e uma sub\'algebra de Lie (o colchete de Lie \'e fechado aqui).
	\item $X\in\mathfrak{X}^L(G)$ \'e completamente determinado por  $X_e$ onde  $e\in G$ \'e a unidade.

	\item Temos um isomorfismo de espaços vetoriais
		\begin{align*}
			\mathfrak{X}^L(G)  &\overset{ \cong }{\longrightarrow}T_eG\\
			X &\longmapsto X_e
		\end{align*}
		e para voltar s\'o pegamos um vetor em  $T_eG$ e espalhamos por todos lados:
\begin{align*}
	T_eG  &\longrightarrow \mathfrak{X}^L(G) \\
	u &\longmapsto u^L
\end{align*}
		Talvez aqui devemos provar que ese espalhamento produz um campo suave.

		Al\'em disso, esse mapa induz uma estrutura de \'algebra de Lie em $T_eG$. Essa  \'e a \textit{\textbf{\'algebra de Lie}} de $G$, denotada por  $\operatorname{Lie}(G)$.
			\end{itemize}
\end{remark}

\paragraph{Extra:} $(\mathfrak{X}(M),-[\cdot,\cdot ])$ \'e a \'algebra de Lie de $\operatorname{Diff}(M)$.

\paragraph{Importante:} Essa associa\c c\~ao \'e funtorial:

\begin{prop}
	Se $\phi:G_1\to G_2$ \'e um morfismo de grupos de Lie, ele induiz uma aplica\c c\~ao entre as algebras de Lie dado por
	\[\operatorname{Lie}(\phi):=d_e\phi:T_eG_1\longrightarrow T_eG_2\]
	que \'e um morfismo de \'algebras de Lie.
\end{prop}

\begin{proof}[Demostra\c c\~ao]
	Muito f\'acil.
\end{proof}

Esse \'e chamado de \textit{\textbf{funtor de Lie}}.  \'E um funtor de diferencia\c c\~ao nesse sentido. Como ir na outra dire\c c\~ao?

\begin{remark}\leavevmode 
	\begin{itemize}
	\item $V$ espaço vetorial, $G\longrightarrow \operatorname{GL}(V)$ morfismo de grupos de Lie se chama de \textit{\textbf{representa\c c\~ao}} de $G$ em $V$.
	\item $\mathfrak{g} \longrightarrow \mathfrak{gl}(V)$ morfismo de \'algebras de Lie se chama de \textit{\textbf{representa\c c\~ao}} de $\mathfrak{g}$ em $V$.
	\end{itemize}
\end{remark}

\begin{example}\leavevmode 
	\begin{itemize}
	\item $\mathfrak{gl}(n,\mathbb{R}) =\operatorname{Lie}(\operatorname{GL}(n,\mathbb{R}))$.
	\item $(V,[\cdot,\cdot ]=0)=\operatorname{Lie}(V,+)$.

	\item Pode ter grupos de Lie diferentes com a mesma \'algebra de Lie, por exemplo $(\mathbb{R},+)$ e $S^1$ tem \'algebra de Lie $(\mathbb{R},[\cdot,\cdot ])$. Outro exemplo \'e $\mathfrak{o}(n)= \operatorname{Lie}(\operatorname{O}(n)$, e como esse grupo de Lie tem dois componentes conexas (det 1 e det -1), e a \'algebra de Lie est\'a determinada s\'o pela componente da identidade, temos que $\mathfrak{o}(n)=\mathfrak{so}(n)= \operatorname{Lie}(\operatorname{SO}(n))$.

	\item $\mathfrak{so}(3) \cong \mathbb{R}^{3}$ com
		\[\begin{pmatrix} 0&a &-b\\-a &0&c\\b-c &0 \end{pmatrix} \longrightarrow \begin{pmatrix} a \\b\\ c \end{pmatrix} \]

	\item Recobrimentos $2:1$.

	 \item Na lista 1 vimos que $\mathfrak{sp}(2n)= \{A\in M_n(\mathbb{R}):A^{\mathbf{T}}J_0+J_0A=0\} =\operatorname{Lie}(\operatorname{Sp}(2n) )$ 

	\item $\operatorname{Lie}(\operatorname{GL}(n,\mathbb{C})) =M_n(\mathbb{C})$.

	\item $\operatorname{Lie}(\operatorname{SL}(n,\mathbb{C})) =\mathfrak{sl}(n,\mathbb{C})=\{A\in M_n(\mathbb{C}):\operatorname{tr}A=0\}$.
	\item $\operatorname{Lie}(\operatorname{U}(n)) =\mathfrak{un}(n) =\{A\in M_n(\mathbb{C}):A=-A^*\}$.
	\item $\operatorname{Lie}(\operatorname{SU}(n)) =\mathfrak{su}(n) =\{A\in M_n(\mathbb{C}):\operatorname{tr}A=0,A=-A^*\}$
	\end{itemize}
\end{example}

\begin{exercise}\leavevmode 
	\begin{itemize}
	\item $u\in\mathfrak{g}$, $u^L\in\mathfrak{X}^L(G)$ \'e completo. Aqui \'e elongar cualquer curva integral usando a trasla\c c\~ao a esquerda.

	Seja $\gamma_e:\mathbb{R}\to G$ curva integral de $u^L$ com  $\gamma_u(0)=e$. 

	\item $\gamma_u(t+s)=\gamma_u(t)\gamma_u(s)$ ($\gamma_u$  \'e um homomorfismo de grupos)

	\item $\gamma_{tu}(1)=\gamma_{u}(t)$ (homogenidade)
	\end{itemize}
\end{exercise}


\subsection{Exponencial}

\begin{defn}
	\begin{align*}
		\operatorname{exp}:\mathfrak{g}  &\longrightarrow G \\
		u &\longmapsto \gamma_u(1)
	\end{align*}
\end{defn}

\begin{remark}
\[\operatorname{exp}((t+s) u)=\operatorname{exp}(tu) \operatorname{exp}(su) \]
\[\operatorname{exp}(-tu) =\operatorname{exp}(tu)^{-1}\]
\end{remark}

\begin{remark}\leavevmode 
	\begin{itemize}
	\item Fluxo de $u^L$ \'e $R_{\operatorname{exp}(tu)}$ :
		\begin{align*}
			\frac{d}{dt}R_{\operatorname{exp}(tu)}(g)&=\frac{d}{dt}\Big|_{t=0}(g\operatorname{exp}(tu) )\\
			&=\frac{d}{dt}\Big|_{t=0}L_g\operatorname{exp}(tu)\\
			 &=\frac{d}{dt}|_{t=0}L_g(u)\\
			 &=u^L|_{g}
		\end{align*}
	\item Fluxo de $u^R$ \'e $L_{\operatorname{exp}(tu)}$.
	\end{itemize}
\end{remark}

\subsection{Propriedades fundamentais}

\begin{enumerate}
	\item $\operatorname{exp}: \mathfrak{g}\to G$ \'e um difeomorfismo entre vizinhançasde $0\in\mathfrak{g}$ e de $e\in G$. ($d_{0}\operatorname{exp}= \operatorname{id}$.

	\item $\phi:G\to H$ morfismo de grupos de Lie, $\Phi=\operatorname{Lie}(\phi) :\mathfrak{g} \to  \mathfrak{h}$. O diagrama comuta:
\[\begin{tikzcd}
	G\arrow[r,"\Phi"]&H\\
	\mathfrak{g}\arrow[u,"\operatorname{exp_G}"]\arrow[r,"\phi"]&\mathfrak{h}\arrow[u,"\operatorname{exp}_H",swap]
\end{tikzcd}\]
\end{enumerate}

\begin{example}
	Para $G=\operatorname{GL}(n)$ e $\mathfrak{g} =M_n(\mathbb{R})$,
	\begin{align*}
		\operatorname{exp}: M_n(\mathbb{R}) &\longrightarrow \operatorname{GL}(n) \\
		A &\longmapsto e^A
	\end{align*}
	Note que
	\[\det :\operatorname{GL}(n)\to \mathbb{R}^*\]
	\[\begin{tikzcd}
		\operatorname{GL}(n)\arrow[r,"\det "]&\mathbb{R}^*\\
		M_n(\mathbb{R})\arrow[u,"\operatorname{exp}"]\arrow[r,swap,"\operatorname{tr}"]&\arrow[u,"\operatorname{exp}",swap]
	\end{tikzcd}\]
	Lembre que $\det(e^A)=e^{\operatorname{tr}A}$.
\end{example}

\begin{remark}
	Numa variedade Riemanniana com uma m\'etrica invariante, a defini\c c\~ao da exponencial usando geod\'esicas coincide com essa daqui.
\end{remark}

\subsection{Teoremas fundamentais de Lie (da \'algebra para o grupo)}

\begin{thm}[Teoremas fundamentais de Lie]\leavevmode
	\begin{enumerate}
		\item[Lie I] $\mathfrak{g} =\operatorname{Lie}(G)$, existe um \'unico grupo de Lie $\tilde{G}$ simplesmente conexo tal que $\mathfrak{g} =\operatorname{Lie}(\tilde{G})$. A ideia e simplificar a topologia do grupo preservando a sua \'algebra. Aqui se usa o recobrimento universal da componente conexa da identidade.
	
		\item[Lie II] $\varphi:\mathfrak{g} =\operatorname{Lie}(G)\to \mathfrak{h} =\operatorname{Lie}(H)$. Se $G$ \'e simplesmente conexo, existe um morfismo $\phi:G\to H$ tal que $\varphi =\operatorname{Lie}(\phi)$.

			\begin{example}
				O fluxo irracional  \'e um morfismo de grupos de Lie que n\~ao se factora (s\'o se factora se a \'orbita fecha)
				\[\begin{tikzcd}
				\mathbb{R}\arrow[r]\arrow[d,dashed,swap]&\mathbb{T}^2\\
				S^1\arrow[ur]
				\end{tikzcd}\]
			\end{example}

		\item[Lie III] Toda \'algebra de Lie de dimens\~ao finita \'e \'algebra de Lie de um grupo de Lie.

			\begin{thm}[Ado]\leavevmode
				Cualquer \'algebra de Lie pode ser vista de maneira fiel dentro de $\mathfrak{gl}(n)$, ie. $\mathfrak{g} \hookrightarrow \mathfrak{gl}(n)$. (isso no acontece para grupos de Lie e $\operatorname{GL}(n)$.)
			\end{thm}
			Da\'i pode usar o teorema de Frobenius para mostrar que toda sub\'algebra de Lie da um subgrupo de Lie: espalha a sub\'algebra (subespaço do tangente em $e$ ) usando multiplica\c c\~ao a esquerda, que \'e Frobenius integr\'avel porque \'e \'algebra de Lie, assim ele vem de uma distribui\c c\~ao, pega a \'orbita que pasa por $e$, essa da\'i \'e um subgrupo.
	\end{enumerate}
\end{thm}

\begin{remark}
	Em dimens\~ao infinita tem obstru\c c\~oes.
\end{remark}

\paragraph{Em resumo:}

\[\begin{tikzcd}
	\{\text{Grupos de Lie simp. conexo} \} \arrow[rr,"\text{Lie} "]&&\{\text{\'Algebras de Lie} \}
\end{tikzcd}\]
\'e uma equivalencia de categor\'ias.


\section{Aula 13}

\subsection{Ações}

\begin{defn}
	Uma \textbf{\textit{ação (à esquerda)} } de $G$ em $M$ é aplicação suave
	\begin{align*}
		\psi: G\times M &\longrightarrow M \\
		(g,x) &\longmapsto gx=\psi(g,x)= \psi_{g}(x)
	\end{align*}
	tal que
	\begin{align*}
		\psi_e(x)&=x,\qquad \forall x\in M\\
		\psi_{gh}&=\psi_g\cdot \psi_h
	\end{align*}
	 que é equivalente a que
	\[\forall g\in G, \quad \psi_g\in\operatorname{Diff}(M), \quad G\mapsto \operatorname{Diff}(M)\]
	é homeomorfismo de grupos.
\end{defn}

\paragraph{Notação} $G\overset{\psi}{\curvearrowright }M$

\begin{remark}
	Análogo para ações à direita com $\psi_{gh}=\psi_h\cdot \psi_g$.
\end{remark}

\paragraph{Terminologia} $M$ + ação por $G$ se llama \textit{\textbf{$G$-variedade}}.

\begin{defn}
	Um \textit{\textbf{morfismo de  $G$-variedades}} é
\[\begin{tikzcd}
	M_1\arrow[r,"\phi"]\arrow[d,"g",swap]&M_2\arrow[d,"g"]\\
	M_1\arrow[r,"\phi"]&M_2
\end{tikzcd}\]
ie.,
\[\phi(g\cdot x)=g\phi(x)\]
ie., \textit{\textbf{$G$-equivariante}}.
\end{defn}

\begin{example}\leavevmode 
	\begin{enumerate}
		\item $M=V$  espaço vetorial. Ações  por transformações  lineares=representações .
			\[G\overset{\psi}{\longrightarrow}\operatorname{GL}(V)\subset \operatorname{Diff}(V)\]
			tem representação dual  $G\curvearrowright V^*$ dada por
			\[G\overset{\psi^*}{\longrightarrow}\operatorname{GL}(V^*)\]
			dada por
			\[\left<(\psi^*)_g(\xi ),v\right> =\left<\xi,\psi_{g ^{-1}(v)}\right> ,\quad \xi \in V^*,v\in V\]
	
		\item $\operatorname{GL}(n)\curvearrowright \mathbb{R}^{n}$ restringe $\operatorname{O}(n)\curvearrowright S^{n-1}$.

		\item $G=\mathbb{R}\curvearrowright M$, $\mathbb{R}$-ação $\longleftrightarrow $ fluxo (completo) pois $\psi_{t+s}=\psi_t\cdot \psi_s$. Além disso, $\mathbb{R}^{n}\curvearrowright M \longleftrightarrow $ $n$-fluxos que comutam, pois
		\[\psi_{(t_1,\ldots,t_n)}=\psi_{t_1}^1\circ \ldots \circ \psi^n_{t_n}\]

	\item $G$ age em si mesmo de várias formas:
		\begin{itemize}
		\item $\psi_g=L_g$, $L_g(a)=ga$, ação por multiplicação à esquerda.

		\item $\psi_g=R_g$ à direita. (Note que  $\psi_g=R_{g ^{-1}}$ é ação à esquerda.)

		\item $\psi_g=I_g=L_g\circ R_{g ^{-1}}$, $ I_g(a)=gaq ^{-1}$, ação por cnjugação (ação por automorphismos de $G$ ).

		\item \textit{\textbf{Ação ou representação adjunta}}. $G$ age em $\mathfrak{g} =\operatorname{Lie}(G)$
			\begin{align*}
				\operatorname{Ad}: G &\longrightarrow \operatorname{GL}(\mathfrak{g})\subseteq \operatorname{Diff}(\mathfrak{g}) \\
				g &\longmapsto d_e I_g:\mathfrak{g} \to \mathfrak{g}
			\end{align*}
			
		\item \textit{\textbf{Representação dual}}. $G\curvearrowright \mathfrak{g}^*$. Usando \href{https://en.wikipedia.org/wiki/Coadjoint_representation#Formal_definition}{wiki},
			 \begin{align*}
				\operatorname{Ad}^*: G &\longrightarrow \mathsf{GL}(\mathfrak{g}^*) \\
				\Big(\operatorname{Ad}^*_g\mu\Big)x &=\mu\Big(\operatorname{Ad}_{g^{-1}}x\Big) 
			\end{align*}
			Em aula,
			\begin{align*}
				\operatorname{Ad}^*: G &\longrightarrow \operatorname{GL}(\mathfrak{g}^*) \\
				g &\longmapsto \begin{aligned}
				\psi_g:\mathfrak{g}^*\longrightarrow \mathfrak{g}^* \\
				\left<\psi_g(\mu),v\right> =\left<\mu,\operatorname{Ad}_{g ^{-1}}(v)\right> 
				\end{aligned}
			\end{align*}
			$(\operatorname{Ad}^* )_g=(\operatorname{Ad}_{g ^{-1}})^*$.


		\end{itemize}
		\item Levantamiento tangente cotangente. $G\overset{\psi}{\curvearrowright }M$.
			\begin{itemize}
			\item $G\curvearrowright TM$
				\[g\mapsto \begin{tikzcd}d\psi_g:TM\arrow[r,"\cong "]\arrow[d]& TM\arrow[d]\\ M \arrow[r,"\psi_g"]&M\end{tikzcd}\]

			\item $G\curvearrowright T^*M$
				\[g \longrightarrow \begin{tikzcd}\Big((d\psi_g){\color{3}^{-1}}\Big)^* :T^* M\arrow[r,"\cong "]\arrow[d]& T^* M\arrow[d]\\ M\arrow[r,"\psi_g"]&M\end{tikzcd}\]
			\end{itemize}

	\end{enumerate}
\subsection{Descrição infinitesimal de $G$-ações}

\begin{upshot}\leavevmode
	Tem uma ação $G\curvearrowright M$. Pegue um vetor $X$ na álgebra de Lie. $\operatorname{exp}(tX)$ da uma curva em $G$. Essa curva em $G$ da uma curva em $M$ quando começamos em qualquer ponto $p\in M$. Isso é um fluxo global porque $\operatorname{exp}((t+s)X)= \operatorname{exp}(tX)\operatorname{exp}(sX)$. Esse fluxo global da um campo vetorial.
	\[\mathfrak{g} \to \mathfrak{X}(M)\]
\end{upshot}

\paragraph{Lembre} $\mathbb{R}$ ação  $\longleftrightarrow $fluxos $\psi_t\rightsquigarrow $campos de vetores $X(x)=\frac{d}{dt}\Big|_{t=0}\psi_t(x)$

\paragraph{Generalização para $G$-ações} $G$ grupo de Lie, $\mathfrak{g} =\operatorname{Lie}(G)$, $G\overset{\psi}{\curvearrowright }M$.

\paragraph{Note}  cada $u\in\mathfrak{g}$ determina uma $\mathbb{R}$-ação = fluxo em $M$.

O que aqui acontece é que cada vetor no tangente à identidade de $G$ genera um fluxo na variedade $M$ quando $G$ age em $M$.

\[\begin{tikzcd}
	\mathbb{R}\arrow[r,"\gamma_u"]\arrow[rr, bend left]& G\arrow[r,"\psi"]&\operatorname{Diff}(M) \\
	t \arrow[rr,maps to]&&\psi_{\operatorname{exp}(tu)}
\end{tikzcd}\]

\begin{defn}
	O \textbf{\textit{gerador infinitesimal} }  de $\psi$ correspondendo a $u\in\mathfrak{g}$ é o campo
	\[u_M(x):=\frac{d}{dt}\Big|_{t=0}\psi_{\operatorname{exp}(tu)}(x)\]
\end{defn}
Em conclusão, uma $G$-ação da lugar a um mapa
\begin{align*}
	\mathfrak{g}  &\longrightarrow \mathfrak{X}(M) \\
	u &\longmapsto u_{M}
\end{align*}

Heurísticamente (só que $\operatorname{Diff}(M)$ não tem dimensão finita), tendo um mapa
\[G\to \operatorname{Diff}(M)\]
derivando obtemos
\[\mathfrak{g} \to \mathfrak{X}(M).\]

\begin{example}\leavevmode 
	\begin{enumerate}
		\item $G\overset{\psi}{\curvearrowright }G$. $\psi_g=L_g\rightsquigarrow \begin{aligned}
			\mathfrak{g}  &\longrightarrow \mathfrak{X}(G) \\
			u &\longmapsto u^R
		\end{aligned}$
	Aqui temos que fazer uma conta:
	\begin{align*}
		u_G(a)&=\frac{d}{dt}\Big|_{t=0}L_{\operatorname{exp}(tu)(a)}\\
		&=\frac{d}{dt}\Big|_{t=0}R_a(\operatorname{exp}(tu))\\
		 &=d_eR_a(u)\\
		 &=u^R(a)
	\end{align*}

	Análogamente,  outros generadores infinitesimais são

\item \[\psi_g=R_g\implies u_G=u^L\]
	\[\psi_g=R_{g ^{-1}}\implies u_G=-u^L\]
\item \[\psi_g=I_g=L_g\circ R_{g ^{-1}}\implies \begin{aligned}
	\mathfrak{g} &\longrightarrow \mathfrak{X}(G) \\
	u &\longmapsto u^R-u^L
\end{aligned} \]
	\end{enumerate}
\end{example}

\subsection{No caso de representações}
\[G \overset{\psi}{\longrightarrow}\operatorname{GL}(V)\subset \operatorname{Diff}(M)\]
diferenciando,
$\mathfrak{g} \longrightarrow \mathfrak{gl}(V) \subset \mathfrak{X}(V)$ 

A partir de uma transformação linear podemos gerar um campo vetorial:
\[\{A:V\to V \text{linear} \} =\mathfrak{gl}(V) \hookrightarrow \mathfrak{X}(V)\]
pega um vetor $v$. Em $v$, o vetor do campo vetorial  é $Av$.
$X_A(v)=Av\in T_vV=V$

\subsection{Geradores infinitesimais das ações adjunta e coadjunta}

Lembre que essas ações correspondem a $G\curvearrowright \mathfrak{g}$ e $G\curvearrowright \mathfrak{g}^*$. Daí,
\begin{itemize}
\item $G\overset{\operatorname{Ad}}{\curvearrowright }\mathfrak{g}$ com
	\begin{align*}
		\operatorname{Ad}: G &\longrightarrow \operatorname{GL}(\mathfrak{g}) \\
		g &\longmapsto (\operatorname{Ad}_g:\mathfrak{g} \to \mathfrak{g} )
	\end{align*}
	"diferenciando" obtemos
	\begin{align*}
		\operatorname{ad}: \mathfrak{g} &\longrightarrow \mathfrak{gl}(\mathfrak{g}) \\
		u &\longmapsto (\operatorname{ad}_u:\mathfrak{g} \to \mathfrak{g} )
	\end{align*}
	
\end{itemize}

\begin{lemma}
	 \begin{align*}
		u_\mathfrak{g} =\operatorname{ad}_u :&\;\mathfrak{g}\longrightarrow \mathfrak{g} \\
		u &\longmapsto u_\mathfrak{g} =[u,\cdot ]
	\end{align*}
\end{lemma}

\begin{proof}[Demostração]
	É só por definição e a regra da cadeia:
	\begin{align*}
		u_{\mathfrak{g}}(v)&=\frac{d}{dt}\Big|_{t=0}\operatorname{Ad}_{\operatorname{exp}(tu)}(v)\\
				   &=\frac{d}{dt}\Big|_{t=0}d_e(R_{\operatorname{exp}(-tu)}\circ  L_{\operatorname{exp}(tu)})(v)\\
		&=\frac{d}{dt}\Big|_{t=0}dR_{\operatorname{exp}(-tu)}(v^L(\operatorname{exp}(tu))\\
		&=[u^L,v^L]|_{e}\\
		&=[u,v]
	\end{align*}
\end{proof}
\end{example}

\subsubsection{Dualização}
Agora considere $G\overset{\operatorname{Ad}}{\curvearrowright }\mathfrak{g}^*$,
\[\left<u_{\mathfrak{g}}^*,v\right> =-\mu\left( [u,v] \right) \]
com $\mu\in\mathfrak{g}^*$, $u_{\mathfrak{g}^*}\in\mathfrak{g}^*=T_\mu\mathfrak{g}^*$.

\begin{question}
	$\mathfrak{g} \to \mathfrak{X}(M)$
\end{question}

\begin{prop}\leavevmode 
	\begin{enumerate}
		\item $M_1,M_2$ $G$-variedades, $ \varphi:M_1\to M_2$ $G$-equivatiante, então $u_{M_1}\overset{\varphi}{\sim}u_{M_2}$, ie; $d\varphi (u_{M_1}(x)=u_{M_2}(\varphi(x))$.

		\item $M$ é $G$-variedad, então $\mathfrak{g} \to \mathfrak{X}(M)$ é (anti!) homomorfismo de álgebra de Lie.

		\item $M$  é $G$-variedade, $(\psi_g)_*(u_M)=(\operatorname{Ad}_g(u))$.
	\end{enumerate}
\end{prop}

\begin{proof}[Demostração]\leavevmode 
	\begin{enumerate}
	\item \[\varphi((\operatorname{exp}(tu) \cdot x)=\operatorname{exp}(tu),\qquad \varphi(x)\]
			\[\frac{d}{dt}\Big|_{t=0}\implies d\varphi (u_{M_1}(x))=u_{M_1}(\varphi(x))\]

		\item Note: resultado vale para \[\phi_g=R_{g ^{-1}}\rightsquigarrow \begin{aligned}
			\mathfrak{g} &\longrightarrow \mathfrak{X}(G) \\
			u &\longmapsto -u^L
		\end{aligned}\]
		Seja $G\curvearrowright \overline{M}=G\times M$ uma ação, \[g(a,z)=(R^{-1}_{g}(a), \qquad u_{\overline{M}}=(-u^L,0)\]

		Note: $[v_{\overline{M}},u_{\overline{M}}=-\left( [u,v]_{\overline{M}} \right) $. Considere
		\begin{align*}
			F: G\curvearrowright \overline{M}=G\times M &\longrightarrow M\curvearrowleft G  \\
			(a,x) &\longmapsto \bar{a}^1x
		\end{align*}
		é $G$-equivariante Daí,
		\[u_{\overline{M}} \overset{\sim}{F}u_M\implies \begin{aligned}
			F_*([u_{\overline{M}},v_{\overline{M}}])&=[u_{M},v_M]\\
			F_*([u,v]_{\overline{M}}&=[u,v]
		\end{aligned}\]
		
		\item A fórmula que eu quero vale em $ \overline{M}$, $F$ é $G$-equivariante, $\implies $ vale em $M$.
	\end{enumerate}
\end{proof}

\paragraph{Próxima aula} 
Estamos usando um funtor
\[\begin{tikzcd}
\{G\text{-variedades} \}\arrow[r]&\{\mathfrak{g} \text{-variedades} \} \arrow[l,bend right, dashed]	
\end{tikzcd}\]
e queremos estudar em que casos podemos voltar.

\section{Aula 14}

$G\overset{\psi}{\curvearrowright }$. Um elemento da algebra de Lie, pega o vetor na identidade, vai para o grupo, daí gera um fluxo na variedade, diferencia o fluxo e já tem um campo vetorial.

\paragraph{Propriedades} 
\begin{itemize}
\item Anti-homomorpfismo de álgebras de Lie.
\item Se $\phi$ é uma aplicação equivariante (preserva ação), os campos vetoriais são relacionados.
\end{itemize}

\subsection{Ações  infinitesimais ($\mathfrak{g}$-ações)}

$\mathfrak{g}$ álgebra de Lie ($\dim <\infty$).

\begin{defn}
	Uma \textit{\textbf{ação (à esquerda)}} de $\mathfrak{g}$ em $M$ é anti-homomorfismo de álgebras de Lie
	\begin{align*}
		\mathfrak{g} &\longrightarrow \mathfrak{X}(M) \\
		u &\longmapsto u_M
	\end{align*}	
\end{defn}

\paragraph{Notação} $\mathfrak{g} \curvearrowright M$, $M$ é $\mathfrak{g}$-variedade.

Para essas variedades exista uma noção de morfismo: suponha que $\mathfrak{g} \curvearrowright M\overset{\phi}{\longrightarrow}M\curvearrowleft \mathfrak{g}$, então $\phi$ é  \textit{\textbf{$\mathfrak{g}$-equivariante}} se
\[u_{M_1}\overset{\phi}{\sim}u_{M_2}\qquad \forall u\in\mathfrak{g}.\]

\paragraph{Em resumo}
\[\begin{tikzcd}[row sep =tiny]
G\text{-variedades} \arrow[rrr,"\text{functor de diferenciação}" ]&&&\mathfrak{g}\text{-variedades, }\mathfrak{g} =\operatorname{Lie}(G) \\ 
g\curvearrowright M\arrow[rrr,maps to]&&&\begin{aligned}
	\mathfrak{g} &\longrightarrow \mathfrak{X}(M) \\
	u &\longmapsto u_M
\end{aligned}\\
\mathfrak{g} \curvearrowright M\overset{\phi}{\longrightarrow}M\curvearrowleft \mathfrak{g}\arrow[rrr]&&&\mathfrak{g} \curvearrowright M\overset{\phi}{\longrightarrow}M\curvearrowleft \mathfrak{g}\arrow[lll,bend left,"\text{integração?}" ]
\end{tikzcd}\]
Nem sempre pode voltar. Precisa que o campo vetorial seja \textbf{completo}, ie. precisamos de uma $\mathbb{R}$-ação infinitesimal $\mathbb{R}\to \mathfrak{X}(M)$. Isso da uma $R$-ação que é um fluxo. Mas, nem tudo campo vetorial é completo.

\begin{thm}[Lie-Richard Palais]\leavevmode
	Seja uma ação 
	\begin{align*}
		\mathfrak{g} &\longrightarrow \mathfrak{X}(M) \\
		u &\longmapsto u_M
	\end{align*}
	seja $G$ um grupo de Lie simplesmente conexo tal que $\mathfrak{g} =\operatorname{Lie}(G)$. Então $\mathfrak{g}$-ação se integra a uma ação $\iff$ cada $u_M$  é completo $\forall u\in\mathfrak{g}$.
\end{thm}

\begin{remark}
	Tem a ver com pensar em uma $S^1$ ação em lugar de uma $\mathbb{R}$-ação. Por isso pedimos simplesmente conexo.
\end{remark}




\paragraph{Sobre sinais/convenção} 

\[\begin{aligned}
	G &\overset{\phi}{\longrightarrow}\operatorname{Diff}(M)  \\
	\psi_g\psi_h &=\phi_{gh} 
\end{aligned}\qquad \qquad \begin{aligned}
g \longrightarrow \mathfrak{X}(M)
\end{aligned}\]
Mas esse daí e um \textit{anti}homomorfismo de álgebra de Lie.

Também temos outros anti-homomorfismos de grupos (é muito importante que aparece uma sinal):
\begin{align*}
	\operatorname{Diff}(M)&\cong \operatorname{Aut}(\mathcal{C}^\infty(M))\\
	\phi :M\to M&\rightsquigarrow \phi^*:\mathcal{C}^\infty(M)\to \mathcal{C}^\infty(M)\\
	\phi \circ \psi&\rightsquigarrow \phi^*\circ \phi^*\\ \\
	\operatorname{Lie}(\operatorname{Diff}(M)) \cong \mathfrak{X}(M) &\overset{\cong }{\longrightarrow}\operatorname{Lie}(\operatorname{Aut}(\mathcal{C}^\infty(M))) \cong \operatorname{Der}(\mathcal{C}^\infty(M))
\end{align*}
Porque de fato a gente já tinha definido o colchete de Lie como uma derivação. Aí o que acontece é que quando voltamos para o lado esquerdo temos que botar uma sinal.

\subsection{Mais sobre ações}

Suponha $G\overset{\psi}{\curvearrowright }M$ e seja $ x\in M$.

\begin{defn}
	\textit{\textbf{Órbita de $x$}}: $\mathcal{O}_x=\{\psi_g(x)=gx:g\in G\} \subseteq M$.
\end{defn}
		A ação da lugar a uma relação de equivalência em $M$ : $x\sim y\iff \exists g\in G$ tal que $y=g\cdot x$. Nesse caso as órbitas são as classes de equivalência (partição de $M$).

\begin{remark}
	Se $M=\mathcal{O}_x$, a ação é dita \textit{\textbf{transitiva}} e  $M$ é \textit{\textbf{espaço $G$-homogéneo}}.
\end{remark}

\begin{example}\leavevmode
	\begin{itemize}
	\item $G\overset{L,R}{G}$, $L_g$ e  $R_g$ são transitivas.
	\item $G\overset{I}{\curvearrowright}G$, $I_g(a)=gag^{-1}$ não é transitiva ($\{e\}$ é órbita).
	\item $\operatorname{SO}(3) \curvearrowright \mathbb{R}^3$ cujas órbitas são esferas e a origem.
	\end{itemize}
\end{example}

\begin{defn} \textit{\textbf{Grupo de isotropía ou estabilizador}}: $G_x=\{g\in G:\psi_g=x\} \subseteq G$.
\end{defn}

\begin{remark}
$x\sim y\implies G_x\cong G_y$.
\end{remark}

\begin{example}
	$\operatorname{SO}(3) \curvearrowright \mathbb{R}^{3}$. Para $x=0$,  $G_x=\operatorname{SO}(3)$, mais para $x \neq 0$ temos $G_x\cong S^1$, as rotações com eixo a linha da origem a $x$.
\end{example}

Agora pegue uma ação infinitesimal $\mathfrak{g} \curvearrowright M$, ie. um mapa
\begin{align*}
	\mathfrak{g} &\longrightarrow \mathfrak{X}(M) \\
	u &\longmapsto u_M
\end{align*}

\begin{defn}\leavevmode
	\begin{enumerate}
		\item \textit{\textbf{Distribuição orbital}}: $D\subseteq TM$ dado por
			\[x\longmapsto T_xM\supseteq D_x=\{u_M(x):u\in\mathfrak{g}\}\]
	
		\item \textit{\textbf{Álebra de isotropía}}: $\mathfrak{g}_{x}\subseteq \mathfrak{g}$ dada por
			\[\mathfrak{g}_x=\{u\in\mathfrak{g} :u_M(x)=0\}\]
	\end{enumerate}
\end{defn}

	Será que tudo encaixa? O estabilizador é um grupo de Lie cuja álgebra de Lie a álgebra de isotropía… a distribuição orbital é o fibrado tangente das órbitas?

\begin{prop}[órbitas, isotropias, … (Warner)]\leavevmode 
	\begin{itemize}
	\item $G_x\subseteq G$ é um subgrupo de Lie e $\operatorname{Lie}(G_x) =\mathfrak{g}_x$.
	\item O conjunto das classes laterias $G/G_x$  é uma variedade $C^\infty$ tal que $G\to G/G_x$ é submersão.

	\item Temos uma aplicação natural 
		 \begin{align*}
			G &\longrightarrow \mathcal{O}_x \\
			g &\longmapsto g\cdot x
		\end{align*}
		E de fato
		\[\begin{tikzcd}[row sep=tiny]
			G\arrow[rr,"\text{sobre} "]&&\mathcal{O}_x\\
			g\arrow[rr,maps to]\arrow[dr]&& g\cdot x\\
			& G/G_x\arrow[ur,"\text{bijeção} "]
		\end{tikzcd}\]
		Assim, $\mathcal{O}_x$ tem estrutura $C^\infty$ tal que $\mathcal{O}_x \overset{\operatorname{dif}}{\cong } G/G_x$.

		E mais. $\mathcal{O}_x\hookrightarrow M$ é uma imersão.
		
	\end{itemize}
\end{prop}

\begin{remark}
	Se $M$ é $G$-variedade homogênea, então $M\cong G/G_x$ para qualquer  $x\in M$.
\end{remark}

\begin{example}\leavevmode
	\begin{itemize}
		\item $S^{n-1}$. Qual é grupo de isotropia de $e_n$?
			\[\begin{pmatrix} &&& 0\\
			& \operatorname{O}(n-1)&&0\\
			&&&0\\
			0&0&0&1\end{pmatrix} \]
		de fato,
		\[S ^{n-1}\cong \dfrac{\operatorname{O}(n)}{\operatorname{O}(n-1)}=\dfrac{\operatorname{SO}(n)}{\operatorname{SO}(n-1)}\]
		já que os dois agem transitivamente na esfera---lembre que a variedade não é mais que o quociente do grupo-estabilizador quando tem uma ação homogênea.

\item \[\mathbb{R}P^{n-1}=\dfrac{\operatorname{SO}(n)}{\operatorname{O}(n-1)}\]
	Pode fazer uma leve meditação ao respeito. Agora tem que fixar uma linha, não só um ponto. Tem que fixar mais pontos, tem menos transformações. Ele colocou uma matriz assim:
				\[\begin{pmatrix} &&& 0\\
			& \operatorname{O}(n-1)&&0\\
			&&&0\\
			0&0&0&\pm 1\end{pmatrix} \]
Porque pode refleitir no ponto antípoda e tudo bem.
	\item \[\mathbb{C}P^{n_01}=\dfrac{\operatorname{SU}(n)}{\operatorname{U}(n-1)}\]
	\end{itemize}
\end{example}

\paragraph{Fato geral}  $G$ grupo de Lie $H\subseteq G$ subgrupo fechado, então $G/H$  é $C^\infty$-variedade e $G\to G/H$ é submersão.

Lembre: uma ação livre, propria… tem o Slice theorem, onde tem que pegar uma vizinhança das órbitas.

\begin{example}[Curiosidade]
	$V$ espaço vetorial, $\operatorname{Fr}_p(V)=\{(v_1,\ldots,v_p):\text{l.i.} \}$. Tem uma ação $\operatorname{GL}(n)\curvearrowright \operatorname{Fr}_p(V)$. Pode calcular o grupo de isotropía $H$, e daí dar uma estrutura diferenciável a $\operatorname{Fr}_p(V)$. Essa variedade se chama \textit{\textbf{variedade de Stiefel}}.
\end{example}

\subsection{De volta à geometria simplética}

Queremos finalmente estudar órbitas coadjuntas e quocientes simpléticos.

Pegue um grupo de Lie $G$ e lembre a ação coadjunta $G\overset{\psi=\operatorname{Ad}^*}{ \curvearrowright }\mathfrak{g}^*$. É assim:
\begin{align*}
	(\operatorname{Ad}^*)_g: \mathfrak{g}^* &\longrightarrow \mathfrak{g}^* \\
	\left<\operatorname{Ad}^*)_g(\xi ),u\right>  &=\left<\xi,\operatorname{Ad}_{g^{-1}}(u)\right> 
\end{align*}
onde $u\in\mathfrak{g}$ e $\xi \in\mathfrak{g}^*$. De fato,
\begin{align*}
	\operatorname{Ad}_g=d_e I_g:\mathfrak{g} &\longrightarrow \mathfrak{g} \\
	I_g(a) &=gag^{-1}
\end{align*}

\begin{remark}
	$I_g:G\to G$ é um homomorfismo de grupos, $\implies \operatorname{Ad}_g:\mathfrak{g} \to \mathfrak{g}$ é homomorfismo de álgebra de Lie, ie. $\operatorname{Ad}_g[u,v]=[\operatorname{Ad}_gu,\operatorname{Ad}_gv]$.

	Também vimos que
	\begin{align*}
		: \mathfrak{g} &\longrightarrow \mathfrak{X}(\mathfrak{g}^*) \\
		u &\longmapsto u_{\mathfrak{g}^*}
	\end{align*}
	\[u_{\mathfrak{g}^*}=-\xi([u,\cdot ])\]
	onde $T_\xi \mathfrak{g}^*=\mathfrak{g}^*$
\end{remark}

Agora. Pegue $\eta\in\mathfrak{g}^*$ fixo. então
\[\mathcal{O}=\mathcal{O}_\eta=\{\operatorname{Ad}^*(\eta):g\in G\} \subseteq \mathfrak{g}^*\]
\[T_\xi \mathcal{O}=\{u_{g^*}(\xi ):u\in\mathfrak{g}\}\]

\begin{defn}
	$\xi \in\mathcal{O}$, $\omega_\xi \in\Lambda^{2}(T_\xi^*\mathcal{O})$,
	\[\omega_\xi(u_{\mathfrak{g}^*} (\xi ),v_{\mathfrak{g}^*}(\xi ):=\left<\xi,[u,v]\right> \]
	Onde tem que mostrar que isso está bem definido. Para isso poderia usar que $\left<\xi,[u,v]\right> =-u_{\mathfrak{g}^*}(\xi )(x)$…
\end{defn}

\begin{thm}[KKS=Kirillov-Kostant-Sourian]
	$\omega\in\Omega^{2}(\mathcal{O})$ é simplética.
\end{thm}

\begin{proof} Seja $\xi \in\mathcal{O}$
	\begin{enumerate}[label=\textbf{Paso \arabic*}]
		\item \textbf{$\omega_\xi$ está bem definido.} $u,u'\in\mathfrak{g}$ tal que $u_{\mathfrak{g}^*}=u'_{\mathfrak{g}^*}(\xi )$, mais isso é equivalente a que $\xi([u,\cdot ])=\xi([u',\cdot ])$.

	\item \textbf{$\omega_\xi$ não degenerada}:
			\begin{align*}
				\omega_\xi(u_{\mathfrak{g}^*}(\xi ),v_{\mathfrak{g}^*}(\xi )&=0\\
				\implies =\xi([u,v])&=u_{\mathfrak{g}^*}(\xi )(x)=0\\
				\iff u_{\mathfrak{g}^*}(\xi)&=0.
			\end{align*}

\item \textbf{$\xi \longrightarrow \omega_\xi$ é suave}. Para isso vamos mostrar que $\omega$ é $G$-invariante, ie. que $(\operatorname{Ad}^*)_g\xi =g\xi$.

	\begin{exercise}
		Numa variedade $G $-homogénea, uma forma $G$-invariante é suave.
	\end{exercise}
\begin{align*}
	(\mathfrak{g}^*\omega )_\xi(u_{\mathfrak{g}^*}(\xi,v_{\mathfrak{g}^*}(\xi))&\overset{\operatorname{def}}{=}\omega_{g\xi}(g_*(u_{\mathfrak{g}^*}(\xi ),g_*(v_{\mathfrak{g}^*}(\xi )))\\
	&=\omega_{g\xi}\left((\operatorname{Ad}_g(u))_{\mathfrak{g}^*}(\xi ),(\operatorname{Ad}_g(v))_{\mathfrak{g}^*}(x)\right)\\
	&\overset{\operatorname{def}\omega}{=}\left<(\operatorname{Ad}^*)_g\xi,[\operatorname{Ad}_g(u),\operatorname{Ad}_g(v)]\right> \\
	&=\omega_\xi(u_{\mathfrak{g}^*}(\xi ),v_{\mathfrak{g}^*}(\xi ))\\
	\implies g^*\omega &=\omega
\end{align*}

\item \textbf{$\omega$ é fechada}. Vamos usar invariânça de novo. A invariância significa que
	\[0=\mathcal{L}_{u_{\mathfrak{g}^*}}\omega=i_{u_{\mathfrak{g}^*}}d\omega +di_{u_{\mathfrak{g}^*}}\omega\]
	\[\implies i_{u_{\mathfrak{g}^*}}d\omega =-d(i_{u_{\mathfrak{g}^*}}\omega )\]
	\begin{claim}
		$i_{u_{\mathfrak{g}^*}}\omega =du$. Já como é o dual do dual, podemos ver a $u$ como uma função $u\in\mathfrak{g} =\mathcal{C}^\infty_{\operatorname{lin}}(\mathfrak{g}^*)\subseteq \mathcal{C}^\infty(\mathfrak{g})$.
	\end{claim}
	Como $u$ é linear, $d|_{\xi}(\eta)=u(\eta)$. Daí,
	\begin{align*}
		du|_{\xi}(v_{\mathfrak{g}^*}(\xi )&=u(v_{\mathfrak{g}^*}(\xi ))=\left<\xi,[u,v]\right> =\omega_{\xi}(u_{\mathfrak{g}^*}(\xi ),v_{\mathfrak{g}^*}(\xi )).
	\end{align*}
	\end{enumerate}
\end{proof}

\section{Aula 15}

Na aula pasada:
\begin{itemize}
\item Ação coadjunta de uma grupo de Lie na sua álgebra de Lie:
	\[G\overset{\operatorname{Ad}^*}{\curvearrowright}\mathfrak{g}^*\qquad \qquad \operatorname{Ad}_{g}=d I_g,\qquad I_g(a)=gag^{-1}\]
\item Numa órbita $\mathcal{O}\hookrightarrow \mathfrak{g}^*$,
	\[\omega\left( u_{\mathfrak{g}^*}(\xi),v_{\mathfrak{g}^*}(\xi) \right) :=\left<\xi,[u,v]\right> \]
	para $u,g\in\mathfrak{g},\xi\in\mathcal{O}\subseteq \mathfrak{g}^*$, é uma estrutura simplética em $\mathcal{O}$.
\end{itemize}
Hoje: exemplos concretos

\subsection{Exemplos concretos (da ação coadjunta)}
\begin{remark}
	Suponha que $k=\left<\cdot,\cdot\right> $ é uma forma bilinear simplética \textit{não degenerada} e \textit{$\operatorname{Ad}$-invariante} em  $\mathfrak{g}$, ie.
	\[\left<\operatorname{Ad}_g(u),\operatorname{Ad}_g(v)\right> \]
	Então podemos identificar órbitas adjuntas e coadjuntas assim:
	\begin{align*}
		k^\flat: \operatorname{Ad}\curvearrowright \mathfrak{g} &\xrightarrow{\cong }\mathfrak{g}^*\curvearrowleft \operatorname{Ad}^*  \\
		u &\longmapsto \left<u,\cdot\right> 
	\end{align*}
	já que é \textit{equivariante}.
\end{remark}

\begin{thing3}{Exemplo concreto da órbita e ação coadjunta}\leavevmode
	\begin{align*}
		\operatorname{SO}(3)&=\{A\in \mathcal{M}_{3\times 3}(\mathbb{R}):AA^{\mathbf{T}}=\operatorname{id},\det A=1\} \\
		\mathfrak{so}(3) &=\{A\in\mathcal{M}_{3\times 3}(\mathbb{R}):A=-A^{\mathbf{T}}\}\\
	\end{align*}
	Lembre que a ação adjunta de $\operatorname{SO}(3)$ em $\mathfrak{so}(3)$ é por conjugação. Agora, podemos identificar
	\begin{align*}
		 \mathfrak{so}(3) &\longrightarrow \mathbb{R}^{3} \\
		\begin{pmatrix} 0&-z&y\\z&0&-x\\-y&x&0 \end{pmatrix}  &\longmapsto \begin{pmatrix} x\\y\\z \end{pmatrix} =u\\
		[\cdot,\cdot ]& \longmapsto (u,v)\mapsto u\times v
	\end{align*}
Aqui a ação é por multiplicação de matrizes! I.e., $\operatorname{SO}(3) \curvearrowright \mathfrak{so}(3) \cong \mathbb{R}^{3}$ é dada por
\begin{align*}
	\operatorname{Ad}_A(v)&=Av,\qquad v\in\mathbb{R}^{3}\\
	\operatorname{ad}_u(v)&=u\times v
\end{align*}
e além disso, se o produto interno canónico de $\mathbb{R}^{3}$ é $\left<\cdot,\cdot\right> $, temos que essa ação é $\operatorname{Ad}_{\mathfrak{so}(3)}$-invariante. Isso nos permete 
\[\begin{tikzcd}
	\mathfrak{so}(3)\arrow[r]\arrow[d,"\cong "]&\mathfrak{so}(3)\arrow[d]\\
	\mathbb{R}^{3}\arrow[r,"\left<\cdot\text{,} \cdot\right>" ]&(\mathbb{R}^{3})^*
\end{tikzcd}\]
	Para entender as orbitas coadjuntas de $\operatorname{SO}(3)$ em $\mathfrak{so}(3)$ basta analizar o caso em $\mathbb{R}^{3}$.

	{\color{5}\bfseries \'Orbitas (co-)adjuntas em $\mathbb{R}^{3}$ ($\cong \mathfrak{so}(3)$).}\hspace{.5em} De fato, isso vai ser um exercício. Vai concluir que a estrutura simplética nas órbitas dessa ação coincide com a estrutura simplética na esfera.

	\begin{exercise}\leavevmode
		Mostre que a forma de área na esfera é um múltiplo constante da forma  KKS. O múltiplo é $1/r$ onde  $r$ é o raio da esfera.
	\end{exercise}
\end{thing3}

\begin{thing6}{Outro exemplo concreto}\leavevmode
	\begin{align*}
		G&=\operatorname{U}(n) =\{A\in\operatorname{GL}(n,\mathbb{C}):A A^*=\operatorname{id}\} \\
		\mathfrak{g} &=\mathfrak{u}(n)=\{u\in\mathcal{M}_{n\times n}(\mathbb{C}):u=-u^*\}
	\end{align*}
	Aqui a ação adjunta $\operatorname{U}(n) \curvearrowright \mathfrak{u}(n)$ é por conjugação. Agora vamos estudar a ação coadjunta.

	Temos um prouto interno $\operatorname{Ad}$-invariante em $\mathfrak{u}(n)$:
	\[(u,v)\longmapsto \operatorname{tr}(u^kv)=-\operatorname{tr}(uv).\]
	Portanto temos uma identificação 
	\[\operatorname{U}(n) \curvearrowright \operatorname{U}(n) \xrightarrow{\cong }\mathfrak{u}(n)^*\curvearrowleft \operatorname{U}(n).\]
	Ou seja, a forma de estudar isto é olhando para a ação de $\operatorname{U}(n)$.

	Note ainda que $\mathfrak{u}(n) =i\mathcal{H}$ onde $\mathcal{H}=\{\xi \in\mathcal{M}_{n}(\mathbb{C}):\xi=\xi^*\}$. Com isso,
	\[\operatorname{U}(n) \overset{\operatorname{Ad}^*}{\curvearrowright}\mathfrak{u}(n)^*\cong \mathfrak{u}(n) \cong \mathcal{H}\overset{\operatorname{Ad}^*}{\curvearrowleft}\operatorname{U}(n)\]
	e então 
	\[\operatorname{Ad}_A^*(\xi)=A\xi^{-1}A.\]
	Portanto,
	\begin{align*}\xi,\xi'\in\mathcal{H}\text{ estão na mesma órbita coadjunta }\mathcal{O}\\
		\overset{\text{teo. espectral} }{\iff}  \operatorname{sp ec}(\xi)=\operatorname{s pec}(\xi')
		\end{align*}
		Ado $\Lambda=(\lambda_1,\ldots,\lambda_n)$, $\lambda_i\in\mathbb{R}$, $\lambda_1\leq \ldots \leq \lambda_n$. Daí as órbitas são
		\[\mathcal{O}_\Lambda=\{\xi \in\mathcal{H}:\operatorname{s pec}(\xi)=\Lambda\}\]
		e se chaman de \textit{\textbf{isoespetrais}}.

	{\color{5}Qual é a topologia dessa órbita?}
\begin{thing3}{Alguns casos}\leavevmode
	\begin{itemize}
	\item $\lambda_1<\lambda_2=\ldots =\lambda_n$ cada $\xi\in\mathcal{O}_\Lambda$ é determinada pelo autoespaço de $\lambda_1$, ou seja, una linha em $\mathbb{C}^{n}$. Defina $\mu=\lambda_2=\ldots =\lambda_n$. (Como o teorema espectral nos da uma descomposição ortogonal com os autoespaços de uma transformação hermitiana, a ação de uma transformação $\xi \in\mathcal{H}$ fica determinada por essa linha.) Daí,
		\[\mathcal{O}_\Lambda\cong \mathbb{C}P^{n-1}\]

	\item $\lambda_1=\ldots =\lambda_k<\lambda_{k+1}=\ldots =\lambda_n$. Obtemos
		\[\mathcal{O}_\Lambda=\operatorname{Gr}(k,n)\]
	\item $\lambda_1<\lambda_2<\ldots <\lambda_n$. Aqui $\xi \in\mathcal{O}_\Lambda$ está determinada pelos autoespaços $L_1\sim\lambda_1,\ldots,L_n\sim\lambda_n$. Equivalentemente,
		\begin{align*}
			E_1&=L_1\\
			E_2&=L_1\oplus L_2\qquad \qquad \text{ aqui só $L_2$ e nova info.}\\
			\vdots \\
			E_j&=L_1\oplus \ldots \oplus L_j\\
			\vdots \\
			E_n&=\mathbb{C}^n
		\end{align*}
		Então temos
		 \[L_1=E_1\subseteq\ldots \subseteq E_n=\mathbb{C}^{n}\]
		isso se chama de \textit{\textbf{variedade bandeira (completa)}}.
	\end{itemize}
\end{thing3}
\end{thing6}

\begin{thing1}{Projeto?}\leavevmode
	Relação entre geometria simplética e álgebra linear. Tem um teorema básico sobre matrizes hermitianas:
	\begin{thm}[Schur-Horn]\leavevmode
		$\mathcal{O}_\Lambda$ uma variedade isoespetral (todos os caras hermitianos cujo espetro e esse $\Lambda$), $\{\xi \in\mathcal{H}:\operatorname{s pec}(\xi)=\Lambda\}$.
		Considere o mapa diagonal
		\begin{align*}
			\operatorname{diag}: \mathcal{O}_\Lambda &\longrightarrow \mathbb{R}^{n}
		\end{align*}
		(É só pegar a diagonal da matriz e colocá-la num vetor.) Então, a imagem de $\operatorname{diag}$ é um poítopo convexo com vertices dados pelas $n!$ permutações  de $\Lambda=(\lambda_1,\ldots,\lambda_n)$.
	\end{thm}
	\begin{remark}
		Esse teorema pode ser generalizado (Atiyah-?) trocando $\mathcal{O}_\Lambda$ por uma variedade simplética com a ação do toro maximal.
	\end{remark}
	\begin{thing3}{Horn's problem}\leavevmode
		Descrever o espectro da soma de dois matrizes hermiticanas com espectros conhecidos. Tem um artigo expositório (survey) de A. Knutson, \textit{AThe symplectic and algebraic geometry of Horn's problem}. Também tem um artigo de pesquisa famoso com Terry Tau, e outro dele com Terry Tau e Chris Woodsword (esses estão citados no survey).
	\end{thing3}
\end{thing1}

\subsection{Ponto de vista "Poisson"}

Lembre:
 \begin{thing7}{Estrutura de Poisson}\leavevmode
 	\begin{align*}
 		\{\cdot,\cdot\}: \mathcal{C}^\infty(M)\times \mathcal{C}^\infty(M) &\longrightarrow \mathcal{C}^\infty(M)
 	\end{align*}
	forma $ \mathbb{R}$-bilinear antisimétrica ($\{f,g\} =-\{g,f\}$) Jacobi e Leibniz ($\{f,gh\} =\{f,g\} h+\{f,h\} g$.
 \end{thing7}

\begin{thing6}{Exemplos}\leavevmode
	\begin{itemize}
	\item $\{\cdot,\cdot\} =0$.
	\item $(M,\omega)$ simplética, $\{f,g\} =\omega(X_g,X_f)$.
	\item (Exemplo nem simplético nem trivial.) $\mathbb{R}^{3}=\{\xi=(x,y,z)\}$,
		\[\{f,g\}(\xi)=\left<\xi,\nabla f(\xi)\times \nabla g(\xi)\right> \]
		onde $\nabla$ denota gradiente.

		Note que pela condição Leibniz do colchete de Poisson, $\{f,\cdot\}$ é uma derivação, que sabemos que é equivalente a dizer que é um campo vetorial. Assim, para toda função $f\in\mathcal{C}^\infty(M)$ existe um campo vetorial Hamiltoniano. De fato, isso é consistente com o caso do colchete de Poisson associado a uma variedade simplética.

		Pegue o Hamiltniano
		\[H(x,y,z)=\frac{x^2}{2I_1}+\frac{y^2}{2I_2}+\frac{z^2}{2I_3}\]
		Daí vêm as equações de Euler-?
	\end{itemize}
\end{thing6}

\subsubsection{Descrição tensorial (do colchete de Poisson)}

Considere $\pi \in\Gamma(\Lambda^{2}(TM))$ e o colchete
\[\{f,g\} =\pi(df,dg).\]
Em coordenadas
\[\pi=\sum_{i<k}\pi_{ij}(x)\frac{\partial }{\partial x_i}\wedge \frac{\partial }{\partial x_j}\]

\subsubsection{Distribuição característica}
Os vectores são vectores  hamiltonianos de alguma função :
\[R\subseteq TM,\qquad \qquad R_x=\{X_f(x):f\in\mathcal{C}^\infty(M)\}\]

\begin{thm}\leavevmode
	\begin{itemize}
	\item $R\subseteq TM$ é integrável.
	\item Cada folha tem estrutura simplética natural. Obtemos a \textit{\textbf{folheação simplética}}. Ela determina, caracteriza, a estrutura de Poisson.
	\end{itemize}
\end{thm}

\subsubsection{Uma clase especial de variedades de Poisson}

Uma clase de estruturas de Poisson à que pertence o exemplo que nem é trivial nem é simplético.

$M=V$ espaço vetorial. $\mathcal{C}^\infty_{\operatorname{lin}}(V)\subseteq \mathcal{C}^\infty(M)$. Note que  $\{\cdot,\cdot\}$ é linear. Suponha que $\left\{ \mathcal{C}^\infty_{\operatorname{lin}}(V),\mathcal{C}^\infty_{\operatorname{lin}}(V) \right\} \subseteq \mathcal{C}^\infty_{\operatorname{lin}}(V)$. Note que $\mathcal{C}^\infty_{\operatorname{lin}}(V)\cong V^*$.

Por outro lado, $(\mathfrak{g},[\cdot,\cdot ]\rightsquigarrow \{\cdot,\cdot\}$ é linear em $\mathfrak{g}^*$, e como $\mathfrak{g} =\mathcal{C}^\infty_{\operatorname{lin}}(\mathfrak{g}^*)\subseteq \mathcal{C}^\infty(\mathfrak{g})$, existe um único $\{\cdot,\cdot\}$ tal que
\[\{\cdot,\cdot\}|_{\mathfrak{g}}=[\cdot,\cdot ]\]

Agora vamos ver a fórmula para o colchete de Poisson associado a $[\cdot,\cdot ]$:
\[\{f,g\}(\xi)=\left<\xi,[df(\xi),dg(\xi)]\right> \]
para $f,g\in\mathcal{C}^\infty(\mathfrak{g}^*)$.

\begin{thing4}{Conclusão}\leavevmode
	\[\substack{\text{estruturas de Poisson}  \\ \text{lineares em $V=\mathfrak{g}^*$} }\longleftrightarrow \substack{\text{álgebra de Lie}  \\ \text{em $V^*=\mathfrak{g}$} }\]
\end{thing4}

Pode ver que
\[\begin{tikzcd}
	&(\mathfrak{g},[\cdot,\cdot ]\text{ álg. de Lie} \arrow[dl]\arrow[dr]\\
	\substack{ G \overset{\operatorname{Ad}^*}{\curvearrowright}\mathfrak{g}^*\\\text{simpléticas}} && \substack{(\mathfrak{g}^* ,\{\cdot,\cdot\})\\\text{folheação simplética}} 
\end{tikzcd}\]
de onde as órbitas coadjuntas correspondem às folhas. O importante aqui é saber que a folheação tem estrutura de Poisson.


\section{Aula 16: ações Hamiltonianas e aplicação momento}

Vimos
\begin{itemize}
\item $G$ grupo de Lie. $\mathcal{O}\hookrightarrow \mathfrak{g}^*$ orbitas coadjuntas possuim uma estrutura simplética canónica. Isso na verdade é um exemplo de um objeto mais geral que a gente vai discutir hoje: ações hamiltonianas.
\end{itemize}
\subsection{Caso mais simples}
Quando $G=\mathbb{R}$. Seja $(M,\omega)$ simplética. Temos:
\[\mathbb{R}\text{-ação =fluxo }\psi \rightsquigarrow X \in\mathfrak{X}(M) \text{campo de vetores}  \]
\begin{defn}\leavevmode
\begin{itemize}
\item Uma $\mathbb{R}$-ação é \textit{\textbf{simplética}} se $\psi_t^*\iff\mathcal{L}_{X}\omega=0$. ($X$ é um campo simplético.)
\item $R$-ação é \textit{\textbf{hamiltoniana}} se  $X$ é hamiltoniano, i.e., existe uma função $H$ tal que $X=X_H$
\end{itemize}
\end{defn}

Agora vamos generalizar isso.

\subsection{Caso geral}

\begin{defn}\leavevmode
	\begin{itemize}
	\item Uma $G$-ação é simplética se $\psi_{g}^*\omega=\omega$ para toda $g \in G$, i.e.
		\[\begin{tikzcd}
			 \psi:G\arrow[r]\arrow[dr]&\operatorname{Dif}(M)\\
			  &\operatorname{S im p}(M,\omega)\arrow[u,hook]
		\end{tikzcd}\]
		
	\item $G$-ação é \textit{\textbf{hamiltoniana}} se existe uma \textit{aplicação momento}.
	\end{itemize}
\end{defn}

Vamos supor a exitência de um mapa

\[\begin{tikzcd}[ column sep=small,row sep=small]
	&\mathcal{C}^\infty(M)\arrow[dd]& f\arrow[dd,maps to]\\
\mathfrak{g} \arrow[ur,"\hat{\mu}"]\arrow[dr]\\
u\arrow[dr, maps to]&\mathfrak{X}(M)&X_f\\
&u_M
\end{tikzcd}\]

\begin{remark}\leavevmode 
	\begin{itemize}
	\item Se $\hat{\mu}$ existe, não é unicamente determinado.
	\item Sempre podemos supor que $\hat{\mu}$ é linear.
	\end{itemize}
	\item É equivalente dar $\hat{\mu}$ como função $\hat{\mu}:\mathfrak{g} \longrightarrow \mathcal{C}^\infty(M)$ ou como uma função
		\begin{align*}
			\mu: M &\longrightarrow \mathfrak{g}^* \\
			\hat{\mu}(u)(x) & =\left<\mu(x),u\right>  
		\end{align*}
		Note que as quantidades nessa igualdade são numeros reais. No lado esquerdo, temos a função $\hat{\mu}(u)\in\mathcal{C}^\infty(M)$ evaluada no ponto $x\in M$. No lado direito, temos o funcional $\mu(x)\in\mathfrak{g}^*$ evaluado no vetor $u\in\mathfrak{g}$.
\end{remark}

\begin{defn}\leavevmode
	Seja $G\overset{\varphi}{\curvearrowright}(M,\omega)$ uma ação simplética. Dizemos que é \textit{\textbf{fracamente hamiltoniana}} se existe  $\mu:M\to \mathfrak{g}^*$ (equivalentemente, $\hat{ \mu}:\mathfrak{g} \longrightarrow \mathcal{C}^\infty(M)$ linear) tal que
	 \begin{equation}\label{eq:a}
	i_{u_M}\omega=d\left<\mu,u\right> \qquad (u_M=X_{\hat{\mu}(u)})	 	
	 \end{equation}
	e a ação é \textit{\textbf{hamiltoniana}} se, além disso, vale
	\begin{equation}\label{eq:b}
		\mu\text{ é equivariante} 
	\end{equation}
	issto é, $\mu:G \curvearrowright M\longrightarrow \mathfrak{g}^*  \overset{\operatorname{Ad}^*}{\curvearrowleft}G$,
	\[\mu\circ \psi_g=(\operatorname{Ad}^*)_g(\mu)\]
	{\color{3}ou seja, a ação de $\mu$ é equivariante com respeito à ação de $G$ em $M$ e à ação coadjunta em $\mathfrak{g}^*$,
	\[\mu(g\cdot x)=\operatorname{Ad}^*_g\mu(x)\]
(lembre que equivariância de uma ação é "linearidade" com respeito as ações).}
\end{defn}

No caso abeliano, que é o caso dos fluxos, essa segunda condição de equivariança é trivialmente satisfeita.

\begin{thing5}{Terminologia}\leavevmode
	\begin{itemize}
	\item $(M,\omega,\mu)$ é um  \textit{\textbf{$G$-espaco hamiltoniano}}.
	\item $\mu$ é \textit{\textbf{aplicação momento}}.
	\item $G$ abeliano: $\mu$ é $G$-equivariante $\iff$ $\mu$ é $G$-invariante.
	\end{itemize}
\end{thing5}

\begin{example}\leavevmode
No caso abeliano, isso vira simplesmente escolher uma função que seja a hamiltoniana do campo vetorial associado à $ \mathbb{R}$-ação. Issto é, a apliação momento para uma $\mathbb{R}$-ação com campo gerador $X$ é uma função $H\in\mathcal{C}^\infty(M)$ tal que $X=X_H$. Note que  $H$ é constante ão longo do fluxo de $X_H$.
	\end{example}

	(Agora?) veremos resultados de existência/obstrução para aplicações momento adiante. (Adiante.)

\begin{remark}\leavevmode
	Em termos do \textit{\textbf{comomento}}  $\hat{\mu}:\mathfrak{g}\longrightarrow \mathcal{C}^\infty(M)$, a equivariância de $\mu$ é equivalente a que
	\begin{align*}
		\iff \psi^* _g(\hat{\mu}(u))&=\hat{\mu}(\operatorname{Ad}_{g^{-1}}(u))\qquad \forall u\in\mathfrak{g}\\
		\implies \mathcal{L}_{v_M}(\hat{\mu}(u))&=\hat{\mu}\left( [u,{\color{3}-}v] \right) 
	\end{align*}
	pegando $g=\operatorname{exp}(tv)$, $\frac{d}{dt}\Big|_{t=0}$. {\color{3}(Literalmente é isso! Do lado esquerdo aparece a derivada de Lie por definição, e do lado direito aparece o colchete por definição de $\operatorname{ad}$)} Lembre que $\{\hat{\mu}(v),\hat{\mu}(v)\} =\mathcal{L}_{v_M}(\hat{\mu}(u))$.
	i.e. $G$-equivariância implica que $\hat{\mu}$ é um \textit{antihomomorfismo} de álgebras de Lie. A conversa é verdade supondo que $G$ é conexo.
\end{remark}

\begin{remark}\leavevmode
	A primeira condição da definição diz que podmoe achar a função que é o hamiltoniano do campo vetorial associado ão elemento em $\mathfrak{g}$. A segunda condição diz que esse mapa $\mathfrak{g} \longrightarrow \mathcal{C}^\infty(M)$ é um antihomomorfismo de álgebras de Lie (ver diagrama)
\end{remark}

\subsection{Princípio de Noether}

$(M,\omega,\mu)$ $G$-espaço hamiltoniano, $H\in\mathcal{C}^\infty(M)$ hamiltoniano. Então,
\[H \text{ é $G$-invariante}\iff \mu\text{ é preservada pelo fluxo hamiltoniano de }H  \]
\begin{proof}\leavevmode
	\begin{align*}
		\mathcal{L}_{X_H}\left<\mu,u\right> &=i_{X_H}d(\left<\mu,u\right> )=i_{X_H}i_{u_M} \omega=-i_{u_M} \underbrace{i_{X_H}\omega}_{dH}=-\mathcal{L}_{u_M}H
	\end{align*}
\end{proof}
(É um desses exemplos onde a definição faz o teorema trivial.)

\subsection{Exemplos de ações Hamiltonianas}

\begin{enumerate}
	\item $M=\mathbb{R}^{2n}=T^*\mathbb{R}^{n}$. $\omega=\sum_{i}dq_i\wedge dp_i$. $G=\mathbb{R}^{n}$. $\mathbb{R}^{n}\curvearrowright \mathbb{R}^{2n}$ por traslação, ie. $g\cdot(p,q)=(q+g,p)$. Isso claramente é uma ação. Mas, como pode achar uma aplicação momento para ela?

		{\color{4}\bfseries Geradores infinitesimais.}\hspace{.5em}$g=tu$,
		\[u_{M}|_{(q,p}=\sum_{i}u_i \frac{\partial }{\partial q_i}\]
		Queremos ver se esse campo é hamiltoniano. De fato,
		\[i_{u_M}\omega=i_{u_M}\left( \sum_{i}dq_i\wedge dp_i \right) =\sum_{i}u_i dp_i=d \left( \left<p,u\right>  \right) \]

	\item  {\color{4}\bfseries Aplicação momento.}\hspace{.5em}
\begin{align*}
	\mu: \mathbb{R}^{2n} &\longrightarrow \mathbb{R}^{n} \\
	(q,p) &\longmapsto p
\end{align*}

\end{enumerate}

\begin{thm}[Noether]\leavevmode
	$H$ invariante por traslação $\iff$ $p$ é preservado pelo fluxo de $H$
\end{thm}

\begin{remark}\leavevmode
	Lembre do ensino medio onde tinhamos energia cinética mas potencial, $K+V$.  O teorema diz que um potencial invariante por traslação iff momento linear é preservado.
\end{remark}

\subsection{De forma mais geral}

Quando temos $(M,\omega)$ e $\omega=-d\theta$ é exata. Com uma ação $G\overset{\psi}{\curvearrowright}(M,\omega)$, temos $\psi_g^*\theta=\theta$, então $\mathcal{L}_{u_{M}}\theta=0$. Nesse caso, a ação é hamiltoniana e a aplicação momento está dada por
\begin{align*}
	M &\longrightarrow \mathfrak{g}^* \\
	\left<\mu,\theta\right>  &=i_{u_{M}}\theta
\end{align*}

\begin{proof}\leavevmode
	\begin{itemize}
	\item 
		 \begin{align*}
			i_{u_M}\omega&\overset{?}{=} d \left<\mu,u\right> =di_{u_M}\theta=\underbrace{\mathcal{L}_{u_{M}}}_{=0}\theta-i_{u_{M}}d\theta=i_{u_{M}}\omega
		\end{align*}
	\item 
		\begin{align*}
			(\operatorname{Ad}^*)_g\mu&\overset{?}{=}\left<\mu\circ \psi_g,u\right> =\psi^* _g(\left<\mu,u\right> )=\psi^* _g(i_{u_M}\theta)=i_{(\psi_{g^{-1}})^*u_M}(\psi_g^*\theta)\\
			&=i_{(\operatorname{Ad}_{g^{-1}}(u))_M}\theta=\left<\mu,\operatorname{Ad}_{g^{-1}}(u)\right> =\left<(\operatorname{Ad}^*)_g\mu,u\right> 
		\end{align*}
	\end{itemize}
\end{proof}

\begin{example}[Levantamiento cotangente]\leavevmode
	$M=T^*Q$, $\omega_{\operatorname{can}}=-d\alpha$.
	\[G\overset{\psi}{\curvearrowright}Q\rightsquigarrow G\overset{\hat{\psi}}{\curvearrowright}T^*Q\]
	\[\hat{\psi}_g=(d\hat{\psi}_{g^{-1}})^*\]
	Y vimos que o levantamiento cotangente preserva a forma tautológica, i.e., $(\hat{\psi}_{g})^*\alpha=\alpha$. Isso implica que temos uma ação hamiltoniana.

	\begin{thing5}{Ação momento}\leavevmode
		\begin{align*}
			\mu:T^*Q  &\longrightarrow \mathfrak{g}^* \\
			\left<\mu,u\right>  &=i_{u_M}\alpha 
		\end{align*}
		
		\textbf{Simplificação:}  $\left<\mu(\xi),u\right> =\xi(u_Q)$.

		Lembre como estava definida a forma tautológica:
		\[\begin{tikzcd}
			\xi \in\arrow[d,maps to]& T^*Q\arrow[d,"\pi"]\\
			x\in&Q
		\end{tikzcd}\]
		onda simplesmente $\alpha_\xi=\pi^*\xi$.
		
	\end{thing5}
\end{example}

\begin{example}[Outros exemplos]\leavevmode
	\begin{itemize}
	\item $\operatorname{SO}(3) \curvearrowright \mathbb{R}^3\overset{\text{levantamento} }{\rightsquigarrow }\operatorname{SO}(3) \curvearrowright \mathbb{R}^{6}\cong T^*\mathbb{R}^{3}\cong \mathbb{R}^{3}\times \mathbb{R}^{3}\times \mathbb{R}^{3}$. Com aplicação momento (dada pelo isomorfismo da álgebra $\mathfrak{so}(3)^*\cong \mathbb{R}^{3}$ que já vimos?)
		\begin{align*}
			\mu: \mathbb{R}^{3}\times \mathbb{R}^{3} &\longrightarrow \mathfrak{so}(3)^*\cong \mathbb{R}^{3} \\
			(q,p) &\longmapsto q\times p
		\end{align*}
		que se chama \textit{\textbf{momento angular}}.
	\end{itemize}

	\item Produtos direitos, restrição a subgrupo (vai vir na lista).

	\item \'Orbitas coadjuntas. $G \curvearrowright(\mathcal{O},\omega_{\operatorname{KKS}})\overset{\mu=i}{\hookrightarrow }\mathfrak{g}^*$.
\end{example}

\subsection{Aplicações momento: quando existem e o que fazer com elas}

Nem sempre pode pegar um quociente simplético (a diferencia do mundo Riemanniano). No mundo simplético, temos que pegar conjuntos de nível da aplicação momento.

\begin{thm}[Marsden-Weinstein, Megre]\leavevmode
	$G\curvearrowright(M,\omega)$, ação hamiltoniana, $\mu:M\to \mathfrak{g}^*$ aplicação momento. Suponha que $0\in\mathfrak{g}^*$ é um valor regular. Então, pela equivariância de $\mu$ temos uma ação $G \curvearrowright \mu^{-1}(0)$:
	\begin{align*}
		x\in\mu^{-1}(0),&\qquad \mu(x)=0\\
		\mu(\psi_g(x))& =\operatorname{Ad}^* _g(\underbrace{\mu(x)}_{=0}=0
	\end{align*}
	(isso significa que tem uma ação de $G$ em $\mu^{-1}(0)$.) Suponha além disso que a ação é boa, i.e. \textit{regular} (definido na seguinte seção), que $\mu^{-1}(0)/G$ é suavem que a projeção  é sumersão. Então, existe uma úmica forma $\omega_{\operatorname{red}}$ simplética em $M_0$ tal que  $i^*\omega=\pi_0^*\omega_{\operatorname{red}}$. Issto é, a forma simplética é o pullback de uma forma no quociente,
	\[\begin{tikzcd}
		G\curvearrowright \mu^{-1}(0)\arrow[r,hook]\arrow[d,"\pi_0"]&M\\
M_0:=\mu^{-1}(0)/G
	\end{tikzcd}\]
\end{thm}

\begin{remark}\leavevmode
	Esse resultado tem muita história: primero o Meyer provou no Salvador, mas foi desapercebido. Daí esse mesmo ano Marsden-Weinsten provarom. Mas também parece que tudo esta relacionado com o Smale, que depois de mostrar a conjetura de Poincaré em dimensões altas estava tentando resolver o problema de treis corpos (talvez usando redução simplética).
\end{remark}

\subsection{Quocientes de ações por grupos de Lie}

\subsubsection{Criterio para ações regulares}

$G \curvearrowright P$, $G$ grupo de Lie, $P$ vareidade, $G/P$ espaço de órbitas com topologua quociente.
\[\begin{tikzcd}
P\arrow[d,"\pi"]\\
P/G
\end{tikzcd}\]
Mas ese espaço pode ser ruim. Por exemplo, $G=\mathbb{R}\curvearrowright R$ com $(t,x)\mapsto e^tx$ gera um espaço de treis órbitas: $R^-$ para $x<0$,  $\mathbb{R}^+$ para $x>0$ e  $0$ para  $x=0$. Mas cualquer aberto contendo  $0$  é todo $\mathbb{R}$: esse espaço nem é Hausdorff.

\begin{defn}\leavevmode
	$G \curvearrowright P$ é \textit{\textbf{regular}} se  $P/G$  é uma $\mathcal{C}^\infty$-variedade e $p\overset{\pi}{\longrightarrow}P/G$ é sumersão.
\end{defn}

\begin{thing2}{Critério}\leavevmode
	$G \curvearrowright P$.
	\begin{itemize}
	\item \textit{\textbf{Livre}}: $\forall  p \in P$, $G_p=\{e\}$ ($g\cdot p=p \implies g=e$).
	\item \textit{\textbf{Própria}}: a função
		\begin{align*}
			G\times P &\longrightarrow P\times P \\
			(g,p) &\longmapsto (g\cdot p,p)
		\end{align*}
		é \textit{\textbf{própria}} se a imagem inversa de compacto é compacto.
	\end{itemize}
\end{thing2}

\begin{example}\leavevmode
	$G$ compacto $\implies $ $G$-ação é própria.
\end{example}

\begin{thm}[Palais]\leavevmode
	$G \curvearrowright P$ livre, própria, então é regular.
\end{thm}

\begin{proof}[Idea da prova]\leavevmode
	Queremos construir uma estrutura suave em $P/G$. Pegue um ponto.  É uma órbita. Pegue uma subvariedade complementar, transversal, à órbita: $T_{x}S\oplus T_{x}\mathcal{O}=T_{x}P$. Faça agir o grupo nesa variedade transversal localmente. O resultado de uma vizinhança da órbita. Essa vai ser a carta. Usar teorema da função inversa.

	Também pode usar criterio de Godement (isso é mais geral). Uma variedade $M$, uma relação de equivalencia $R \subseteq M\times M$ (pode ser foleação, …). Então o quociente é variedade suave se e só se a relação de equivalência é fechada e mergulhada e a projeção a um dos fatores é sumersão. Então aqui a gente só prova que o mapa da definição de mapa próprio é um mergulho  fechado.
\end{proof}




\section{Aula 17: redução simplética (quociente simplético)}

Lembre:
\begin{itemize}
\item \textit{\textbf{Ação hamiltoniana}}: existe uma pliação momento $\mu:M\longrightarrow \mathfrak{g}^*$ satisfazendo duas propriedades: $i_{u_M}\omega=\partial \left<\mu,u\right> $ e $G$-equivariância de $\mu$, i.e. $\mu\cdot \psi_g=(\operatorname{Ad}^*)_g\mu$
\end{itemize}

Hoje: Redução simplética (quociente simplético)

\subsection{Uma proposição parecida ão exerício 5 da lista 3}

Precisaremos considerar a seguinte situação:
\[\begin{tikzcd}
G\curvearrowright P\arrow[d]\\
B=P/G
\end{tikzcd}\]
que é uma ação \textit{\textbf{regular}} (livre, própria).

\begin{question}\leavevmode
	Quando $\omega\in\Omega^{2}(P)$ é \textit{\textbf{básica}}, ie. da forma $\pi^*\bar{\omega}$, $\bar{\omega} \in\Omega^{2}(B)$?
\end{question}

\begin{defn}\leavevmode
	$\omega$ da forma $\pi^*\omega$ é dita \textit{\textbf{básica}} (com respeito à ação ).
\end{defn}

\subsection{$G$ conexo}


\begin{remark}\leavevmode
	$G$ conexo $\implies $ órbitas conexas $\implies $ $\pi$ é sumersão com fibras conexas. Nesse caso, (lista 3)
	\[\omega\text{ é básica} \iff\begin{cases}
		\mathcal{L}_{u_M}\omega=0,\qquad \forall u\in\mathfrak{g}\\
		i_{u_M}\omega=0
	\end{cases}\]
\end{remark}

\subsection{$G$ qualquer}

Primeiro lembre que

\begin{prop}\leavevmode
	\begin{enumerate}
		\item $\omega$ é basica se e somente se $\omega$ é {\color{1}$G$-invariante (em caso de $G$ ser conexo isso é equivalente a $\mathcal{L}_{u_M}\omega=0$}), e $i_{u_M}\omega=0$, i.e. $\ker \pi_* \subseteq \ker \omega$.

			E mais, se $\omega=\pi^*\bar{\omega} $, vale
		\item $ d \omega=0\iff d \bar{\omega} =0$
		\item $\bar{\omega} $ não degenerada $\iff$ $\ker \pi_* =\ker \omega$.
	\end{enumerate}
\end{prop}

\begin{prop}\leavevmode
	\begin{enumerate}
		\item Só tem um jeito de definir $\bar{\omega} $:
	\[\bar{\omega}_b (\bar{X},\bar{Y})=\omega_p(X,Y)\]
	onde $\pi_*(X)=\bar{X}$ e $\pi_*(Y)=\bar{Y}$ e $b=\pi(p)$. Daí só tem que verificar que $\bar{\omega} $ está bem definido:
	\begin{itemize}
	\item Não depende das escolhas de $X$ e $Y$. Fixando $p$, temos que
		\[X+\ker \pi_*,Y+\ker \pi_*\]
	\item Não depende de $p$ tal que $\pi(p)=b$. Aqui usamos que $\omega$ é $G$-invariante, {\color{4}essa é a única diferencia com o exercício 5 da lista 3}. Temos
		\[\pi(p)=\pi(p')\iff p=g\cdot p'\]
		\[\omega_p (X,Y)=\omega_{p'}(X',Y')\]
		\[X=d\psi_g(X'),\qquad Y=d\psi_g(Y)\]
		assim, $\overline{\omega}$ está bem definido, é suave e $\pi^*\overline{\omega}=\omega$ 
	\end{itemize}
	\item Como na lista 3
	\item Como na lista 3.
	\end{enumerate}
\end{prop}

\subsection{Redução simplética}

Pegue $G\overset{\psi}{\curvearrowright}(M,\omega)$, a aplicação momento, $x\in\mathfrak{g}^*$, e algumas outras coisas mais que não consegui escrever (é o escenário para:)

\begin{thm}[Marsden-Weinsten, Meyer]\leavevmode
{\color{4}Por cada elemento no dual da algebra de Lie, temos uma variedade pegando imagem inversa do mapa momento e quocientando.}

	$G\curvearrowright(M,\omega)$, $\mu:M\longrightarrow \mathfrak{g}^*$. Seja $\xi\in\mathfrak{g}^*$ valor regular de $\mu$, ie. $\mu^{-1}(\xi)\twoheadrightarrow M_\xi=\mu^{-1}(\xi)/G_\xi$. Suponha $G_\xi\curvearrowright \mu^{-1}(\xi)$ é uma subvariedade regular.

Então existe uma única forma simplética $\omega_{\operatorname{red}}$ em $M_\xi=\mu^{-1}(\xi/G_\xi$ tal que
\[\pi^* _\xi\omega_{\operatorname{red}}=i^*_\xi\omega\]
ie,
\[\begin{tikzcd}
	\mu^{-1}(\xi)\arrow[r,hook,"i_\xi"]\arrow[d,"\pi_\xi"]& M^\omega\\
	(M_\xi,\omega_{\operatorname{red}})
\end{tikzcd}\]
\end{thm}

\begin{remark}\leavevmode
	A diferencia com a outra vez que enunciamos o teorema é que quando pegamos um ponto arbitrário no $\mathfrak{g}^*$ temos que ter cuidado com qual é o grupo que é invariante na variedade para quocientar: pode não ser o $G$ tudo.
\end{remark}

\begin{proof}\leavevmode
	Só temos que checar que
	\[i^*\xi(\omega)\in\Omega^{2}(\mu^{-1}(\xi) \text{ é básica para $G_\xi$-ação} \]
	ou seja, de acordo com a proposição, queremos ao fim de contas, que
	\begin{itemize}
	\item $G_\xi$-invariante. (Isso é muito simples.)
	\item $\ker(\pi_\xi)_* =\ker(i^*_\xi\omega)$. Só falta ver isso. Note que
		\[T_x(G_\xi\cdot x)=\ker(\pi_\xi)_* |_{x},\qquad \ker(i_\xi^*\omega)|_{x}=T_x(\mu^{-1}(\xi))\cap T(\mu^{-1}(\xi))^\omega\]
		Agora vamos trabalhar com isso na forma de um lema:

	\begin{lemma}\leavevmode
		\begin{enumerate}
			\item $T_x(\mu^{-1}(\xi))=T_x(G\cdot x)^\omega$.
		
			\item $T_x(G_x\cdot x)=T_x(G\cdot x)\cap T_x(\mu^{-1}(\xi))$
		\end{enumerate}
	\end{lemma}
	Daí podemos pegar  $\omega$-complemento na igualadade 1., daí sustituimos na igualdade 2. e o teorema segue.

	\begin{proof}[do Lema]\leavevmode 
	\begin{enumerate}
		\item Isso é provar que $\ker d_x\mu=T_x(G\cdot x)^\omega$. Isso é verdade porque (propriedades do mapa momento)
			\begin{align*}
				x\in\ker d_x\mu&\iff \left<d \mu(X),u\right> =0\quad \forall u\\
				&\iff i_X \underbrace{d \left<\mu,u\right> }_{i_{u_M}\omega}=0\quad \forall u\\
				&\iff X\in T_x(G \cdot x)^\omega
			\end{align*}
		\item Por um lado,
			\[T_x(G_\xi\cdot x)=\{u_M(x)|u\in\mathfrak{g}\}=\{u_M(x)|u\text{ satisfaz }u_{\mathfrak{g}^*}(\xi)=0 \}\]
		Por outro lado, (equivariância de $\mu$)
		\[T_x(G\cdot x)=\{u_M(x)|u\in\mathfrak{g}\}\]
		\[\ker d_* \mu=T_x(\mu^{-1}(\xi)=\{X|d\mu(X)=0\}\]
		Intersectando ambos obtemos
		\[\{u_M(x)|d\mu(u_M(x)=0\}\]
		e pela equivariância de $\mu$, $d\mu(u_M(x)=0\iff u_{\mathfrak{g}^*}(\xi)=0$.
	\end{enumerate}
	\end{proof}
	\end{itemize}
\end{proof}


\subsection{Alguns comentários sobre o teorema}

\begin{thing4}{Para que serve?}\leavevmode
Por exemplo, para disminuir os graus de liberdade de um sistema Hamiltoniano em uma variedade. Também pode ser que o espaço de moduli seja esse quociente.
\end{thing4}

\begin{thing5}{Em vários casos a redução é "singular"}\leavevmode
	É quando a ação não é regular, não tem quociente suave. Aqui pode usar "espaços estatificados". Também pode usar um approach mas algébrico: redução "homológica". Isso último leva aos métodos BRST e BV.
\end{thing5}

\begin{exercise}\leavevmode
	$x\in M$, $\mathfrak{g}_x\subseteq \mathfrak{g}$, então vale
	\[\operatorname{Ann}(\mathfrak{g}_x) =\operatorname{img} d_x\mu\]
	daí, $\xi$ é um valor regular $\implies $ $\mathfrak{g}_x=\{0\}$ $\implies $ $G_x$  é discreta $\implies $ $G_\xi\curvearrowright \mu^{-1}(\xi)$ é localmente livre. Se $G_\xi$ é compacto, $M_\xi$ é um \textit{\textbf{orbifold simplético}}.
\end{exercise}

\begin{thing1}{Outro comentário}\leavevmode
	$\xi$ valor regular, $M_\xi=\mu^{-1}(\xi)/G_\xi$ onde $G_\xi$ age livremente em $\mu^{-1}(\xi)$.
	\begin{align*}
		\dim M_\xi&=\dim \mu^{-1}(\xi)-\dim G_\xi\\
		&=\dim M-\underbrace{\dim \mathfrak{g}^*}_{\dim G}-(\dim G-\dim \mathcal{O}_\xi)\\
		&=\dim M+\dim \mathcal{O}_\xi-2\dim G
	\end{align*}
\end{thing1}


\begin{thing2}{Shift trick}[Vai vir na lista]…\leavevmode
	A redução simplética depende da escolha do $\xi$. Se pegamos $\xi,\xi'$ na mesma órbita coadjunta, então $M_\xi\cong M_{\xi'}$. Tem mais: bota tudo em um espaço mais grande para mostrar que todo quociente é um quociente do zero naquele espaço maior… vai vir na lista
\end{thing2}

\begin{thing6}{Hamiltoniano em $M$ da hamiltoniano em $M_\xi$}\leavevmode
	Isso mem
\end{thing6}

\subsection{Exemplos simples}

\begin{thing4}{Um exemplo muito simlpes}\leavevmode
	
A reta agindo numa variedade simplética, que é o fluxo de algum campo $X$. A aplicação momento é o Hamiltoniano $H\in\mathcal{C}^\infty(M)$, $X_H=X$, a aplicação momento é $H:M\longrightarrow \mathbb{R}$. Isso daqui é uma ação hamiltoniana no caso mais simples.

Agora pegue um valor regular de  $H$, $\xi$. Se a ação $\mathbb{R}\curvearrowright H^{-1}(\xi)$ é boa, obtemos que $H^{-1}(\xi)/\mathbb{R}$ é uma variedade simplética com $\omega_{\operatorname{red}}$.
\end{thing4}

\begin{thing3}{Exemplo}\leavevmode
	Muito parecido: considere $R\dashrightarrow \mathbb{S}^1\curvearrowright \mathbb{C}^{n+1}=\mathbb{R}^{2n+2}$, a ação do círculo assim:
	\[e^{i\theta}(z_0,\ldots,z_n)=(e^{i\theta}z_0,\ldots,e^{i\theta}z_n)\]
Agora, pelo teorema da aplicação momento (?), temos que
\begin{align*}
	\mu: \mathbb{C}^{n+1} &\longrightarrow \mathbb{R} \\
	\mu(z) &=-\frac{1}{2}\|z\|^2\underbrace{+1/2}_{\text{opcional} }
\end{align*}
assim, $\mu^{-1}(0)/\mathbb{S}^1=\mathbb{C}P^{n},\omega_{\operatorname{red}}$, i.e., {\color{3}pelo teorema de Marsden-Weinsten, o $\mathbb{C}P^{n}$ tem uma estrutura simplética}.
\end{thing3}

\begin{thing4}{Mesmo mas mais geral}\leavevmode
	Agora pegue
	\[e^{i\theta}(z_0,\ldots,z_n)=(e^{ik_0\theta}z_0,\ldots,e^{ik_n\theta}z_n)\]
	Considere caso $n=1$, temos
	\[e^{i\theta}(z_0,z_1)=(e^{k\theta}z_0,e^{i\theta}z_1)\]
	obtemos
	\[\mu(z_0,z_1)=-\frac{1}{2}\left( k\|z_0\|^2+\|z_1\|^2 \right) \]
	Pegue $\xi<0$ valor regular. Se $k=1$ obtemos um quociente que  é topologicamente uma esfera, mas tem uma singularidade: a \textit{\textbf{gota de lágrima}}---uma esfera com um pontinho singular.
\end{thing4}

\section{Aula 18}

Lembre:
\begin{itemize}
\item Ação de um grupo de Lie numa variedade simplética, o mapa momento $\mu:M\longrightarrow \mathfrak{g}^*$.
\item Redução simplética
\item Exemplos
\end{itemize}

\subsection{Outros exemplos}

\subsubsection{Reduções de fibrados cotantentes}

Pegue uma variedade com uma ação de grupo de Lie, $Q\curvearrowright G$, a gente já viu (quando?) passa para o levantamento cotangente obtendo uma ação $G\curvearrowright T^*M$. Se a ação em $Q$ é livre e própria, também a ação no cotangente.

O mapa momento em $T^*Q$ é

\begin{align*}
	\mu:T^*Q  &\longrightarrow \mathfrak{g}^* \\
	\xi_x &\longmapsto \mu(\xi_x)\\
	\left<\mu(\xi_x,u\right> &:=\xi(u_Q(x))
\end{align*}


Então, quando reduzimos no $\xi=0$ obtemos
\[(T^* Q)_{\xi=0}=T^*(Q/G)\]

Mas isso é mais interessante para $\xi\neq 0$. Considere a  ação de um grupo de Lie $G$ nele mesmo pela multiplicação a esquerda, $G\overset{L_g}{\curvearrowright G}$. Note que $G/G=\operatorname{pt}$.

A aplicaco momento é
\begin{align*}
	\mu: T^*G &\longrightarrow \mathfrak{g}^* \\
	\xi &\longmapsto (d_eR_g)^*(\xi)
\end{align*}
Nesse caso, todo $\xi$ é valor regular e
\[\mu^{-1}(\xi)=\{(g,\xi^R|_{g}),\forall g\in G\}\]
Que é basicamente pegar uma forma e espalhar $\xi$ por todos lados usando o grupo. Fazendo redução obtemos
\[\mu^{-1}(\xi)/G_\xi\cong G/G_\xi\cong \mathcal{O}_\xi\]

\begin{exercise}[Vai vir na lista]\leavevmode
	De fato, $\omega_{\operatorname{red}}=\omega_{K K S}$. Mas ainda, toda órbita KKS pode ser vista desse jeito.
\end{exercise}

\begin{remark}\leavevmode
	Lembre que toda variedade de Poisson tem uma descomoposição em folhas simpléticas (isso é totalmente geral). No caso de $\mathfrak{g}^*$, essas folhas são as órbitas coadjuntas. Aqui temos que de fato  $\mathfrak{g}^*=(T^*G)/G$ quando pegamos o quociente do grupo de Lie $G$ levantado agindo em $T^*G$.

	Isso é mas geral. Pode ter qualquer variedade simplética com uma ação Hamiltoniana boa. Pega quociente: $M/G$  não tem por que ser simplético, mas ele tem uma estrutura de Poisson, e então uma descomposição em folhas que são simpléticas. O interessante é que essas folhas são da forma $\mu^{-1}(\xi)/G_\xi$, i.e. são órbitas coadjuntas.
\end{remark}

\subsubsection{Um espaço de moduli}

Considere alguns caras positivos assim: $\lambda_1,\ldots,\lambda_n>0$. Para $\lambda>0$, defina $\mathbb{S}^2_\lambda$ a esfera untária com forma simplética $\lambda\omega_{\operatorname{area}}$. Considere a variedade $M$ seguinte com a ação seguinte:
\[\mathsf{SO}(3) \curvearrowright \underbrace{\mathbb{S}^2_{\lambda_1}\ldots \times \mathbb{S}^2_{\lambda_n}}_{M}\]
e a hamiltoniana (tu pode ver que isso é certo)
\begin{align*}
	\mu: M &\longrightarrow \mathfrak{so}(3)^*=\mathbb{R}^3 \\
	(x_1,\ldots,x_n) &\longmapsto \lambda_1x_1+\ldots +\lambda_nx_n
\end{align*}
e agora considere o espaço reduzido em $\xi=0$,
\[\mathcal{M}_{\lambda_1,\ldots,\lambda_n}:=\mu^{-1}(0)/\mathsf{SO}(3) =\{(x_1,\ldots,x_n)\in(\mathbb{S}^2)^n|\lambda_1x_1+\ldots +\lambda_nx_n=0\} \Big/\mathsf{SO}(3)\]
isso tá bom pelo teorema. Vai ser o \textit{\textbf{espaço de polígonos espaciais em $\mathbb{R}^{3}$}}.

\begin{defn}\leavevmode
	Um \textit{\textbf{polígono espacial em $\mathbb{R}^{3}$}} é uns pontos $(v_1,\ldots,v_n)\in\mathbb{R}^{3}$, lados $(\ell_1,\ldots,\ell_n)$ que conectam vértices ciclicamente. Os $\ell$s são as diferencias dos $v$s. A condição do cara fechar é $\ell_1+\ldots +\ell_n=0$
\end{defn}

Consideremos o espaço de polígonos a menos de movimentos rígidos (traslações, rotações) no $\mathbb{R}^{3}$. Agora note
\begin{itemize}
\item Por traslação, podemos assumir que $v_1=0$.
\item Polígonos com $v_1=0$ são representados pelos lados $\ell_1,\ldots,\ell_n$ a menos de rotação.
\item Defina $\lambda_i=\|\ell_i\|$, $x_i=\ell_i/\lambda_i$.
\item obtemos $\mathcal{M}_{\lambda_1,\ldots,\lambda_n}$. É simplético e suave a menos de polígonos degenerados. Topologia dele: Kavitch Millson. "Moment maps and GIT quotients": C. Roodwand. Até é Kähler. Relação com outros espaços moduli como espaço de conexões planas na esfera com furos. Tem ação tórica. São exemplos de variedades quiver (representações).
\end{itemize}

O punto de tudo isso é que esse espaço aparece como uma redução simplética relativamente simples.

\subsubsection{Um exemplo de mecânica}

Lembre que tudo isso vem do trabalho do Smale quando ele tava estudando mecânica celeste.

Considere o problema de dois corpos. O espaço de fase de dois partículas é
\[Q=\{(r_1,r_2),r_i=(q_i^1,q_i^2,q_i^3)\} =\mathbb{R}^6\]
Daí,
\[M=T^* Q=\mathbb{R}^{12}=\{(r_1,r_2,P_1,P_2)\}\]
vamos denotar a estrutura canónica assim: $\omega=\sum dr_jdP_j$.

Vamos botar uma ação:
\begin{align*}
	G&=\mathbb{R}^{3}\curvearrowright Q=\mathbb{R}^{6}\\
	(r_1,r_2)&\longmapsto(r_1+g,r_2+g),\qquad g\in G=\mathbb{R}^{3}
\end{align*}
O levantamento cotangente da:
\[(r_1,r_2,P_1,P_2) \overset{g}{\mapsto }(r_1+g,r_2+g,P_1,P_2)\]
E a aplicação momento é
\begin{align*}
	\mu: \mathbb{R}^{12}  &\longrightarrow \mathbb{R}^{3} \\
	(r_1,r_2,P_1,P_2) &\longmapsto P_1+P_2
\end{align*}
O que é $\mu^{-1}/G_\xi$?

Vamos fazer uma mudança de coordenadas em $Q$ (um difeomorfismo):
\begin{align*}
	\mathbb{R}^{6} &\overset{\operatorname{\cong }}{\longrightarrow} \mathbb{R}^{6} \\
	(r_1,r_2) &\longmapsto (r_{\operatorname{cm}},r_{\operatorname{rel}})
\end{align*}
onde
\[r_{\operatorname{cm}}=\frac{m_1r_1+m_2r_2}{m_1+m_2}\qquad \qquad r_{\operatorname{rel}}=r_1-r_2\]
são o centro de masa e a posição relativa.

Você calcula que o levantamento cotange é
\begin{align*}
	\mathbb{R}^{12} &\longrightarrow \mathbb{R}^{12} \\
	(r_1,r_2,P_1,P_2) &\longmapsto \left( r_{\operatorname{cm}}:=r,r_{\operatorname{rel}}:=\bar{r},P_1+P_2:=P,\frac{1}{m_1+m_2}(m_2P_1-m_1P_2):=\bar{P} \right) 
\end{align*}
e
\[\omega=dr\cdot dP+d\bar{r}\cdot d\bar{P}\]
Temos uma ação levantada certo, as coordenadas dela nas novas coordenadas é
\[(r,\bar{r},p,\bar{p})\overset{g}{\longmapsto}(r+g,\bar{r},p,\bar{p})\]
A ideia é usar um sistema de coordenadas com o centro de masa na origem, e isso simplifica o negocio. Pode reduzir a dimensão para $\mathbb{R}^{6}$.

Agora vamos botar um hamiltoniano para transportar tudo ao reduzido. Lembre a redução de hamiltoniano é quando temos uma função $H\in\mathcal{C}^\infty(M)$ que é $G$-invariante, ie. constante ao longo das órbitas, assim está bem definido no quociente $\mu^{-1}(\xi)/G$. Pegue o Hamiltoniano
\[H(r_1,r_2,P_1,P_2)=\frac{1}{2}\left( \frac{|P_1|^2}{m_1}\frac{|P_2|^2}{m_2} \right) -\frac{Km_1m_2}{|r_1-r_2|}\]
é um hamiltoniano classico para o problema de dois corpos. Note que é $\mathbb{R}^{3}$-invariante olhando a definição da ação que dimos acima. Como ele vira no quociente? Assim:
\[H_{\operatorname{red}}(\bar{r},\bar{P})=\frac{1}{2}\frac{|P|^2}{\bar{m}}-\frac{K(m_1+m_2)\bar{m} }{\bar{r}}\]
Alí o que a gente consigui foi levar um problema de dois corpos a um problema de um corpo. Mas ainda, nesse quociente tem uma $\mathsf{SO}(3)$-ação, assim pode reduzir ainda mais, e resolver completamente o problema de dois corpos.

\subsubsection{Um exemplinho de mecânica com geometria riemanniana}

Porque achar geodésica é um problema de mecânica. $G$ grupo de Lie, $\left<\cdot,\cdot\right> $ uma métrica invariante à esquerda. $\left<\cdot,\cdot\right> $ é determinada um produto interno $\left<\cdot,\cdot\right> $ em $\mathfrak{g}$, que por sua vez é determinado por um operador simétrico autoadjunto $I:\mathfrak{g} \longrightarrow \mathfrak{g}^*$ de maneira que
\[\left<u,v\right>_{g}=\left<I(u),I(v)\right> \]
i.e. é o mesmo ter esse produto interno que ter a outra coisa.

Considere $G\curvearrowright G$, $G\curvearrowright T^*G$, $\mathcal{O}_\xi=\mu^{-1}(\xi)/G_\xi$ para $\xi \in\mathfrak{g}^*$.

O sistema hamiltoniano reduzido em $\mathcal{O}_\xi$ é
\[\dot \eta(t)=-\operatorname{ad}^*_{I^{-1}(\eta(t))}(\eta(t))\]
que é a equação de Euler, uma versão reduzida das equações das geodésicas.
\[dH_{\operatorname{red}}|_{\eta}=I^{-1}(\eta)\]
\[H_{\operatorname{red}}=\frac{1}{2}(\eta(I^{-1}(\eta))\]
a parte da mecânica entra quando consideramos $G=\mathsf{SO}(3)$. Pense num corpo que tá voando por todos lados mas tem um ponto fixo num lugar. A solução desse corpo rígido é dada por geodésicas no $\mathsf{SO}(3)$.

O que acontece é que
\[I=\begin{pmatrix} I_1&0&0\\0&I_2&0\\0&0&I_3 \end{pmatrix} ,\qquad I_i>0\]
é o \textit{\textbf{momento de métrica}} cujos autovalores se chaman de \textit{\textbf{valores principais?}}. Esse cara  é o que vai determinar a métrica. No caso desse corpo rígido, temos a \textit{\textbf{equação de Eulor para o corpo rígido}}:
\[\frac{d\eta_i}{dt}(t)=(I_k^{-1}-I_j^{-1})\eta_j\eta_k\]
\[\dot \eta(t)=\eta\times I^{-1}\eta\]
que são esferas. As soluções da equação de Euler é um negocio bem bonito: pega uma esfera, tem 6 pontos de equilíbrio que estão em pares, dois em cada autoespaço. Quatro estáveis, dos inestáveis. {\color{4}É como rodar um livro no eixo diagonal…}

\begin{remark}\leavevmode
	Tem muitíssimas EDPs relacionadas a equações  reduzidas de geodésicas quando $G$ é de dimensão infinita. Por exemplo $G=\operatorname{Dif}_{\operatorname{Vol}}(M,\left<\cdot,\cdot\right> )$. Por exemplo a de Burger aparece com $M=\mathbb{S}^1$.
\end{remark}

\section{Aula 19: comentários de possíveis projetos}

Lembre:
\begin{itemize}
\item Ações hamiltonianas e redução.
\item Aplicação momento, estrutura simplética nas órbitas da ação coadjunta.
\item Diversos exemplos.
\end{itemize}

Hoje:
\begin{itemize}
\item Super exemplo.
\item $G=\mathbb{T}^n$ toro.
\end{itemize}

\subsection{Super exemplo (de ação hamiltoniana e redução simplética)}

Advertencia: os objetos não tem sentido mas são expressoes de objetos geométricos.

{\color{5}\bfseries Objetos.}\hspace{.5em} $\Sigma$ variedade compacta conexa orientada, $G$compacto grupo de Lie, $\mathfrak{g}$ munida de um produto interno $\left<\cdot,\cdot\right> $ $\operatorname{Ad}$-invariante.
\begin{itemize}
\item $\mathcal{A}=\Omega^{1}(\Sigma,\mathfrak{g})$. Formas que tomam valores em $\mathfrak{g}$ em lugar de em $\mathbb{R}$. \textbf{Interpretação geométrica:}  espaço das conexões em 
	\[\begin{tikzcd}
	\Sigma\times G\arrow[d]\\
	\Sigma
	\end{tikzcd}\]
	\begin{remark}[Conexões]\leavevmode
		\[\begin{tikzcd}
		G\curvearrowright P\arrow[d,"\pi"]\\
		B=P/G
		\end{tikzcd}\]
		com ação livre, própria. Temos o \textit{\textbf{fibrado vertical}}, que  $\ker \pi$, o espaço tangente às fibras de $\pi$. Uma \textit{\textbf{conexão}} é a escolha de um subfibrado $H$ que complemente o fibrado vertical, dizemos que é \textit{\textbf{horizontal}}, e que  é $G$-invariante. Em soma, $TP=V\oplus H$ e $g_* H=H$.

		A ideia é que o fibrado horizontal permite conectar fibras, transporte paralelo.

		Comparando com a noção de conexão em geometria riemanniana, uma conexão em
		\[\begin{tikzcd}
		E\arrow[d]\\
		M
		\end{tikzcd}\]no sentido riemanniano, é uma conexão no sentido daqui em
		\[\begin{tikzcd}
		\operatorname{Fr}(E)\arrow[d]\\
		M
		\end{tikzcd}\]onde $\operatorname{Fr}(E)$ é o bundle de bases das fibras de $E$, que é um fibrado principal com a ação de $\mathsf{GL}(n)$.
	\end{remark}
	A ideia será que os núcleos das formas em $\mathcal{A}$ são conexões. Em fim, é só para menter em mente que esse $\mathcal{A}$ define conexoes em $\Sigma\times G$.

	\item  $ \mathcal{G}=\mathcal{C}^\infty(\Sigma,G)$. \textbf{Interpretação geométrica}: é um \textit{\textbf{grupo de calibre}} quando temos
		 \[\begin{tikzcd}
			G\curvearrowright P\arrow[r]\arrow[d]&P\curvearrowleft G\arrow[d]\\
			\Sigma\arrow[r,"\operatorname{id}"]&\Sigma
		\end{tikzcd}\]

	\item $\mathcal{G}\curvearrowright \mathcal{A}$, $gA=\operatorname{Ad}_g(A)-g^*\theta^R$. \textbf{Interpretação geométrica}: se chama de \textit{\textbf{ação de calibre}}. São transformações  muito naturais no mundo dos grupos de calibre.

	\item $\Omega^{1}(\Sigma,\mathfrak{g}) \to \Omega^{2}(\Sigma,\mathfrak{g})$. Dada por  $A\mapsto dA+\frac{1}{2}[A,A]$. Aqui $[A,A](X,Y):=[A(X),A(X)]-[A(Y),A(X)]$

		\begin{remark}[Curvatura]\leavevmode
			A importança da conexão e conectar fibras. É um levantamento de campos vetorias na base. Como muda o colchete? Isso é o que mede a curvatura. Se a curvatura é zero, é plana, a distribuição no espaço total é integrável.
		\end{remark}
\end{itemize}

{\color{8}\bfseries Portanto}\hspace{.5em}$\mathcal{A}$ é uma variedade simplética com a forma, para $A\in A$ e $u,v\in T_A \mathcal{A}\cong \mathcal{A}=\Omega^{1}(\Sigma,\mathfrak{g}) $
\[\omega_A(u,v)=\int_{\Sigma}\left<u,v\right> \]
onde definimos $\left<u,v\right>(X,Y)=\left<u(X),v(Y)\right> -\left<u(Y),v(X)\right> $ ussando o produto interno que botamos em $\mathfrak{g}$.

{\color{3}\bfseries Identificações.}\hspace{.5em}
\begin{itemize}
\item $\mathsf{Lie}(\mathcal{G}) =\Omega^{0}(\Sigma,\mathfrak{g})$.
\item $\mathsf{Lie}(\mathcal{G})^* =\Omega^{2}(\Sigma,\mathfrak{g})$.
\end{itemize}

\begin{thm}[Atiyah-Bott]\leavevmode
	$\mathcal{G}\curvearrowright \mathcal{A}$ é hamiltoniana. A aplicação momento é 
	\begin{align*}
		\operatorname{curv}: \mathcal{A}=\Omega^{1}(\Sigma,\mathfrak{g}) &\longrightarrow \mathsf{Lie}(\mathcal{G})^* =\Omega^{2}(\Sigma,\mathfrak{g}) \\
	\end{align*}
	Fazendo redução, de acordo com o teorema de Marsden-Weinsten,
	\[\operatorname{curv}^{-1}(0)\Big/\mathcal{G}=\mathcal{M}\]
	é um espaço simplético, o espaço moduli de conexões planas.
\end{thm}

{\color{8}\bfseries Onde é que isso aparece?}\hspace{.5em}

\begin{itemize}
\item Por qué esse espaço tem dimensão finita? porque
	\[\mathcal{M}=\dfrac{\operatorname{Hom}(\pi_{1}(\Sigma),G)}{G}\]

\item $\mathcal{M}$ e o espaço de fase de \textit{\textbf{Chern-Simons}} (física).

\item Em topologia, esse espaço da invariantes de 3-variedades. Corta uma 3-variedade. A intersação é uma superficie. Calcula o seu $\mathcal{M}$. Olha para as conexoes que são restrição es de conexoes planas em $M^3$ de um lado do plano com que cortou. Isso é uma subvariedade lagrangiana de $\mathcal{M}$. Otro lado. Interseção das lagrangianas é invariante simplético.

\item Se a variedade tem bordo, $\mathcal{M}$ é só de Poisson. Tem um negocio que não consegui entender.

\item Também pode botar mas estrutura. Em alguns casos $\mathcal{M}$ é Kähler.
\end{itemize}

\subsection{Aplicação de redução: cortes simpléticos}

Pegue uma variedade simplética conexa com uma ação hamiltoniana de $S^1$. $S^1\curvearrowright(M,\omega)\overset{H}{\longrightarrow}\mathbb{R}$. Considere $M\times \overline{\mathbb{C}}\cong \mathbb{R}^{2}$. Alí também pode agir com $S^1$ onde $S^1$ é rotação em $\mathbb{C}$.Você vai ver na lista que em casos assim, a aplicação momento é $H(x)-\frac{1}{2}|z|^2$.

Agora pode olhar para
\[\mu^{-1}(\xi)= \underbrace{H^{-1}(\xi)\times \{0\}}_{\cong H^{-1}(\xi)}\sqcup \left( \bigsqcup_{r>0}H^{-1}(\xi+r)\times  \left\{ \frac{1}{2}|z|^2=r \right\}  \right) \]
Esse segundo conjunto é também $=\{(x,z)|H(x)>\xi,z=e^{i\theta}\sqrt{H(x)-\xi} \} \cong \{H>\xi\} \times S^1$.

Agora vamos supor que $\xi$ é um valor regular de $H$, e $S^1\curvearrowright H^{-1}(\xi)$ é livre, própria. Mesmo para $\mu$. Se consigo fazer redução simplética para $H$ e $(M,\omega)$, vou conseguir também para $M\times \overline{\mathbb{C}}$. Vamos:
\[\mu^{-1}(\xi)/S^1\cong H^{-1}(\xi)/S^1\sqcup\left( \dfrac{\{H>\xi\} \times S^1}{S^1} \right) \cong M_\xi\sqcup \{H>\xi\}\]
Qué é que estou dizendo? A variedade simplética $\mu^{-1}(\xi)/S^1$ é uma união de dois coisas: níveis da aplicação momento, e, no bordo $\mu^{-1}(\xi)$ fazemos redução para ter também estrutura simplética. Também poderiamos olhar para niveis que estão abaixo de $\xi$. Uma variedade podemos chamar $M_{\geq \xi}$ e a outra $M_{\leq \xi}$. Então separamos $M$ em duas variedades ao longo de $\mu^{-1}(\xi)$.

{\color{1}\bfseries Colagem simplética (Symplectic glueing).}\hspace{.5em} O Eugene Lerman introduiz essa operação muito linda para desfazer o processo anterior (acho, não será que o Lerman introduiz o corte simplético mesmo?). Isso é chamado \textit{\textbf{blow-up simplético}}, onde tiramos não um ponto senão uma bola e pegamos o fibrado de Hopf…

\subsection{Dois resultados de ações hamiltonianas por toros}

Lembre: tudo isso é o que a gente não vai falar.

Vamos falar rapidamente de convexidade. $(M,\omega)$ simplética, conexa e compacta. Ação hamiltoniana de algum toro $\mathbb{T}^n\curvearrowright(M,\omega)$ $\mu:M\to \mathbb{R}^{m}$.

\begin{thm}[Atiyah/Guilleimen-Steinberg]\leavevmode
	A imagem da aplicação momento é convexa.

	De fato, a imagem da aplicação momento é o fecho convexo da imagem dos pontos fixos da ação. Esse é o \textit{\textbf{politopo momento}}.
\end{thm}

\begin{example}[Que de fato motivou isso, Horn-Schur]\leavevmode
	Pegue $\lambda_1,\ldots,\lambda_n$ é olha as matrizes hermitianas que tem esse especto, $\mathcal{O}_{\lambda_1,\ldots,\lambda_n}$. Quer saber a diagonal, ela fica na envoltura conexa dos $\lambda_{i}$ (acho).
\end{example}

\subsection{Variedade simpléticas tóricas}

\begin{defn}\leavevmode
\textit{\textbf{Variedade simplética tórica}}: $M^{2n}$ variedades conexas compactas com ação hamiltoniana (efetiva=qualquer elemento do grupo move pelo menos um elemento) de $\mathbb{T}^n$.	
\end{defn}

\begin{example}\leavevmode
\begin{itemize}
	\item $S^1\curvearrowright S^2$ por rotações aplicação momento é função altura (polo norte, polo sur). O politopo momento é o intervalo $[-1,1]$.
\item $\mathbb{T}^2\curvearrowright \mathbb{C}P^{2}$ dada por
	\[(e^{i\theta_1},e^{i\theta_2})([z_1:z_2:z_3])=([z_0,e^{i\theta_1}z_1,e^{i\theta_2}z_2])\]
	a aplicação momento resulta ser
	\[\mu([z_0,z_1,z_2])=-\frac{1}{2}\left( \frac{|z_1|^2}{|z_0|^2+|z_1|^2+|z_2|^2} ,\ldots\right) \]
	e o politopo momento é um triangulsinho conectando $(-1/2,0)$ e  $(0,-1/2)$.
\end{itemize}\end{example}

Para variedades tóricas,
\begin{itemize}
\item É possível caracterizar os politopos que são imagem de aplicaçã oes momento, se chamam de \textit{\textbf{politopos de Delzant}}.
 \item As variedades simpléticas tóricas estão classificadas opr politopos de Delzant
\end{itemize}

\begin{thm}[Delzant]\leavevmode
	\[\left\{ \text{variedades simpléticas tóricas}  \right\} \overset{1-1}{\longleftrightarrow} \{\text{politopos de Delzant} \}\]
\end{thm}

\begin{remark}\leavevmode
	Cortar variedades simpléticas como falamos antes corresponde com cortar o politopo!
\end{remark}

\begin{remark}\leavevmode
	Olha para um livro de Guilleimin. A Anna pede para fazer muitos exercícios se quer  entender bem. Também o Audin.
\end{remark}

\section{Aula 20: Exitência e unicidade de aplicações momento}

\begin{thing8}{Projeto final}\leavevmode
	Para este domingo: tema e uma ou duas referencias. Aproximadamente um mes depois (22 novembro): primeiro rascunho. Última semana de novembro e primeira de dezembro: apresentações. Uns 40 minutos de apresentação. Não tem que apresentar tudo que está no trabalho escrito.

	

\end{thing8}

\subsection{A primeira obstrução : um diagraminha commutativo}

Lembre:
\begin{itemize}
\item $G\curvearrowright(M,\omega)$ ação simplética é \textit{\textbf{Hamiltoniana}} se exite  $\mu:M\to \mathfrak{g}^*$ ou, equivalentemente, $\hat{\mu}:\mathfrak{g} \to \mathcal{C}^\infty(M)$ tal que
	\begin{enumerate}
		\item (Diagraminha commutativo aqui embaixo)
	
		\item (Seguinte seção)
	\end{enumerate}
\end{itemize}

Queremos achar uma aplicação $\hat{\mu}:\mathfrak{g} \to \mathcal{C}^\infty(M)$ tal que
\[\begin{tikzcd}
&\mathcal{C}^\infty(M)\arrow[dr]\\
\mathfrak{g}\arrow[rr,""]\arrow[ur,"\hat{\mu}"]&&\mathfrak{X}(M)
\end{tikzcd}\]
onde o mapa $\mathcal{C}^\infty(M)\to  \mathfrak{X}(M)$ é $f\mapsto X_g$.
\begin{prop}\leavevmode
	Existe $c_1 \in [\mathfrak{g},\mathfrak{g} ]^0\otimes H^{1}(M)$ tal que 
	\[\hat{\mu}:\mathfrak{g} \to \mathcal{C}^\infty(M)\text{ satsifazendo o diagrama anterior} \iff c_1=0 \]
	onde $[\mathfrak{g},\mathfrak{g} ]^0\subset \mathfrak{g}^*$ é o anulador do commutador de $\mathfrak{g} $.
\begin{proof}\leavevmode
	A existencia de $\hat{\mu}$ é equivalente a que a imagem do mapa em baixo seja de campos hamiltonianos: pegue uma base de $\mathfrak{g}$, chega a campos hamiltonianos e escolhe uma função para cada um deles. Alí já determinou um $\hat{\mu}$.

	Considere
\[\begin{tikzcd}
	\mathfrak{g} \arrow[rr,bend right,"\bar{\psi}" ]\arrow[r,"\psi"]&\mathfrak{X}(M)\arrow[r]&\mathfrak{X}_\omega(M)/\mathfrak{X}_{\operatorname{Ha m}}(M)\cong H^{1}_{\operatorname{dR}}(M)
\end{tikzcd}\]
Note que $\bar{\psi}([\mathfrak{g},\mathfrak{g} ] =0$ já que a gente viu no principio do curso que $\mathfrak{X}_\omega(M),\mathfrak{X}_\omega(M)\subseteq \mathfrak{X}_{\operatorname{HAM}}(M)$. Daí a aplicação induzida é
\[\mathfrak{g} /[\mathfrak{g},\mathfrak{g} ]\overset{c_1}{\longrightarrow}H^{1}(M)\]
se e somente se
\[c_1\in\left( \frac{\mathfrak{g}}{[\mathfrak{g},\mathfrak{g} ]} \right)^*\otimes H^{1}(M) \]
já que em geral $T:V\to W$ linear se e somente se $T\in V^*\otimes W$
\end{proof}
\end{prop}

\begin{remark}\leavevmode
	 Uma ação simplética fracamenta hamiltoniana satisfaz isso (acho).
\end{remark}

\subsubsection{Interpreatação coholomogica}
$g$ álgebra de Lie,
\[C^k=\Lambda^k\mathfrak{g}^*,\qquad \delta:\Lambda^k\mathfrak{g}^*\to \Lambda^{k+1}\mathfrak{g}^*\]
\[\delta_C(u_0,u_k)= \sum_{i<j}(-1)^{i+j}\subset([u_i,u_j],u_1,\ldots,\hat{u}_i,\hat{u}_j,\ldots,u_k)\]
que é o operador de \textit{\textbf{Chavalley-Eilenberg}}.

\begin{remark}\leavevmode
	$\Lambda^k\mathfrak{g}^*\cong \Omega^k(G)^G\subseteq \Omega^{k}(G)$ onde esse na esqueda são as formas $G$-invariantes. Quando $G$ é compacto tem igualdade
\end{remark}

 \begin{exercise}\leavevmode
	 $\delta^2=0\iff$ identidade de jacobi em $[\cdot,\cdot ]$. Desse jeito é equivalente dar um espaço vetorial que a forma colchete seja jacobi $(V,[\cdot,\cdot ])$ ou pedir que $\Lambda^\bullet V^*$ tenha uma derivação $\delta$  tal que $\delta^2=0$
\end{exercise}

\begin{remark}\leavevmode
	$\delta:\mathfrak{g}^*\to \Lambda^2\mathfrak{g}^*$ é dual (com signo $-$) de $[\cdot,\cdot ]:\Lambda\to \mathfrak{g}$,
	\[\delta c[u,v]=-c([u,v])\]
	tem o colchete tem delta, tem delta tem colchete, é só isso.
\end{remark}

\begin{defn}\leavevmode
	\[H^{2}(\mathfrak{g})=\frac{\ker \delta}{\operatorname{img} \delta}\]
\end{defn}

\begin{remark}\leavevmode
	\begin{itemize}
	\item Em grau zero $\delta$ é zero: $\delta:\mathbb{R}\to \mathfrak{g}^*$, $\delta=0$
	\item Em grau um, $\delta:\mathfrak{g}^*\to \Lambda^2\mathfrak{g}^*$, $\delta c(u,v)=-c([u,v])$.
	\end{itemize}
	Assim,
	\[H^{1}(\mathfrak{g})=\ker \delta=[\mathfrak{g},\mathfrak{g} ]^0\]
\end{remark}
Por tanto, o resultado da seção anterior pode ser escrito como
\[c_1\in H^{1}(\mathfrak{g})\otimes H^{1}(M)\]

\subsection{A segunda obstrução: equivariância}
A outra condição da aplicação momento é que
\[\hat{\mu}([u,v])=-\{\hat{\mu}(u),\hat{\mu}(v)\}.\]

Vamos supor que conseguimos resolver a primeira obstrução. Vimos que não existe só uma função.  Agora vamos ver cómo fazer uma escolha de $\hat{\mu}$ que permita satisfazer a segunda obstrução.

Fixe $\hat{\mu}$ satisfazendo a primeira obstrução. Considere a \textit{\textbf{falha}}:
\[c^\mu_2=\hat{\mu}([v,u])-\{\hat{\mu}(u),\hat{\mu}(v)\} \in\mathcal{C}^\infty(M)\]
for $u,v\in\mathfrak{g}$.

\begin{remark}\leavevmode
	\[X_{c^\mu_2(u,v)}=\underbrace{\psi([v,u])}_{=[\psi(u),\psi(s)]}-\underbrace{X_{\{\hat{\mu}(u),\hat{\mu}(v)\}}}_{[X_{\hat{\mu}(u)},X_{\hat{\mu}(v)}}=0\]
	asi que $c_2^\mu(u,v)$ é constante.

	Se $M$  é conexo, $c_2^\mu(u,v)\in\mathbb{R}$. temos um operador antisimétrico
	\[c_2^\mu:\mathfrak{g} \times \mathfrak{g} \to \mathbb{R}\]
	i.e., um elemento $c_2^\mu\in \Lambda^2\mathfrak{g}^*$, i.e. uma classe $[c_2^\mu\in H^{2}(\mathfrak{g})$.
\end{remark}

\begin{thing4}{Lema (Exercício)}\leavevmode
	$\delta c^\mu_2=0$. Segue da identidade de Jacobi
\end{thing4}

\begin{lemma}\leavevmode
	$[c_2^\mu]$ não depende de $\hat{\mu}$ (só depende da ação simplética).
\end{lemma}

\begin{proof}\leavevmode
	Suponha $\hat{\mu}_1,\hat{\mu}_2$ que resolvem a primeira obstrução. Bom isso implica que $X_{\hat{\mu}_1}(u)=X_{\hat{\mu}_2(u)}\implies \hat{\mu}_1(u)-\hat{\mu}_2(u)\in\mathbb{R}$ para toda $u\in\mathfrak{g}$. E isso implica que $b:=\hat{\mu}_1-\hat{\mu}_2\in\mathfrak{g}^*$.

Então,
\[(c_2^{\mu_1}-c_2^{\mu_2}=\hat{\mu}_1([v,u])-\hat{\mu}_2([v,u])- \underbrace{\{\hat{\mu}_1(u),\hat{\mu}_1(v)\} +\{\hat{\mu}_2(u),\hat{\mu}_2(v)\}\}}_{=0}=\delta b(u,v)\]
Agora lembre que por definição  de $\delta$, $\delta b (u,v)=\big([v,u])$. Isso implica que
	\[ [c_2^{\mu_1}]=[c_2^{\mu_2}=c_2\in H^{2}(\mathfrak{g})\]
\end{proof}

\begin{prop}\leavevmode
	Seja $G\curvearrowright(M,\omega)$ fracamente hamiltoniana (ou seja que a primeira obstrução ta resolvida) ($\iff c_1-0$). Então existe  $\hat{\mu}$ satisfzendo as duas obstruções  $\iff c_2=0$.
\end{prop}

\begin{thing8}{Ou seja, $c_2$ mede a segunda obstrução}\leavevmode
\end{thing8}

\begin{proof}[Demostração]
	\begin{itemize}
	\item Se $\hat{\mu}$ satisfaiz a segunda obstrução $\implies c_2^\mu=0\implies c_2=0$.
	
		\item Se $c_2=0$, então existe  $\hat{\mu}$ com $c_2^\mu=\delta b$. Tome
			\[\hat{\mu}'=\hat{\mu}-b\text{ satisfaz as duas obstruções} \]
	\end{itemize}
	{\color{3}\begin{quotation}
		ou seja, se tem um cada que resolve a primeira obstrução e o $c_2$ é $\delta b$, então modifica a $\hat{\mu}$ tirando $b$ e já com isso vai resolver também a segunda obstrução.
	\end{quotation}}
\end{proof}

\begin{thing5}{Consequencia}\leavevmode
	$M$ y $G$ conexos. Toda ação simplética em $(M,\omega)$ satisfazendo
	\begin{itemize}
	\item $H^{1}(\mathfrak{g})=0$.
	\item $H^{2}(\mathfrak{g})=0$
	\end{itemize}
é hamiltoniana {\color{3}com única aplicação momento (ainda falta falar de unicidade)}.
\end{thing5}

\begin{remark}\leavevmode
	$\mathfrak{g}$ semisimple implica $H^{1}(\mathfrak{g})=0$ e $H^{2}(\mathfrak{g})=0$. Isso é um caso particular dos lemas de Whitehead.
\end{remark}

Ou seja que tudo isso para ver quando existe a aplicação momento.

\subsection{Faltou falar de unicidade}

Suponha $\mu_1,\mu_2$ são aplicações  momento. $b=\hat{\mu}_1-\hat{\mu}_2\in\mathfrak{g}^*$, $M$ conexo. Então, pela equivarância, $b([u,v])=0$, i.e. $b\in [\mathfrak{g},\mathfrak{g} ]^0=H^{2}(\mathfrak{g})$.

Por outro lado, $\mu$ aplicação momento equivarânte e pegamos $b\in [\mathfrak{g},\mathfrak{g} ]^0$, então $\hat{\mu}+b$ é uma nova aplicação momento equivariante. Portanto,

{\color{8}\begin{quotation}
	aplicações  momento são únicas modulo $H^{1}(\mathfrak{g})$.
\end{quotation}}

\subsection{Comentarios finais}

\begin{itemize}
\item $G$ abeliano (e.g. $G=\mathbb{T}^n$. Lembre que no príncipio do curso a gente viu que rotar ão longo de um dos círculos do toro da um campo que é simplético mas não hamiltoniano. Essa ação não admite uma aplicação momento, i.e. $ H^{1}(\mathbb{T}^2)\neq 0$).
	\begin{itemize}
	\item Ação simplética não precisa ser hamiltoniana.
	\item Se  aplicação momento existe, está definida a menos de $\mathfrak{g}^*$. (caso de um toro tem um $\mathbb{R}^{n}$ de escolhas de aplicações  momento; $H^{1}(\mathfrak{g})$ é $\mathbb{R}^{n}$ acho).
	\end{itemize}

\item Existem condições em $M$ garantindo que toda ação simplética é hamiltoniana. Por exemplo, se $H^{1}(\mathfrak{g})=0$ e além disso $M$ conexa  e compacta.
\end{itemize}

\section{Aula 21}

\subsection{Hiperfície de contato}

\begin{defn}\leavevmode
	Uma \textit{\textbf{variedade de contato}} $M^{2n-1}$ é uma variedad com uma \textit{\textbf{1-forma de contato}} $\alpha$, ou seja, $\alpha\wedge (d\alpha)^{n-1}$ é uma forma de volume (diferente de zero em todo ponto). Chamamos $\xi=\ker \alpha$ de \textit{\textbf{estrutura de contato}}.
\end{defn}

Note que $d\alpha|_{\xi_p}$ é não degenerada. $\alpha |_{\ker d\alpha}$ é não nula pelo fato de que $\alpha\wedge (d\alpha)^{n-1}$. Aí podemos dividir essa distribuição de dimensão 1 no lado onde $\alpha$ é  positiva e onde é negativa.

\begin{defn}\leavevmode
	O \textit{\textbf{campo de Reeb}} de $\alpha$ é o campo vetorial $R_\alpha$ únicamente caracterizado pelas equações 
	\begin{align*}
		i_{R_\alpha}d \alpha&=0,\qquad \alpha(R_\alpha)\equiv 1.
	\end{align*}
\end{defn}

\begin{remark}\leavevmode
	Esee é um campo canônicamente associado à forma de contato. Não tem algo assim em variedades simpléticas.
\end{remark}

\begin{remark}[dani]\leavevmode
Uma variedade de contato tem uma estrutura de contato que é uma distribuição $\mathcal{C}=d\alpha$. $d\alpha$ é simplética em $\mathcal{C}$. O fato de $i_{R_\alpha}d\alpha=0$ significa que $R_\alpha|_{\mathcal{C}}=0$ porque essa forma simplética $d\alpha$ é não degenerada. Então como $R_\alpha$ não tá em $\mathcal{C}=\ker \alpha$, ele fica determinado pela condição $\alpha(R_\alpha)\equiv 1$.
\end{remark}

\begin{thing4}{Referência}[Para quem quiser estudar mais sobre isso]\leavevmode
	Geiges, \textit{Introduction to Contact Topology} 
\end{thing4}

Agora pegue $(M^{2n},\omega)$ simplética e $S\subset M$ hiperfície.

\begin{defn}\leavevmode
	$S$ é uma \textit{\textbf{hiperfície de contato}} se  $S$ admite uma forma de contato $\alpha$ tal que $d\alpha=\omega|_{S}=i^*\omega$.
\end{defn}

\begin{thm}\leavevmode
	Uma hiperfície de contato compacta $S$ de $M$ é de contato se e somente se existe um campo vetorial $X$ definido em uma vizinhança $U$ de $S$ tal que
	\begin{enumerate}
		\item $X$ é \textit{\textbf{conformemente simplético}}, ou seja, $\mathcal{L}_X\omega=\omega$.
		\item $X\pitchfork S$ ($X$ é transversal a $S$)
	\end{enumerate}
\end{thm}

\begin{exercise}\leavevmode
	 Mostre que  $\mathcal{L}_X\omega=\omega\iff \varphi_t^*\omega=e^t \omega$.

	 Compare com o caso simplético: aí a gente tinha que $\mathcal{L}_X\omega=0$, i.e. que o fluxo puxa $\omega$ em $\omega$. Aqui estamos multiplicando por factor, e por isso dizemos que é conforme.
\end{exercise}

\begin{proof}[Prova do teorema]\leavevmode
	($\impliedby$) Defina $\alpha=i_X\omega$. Vamos ver que $\alpha$ restringida a $S$ é uma forma de contato. Pela fórmula de Cartan,
	\begin{align*}
		d\alpha&=\mathcal{L}_X\omega=\omega
	\end{align*}
	Como $S$ possui dimensão ímpar $\exists T\in T_xS$ tal que $\omega_x(Y,v)=0$ $\forall  x\in S$ e $\forall v \in \in T_xS$. Como $\omega$ é não degenerada, $\alpha_x(Y)=\omega_x(x),Y)\neq 0$.

	Portanto,
	\[D_s=\{v\in T_xS;\omega_x(X(x),v)=0\}\]
	define uma distribuição de hiperplanos em $S$ complementar a $Y$. (Note que $D_x$ é o ortogonal simplético do plano simplético $ \operatorname{s p a n}\{X(x),Y\}$.)

Consequentemente $\alpha\wedge (d\alpha)^{n-1}\neq 0$ em todo ponto, i.e., forma de volume.

($\implies $) Vamos contruir uma primitiva de $\alpha$ i.e. que $d\alpha$ seja $\omega$. Vamos precisar do seguinte

\begin{lemma}\leavevmode
	Existe uma vizinhança $U$ de $S$ e uma 1-forma $\tau$ em $U$ tal que $\tau$ é uma primitiva de $\omega$, i.e. $d\tau=\omega$ e $i^*\tau=\alpha$
\end{lemma}

\begin{proof}[Prova do lema]\leavevmode
	Pelo teorema da vizinhança tubular, existe uma vizinhança $U$ de $S$ e um difeomorfismo \[\psi:S\times (-1,1)\to U\] tal que \[\psi(x,0)=x\qquad  \text{ para toda}x \in S\]. Compondo a inversa de $\psi$ com a projeção no primeiro fator temos um mapa $r:U\to S$.

	Seja $\mu=i^*\alpha$. Temos que $i^*\mu=i^*r^*\alpha=\alpha$ pois $\underbrace{\pi_1 \circ \psi^{-1}}_{r}\circ i=\operatorname{Id}$.

	Considere a 2-forma
	\[\eta:=\omega-d\mu\]
	Repare que $d\eta=0$ e $i^*\eta=0$, pois
	\[i^*\eta=i^*\omega-di^*\mu=d\alpha-d\alpha=0\]

	{\color{3}$\eta$ é uma 2-forma definida numa vizinhança $U$ de $S$, ela é fechada, e restrita a $S$ é zero. Então pode achar uma primitiva de $\eta$ numa vizinhança de $S$ que se anula quando restrita a $S$ pelo lema de Poincaré generalizado.}

	Pelo lema de Poincaré generalizado, existe uma 1-forma $\xi$ em $U$ tal que
	 \[\eta=d\xi\qquad \text{e}\qquad i^*\xi=0 \]
	 Defina então
	 \[\tau:=\mu+\xi.\]
	 Como
	 \[d\tau=d(\mu+\xi)=d\mu \omega-d\mu=\omega\]
	 e
	 \[i^*\tau=i^*\mu=\alpha,\]
	 concluimos a prova do lema.
\end{proof}

{\color{7}Como terminhamos a demostração do teorema ussando o lema?} Seja $X$ o campo em $U$ tal que $i_X\omega=\tau$. Vamos ver que esse $X$ satisfaiz a duas propriedades.

Primeira propriedade e imediata:
\[\mathcal{L}_X\omega=d(i_X\omega)=d\tau=\omega.\]
Agora vejamos que $X\pitchfork S$. Seja  $Y$ o campo de Reeb em $S$ i.e. dado pelas equações
\begin{align*}
	\alpha(Y)&\equiv 1,\qquad  i_Yd\alpha=0
\end{align*}
Note que
\[\omega(X,Y)=\alpha(Y)=1\]
Mas se $X$ fosse tangente a $S$ em algum ponto $x\in S$, teríamos que
\[\omega(X(x),Y(x))=d\alpha(X(x),Y(x))=0\]
{\color{3}A $\alpha$ restrita a $S$  é $\omega$.} Isso é uma contradição.
\end{proof}

\subsection{Exemplos de hiperfícies de contato}

Agora vamos ussar esse teorema para construir exemplos simples de variedades de contato.	

\begin{enumerate}
	\item O campo radial $X(x)=\frac{x}{2}$, onde $x=(q,\gamma)\in\mathbb{R}^{2}$ é claramente conformemente simplético com respeito à $\omega=\sum dq_i\wedge dp_i$ pois $i_X\omega=\frac{1}{2}\sum q_idp_i-p_idq_i$.

	{\color{3}Portanto, qualquer hiperfície compacta star-shaped (i.e. $\pitchfork X$) é tipo contato.}

	Em particular, as esferas vão ser de contato.

	\item Depois de $\mathbb{R}^{2n}$, a seguinte variedade de contato mais natural é o fibrado contangente $T^*Q$ com a forma simplética canónica. Seja $X(q,p)=p$ o campo radial  ão longo das fibras. Temos que
		\[i_X\omega=\sum p_idq_i=-\alpha\]
		one $\alpha$ é a forma tautológica. Em particular o fibrado unitário $S^*Q=\{(q,p)\in T^* Q:\|p\|=1\}$ com qualquer métrica Riemanniana em $Q$ é uma hiperfície tipo contato.

		{\color{3}Níveis de energia do fluxo geodésico são hipersuperfícies de contato. O mundo natural onde mora o fluxo geodésico é o fibrado unitário. Em Da Silva: o fluxo geodésico é um sistema Hamiltoniano.}
\end{enumerate}



\section{Aula 22}

Hoje vamos terminhar de falar de hiperfícies de contato. Interessantemente não é fácil construir exemplos de hiperfícies que não são de contato. Vamos dar um exemplo agora

\subsection{Exemplo de hiperfície que não é de contato}

\begin{example}[De hiperfície em $\mathbb{R}^{2n}$ que não é de contato]\leavevmode
	Pegue duas constantes positivas $R>r>0$ e considere a hiper ífice $S$ em $\mathbb{R}^{2n}$ difeomorfa a $S^{2n-1}$ que tangencia os bordos dos cilindros $Z(R)$ e $Z(r)$ onde
	 \[Z(r)=\{(p,q)\in\mathbb{R}^{2n}:q_1^2+p_1^2\leq r\}\]
	 i.e. pegue um disco de raio $r$ no plano \textit{simplético} (é importante que seja simplético) e varra ele por tudo o $\mathbb{C}^{n-1}$ restante. E também tem o cilindro de raio $R$. Ahí pode botar uma esfera que contenha o círculo de raio $R$, e deformar ela de um jeito simpático para fazé-la ser tangente também ao cilindro pequeno que esta dentro do cilindro grande.

	 A esfera normal sim é de contato. Mas essa não.

	 Seja $N$ um campo normal a $S$ que em $S\cap Z(R)$ aponta para fora e em $S\cap Z(r)$ aponta para dentro. (Qualquer campo normal a essa esfera tem essa propiedade, ou a inversa---dentro/fora, tanto faiz.) {\color{4}Parece que a esfera é deformada ainda mais, de tal jeito que ten regiões onde coincide com os cilindros---com os dois.} Note que tal campo pode ser dado por
	 \[N( q_1,\ldots,q_n,p_1,\ldots,\mathbb{P}^n)=(q_1,0,\ldots,p_2,0,\ldots,), \qquad \forall (q_1,p_1)\in Z(R)\]
	 pois é porque simplesmente é não fazer nada nas coordenadas de dimensão maior e, suponho, no círculo esse é o campo normal.

	 Note que $J_0N$ é tangente a $S$ ($J_0$ é a matriz simplética que usamos muito no principio do curso). Vamos a ussar a convenção (acho que coincide com a usamos antes) de que métrica compatível é quando
	 \[g=\omega(\cdot,J_0\cdot )\]
	 é um produto interno, i.e.
	 \[g(\cdot,J_0\cdot )=-\omega(\cdot,\cdot ).\]
	 Em fim, $J_0N$ é tangente a $S$ porque
	 \[\left<J_0N,N\right> =\omega(N,N)=0\]
 e como \[\omega_0(N,J_0v)=-\omega_0(J_0N,v)=g(N,v)=0\;\forall v\in TS\]
 i.e. $J_0N$ gera $\ker \omega|_{S}$.

 {\color{8} O ponto é que o núcleo dessa forma é gerado por $J_0N$.}

 {\color{8}Então como $J_0N$ ta no kernel, se tem uma forma de contato $\beta$, ela é não zero no kernel, assim $\beta J_0N$ e sempre positiva ou sempre negativa.}

 {\color{8}A gente argumentou que $J_0N$ ta no kernel. Como é que $\beta(J_0N)$ é sempre positivo/negativo? }

 Suponha que $S$ possui uma forma de contato $\beta$ de forma que $\beta(J_0N(r))>0$ ou $ \beta(J_0N(r))<0$ para toda $r \in S$. Suponha, sem perda de generalidade, que $ \beta(J_0N)>0$.

{\color{7}Ahora note que $J_0$ rota $N$ no sentido antihorario, i.e. vai ser tangente  ao círculo, e a curva integral (ele chama de órbita) vai ser o círculo. Mas como no círculo grande/pequeno, o campo normal muda de ir pra dentro e pra fora, aí o fluxo vai mudar de direção.}

Pela construção, $S$ possui duas órbitas periódicas $\gamma_R$ e $\gamma_r$ em $S\cap Z(R)$ e $S\cap Z(r)$ no plano (afim) $(q_1,p_1)$ quando em direções  opostas.

Seja $\alpha=\sum p_idq_i$ de forma que $d\alpha=-\omega_0$. É fácil ver que
\[\int_{\gamma_r}\alpha=-\pi r^2\qquad \text{ e }\qquad \int_{\gamma_R}\alpha=\pi R^2 \]
{\color{8}Aqui é Stokes acho}, é a área dos círculos.

Por outro lado, como $S\cong S^{2n-1}$, $H^{1}(S)=0$, então toda forma fechada é exata. A forma $\alpha+\beta$ é exata pois $d(\alpha+\beta)=-\omega_0+\omega_0=0$. Daí, existe $f \in\mathcal{C}^\infty(S)$ tal que $\alpha+\beta=df$. Consequentemente,
\[\int_{\gamma_r}\beta=-\int_{\gamma_r}\alpha>0 \qquad \text{ e }\qquad \int_{\gamma_R}\beta=-\int_{\gamma_R}\alpha<0 \]
mais esses campos são soluções  de $J_0N$, e tamos supondo que  $\beta(J_0N)>0$, que é absurdo. Portanto, $\beta$ não existe e $S$ não é de contato.

{\color{8}Mas qual é a contradição? A condição $\beta(J_0N)>0$ diz o que sobre essas integrais?}
Ah porque as integris são, por definição 
\[0<\int_{\gamma_r}\beta=\int_{0}^{T_r}\beta(\dot \gamma_r)=\int_{0}^{T_r}\beta(J_0N)\]
onde $T_r$  é o período de $\gamma_r$. A outra integral é análoga. Assim a gente vai ter que essas integrais são uma positiva e a outra negativa. Mas a gente viu que $\beta(J_0N)$ é sempre positiva ou negativa---as integrais deveriam ser do mesmo signo.
\end{example}

\subsection{Simplectização}

Seja $(N,\alpha)$ uma variedade de contato.

\begin{defn}\leavevmode
	A \textit{\textbf{simplectização}} de $(N,\alpha)$ é a variedade de contato
	\[(N\times \mathbb{R},d(e^t\alpha))\]
	onde $t$ é a coordenada em $\mathbb{R}$. {\color{4}Exerício simples:} verifique que essa forma é simplética.
\end{defn}

\begin{thm}\leavevmode
	{\color{4}Toda hiperfície de contato e localmente a simplectização.} Seja $S$ hiperfície de contato compacta de $(M,\omega)$. Então existe uma vizinhança $U$ de $S$, $\varepsilon>0$ e um simplectomorfismo
	\[\psi:(U,\omega)\longrightarrow (S\times \Big((-\varepsilon,\varepsilon),d(e^t\alpha)\Big)\]
\end{thm}

\begin{proof}\leavevmode
	Lista.
\end{proof}

\subsection{Sistemas Hamiltonianos}

$(M,\omega)$. Pegue um campo simplético, i.e.
\[\mathcal{L}_X\omega=0\qquad \iff\qquad d(i_X\omega)=0\]
A gente vai ter uma função $H:M\to \mathbb{R}$ e o campo hamiltoniano dela satisfazendo $i_{X_H}\omega=dH$.

\begin{remark}\leavevmode
	$\mathcal{L}_{X_H}H=dH(X_H)=\omega(X_H,X_H)=$
\end{remark}

\begin{coro}[Teorema de Liouville]\leavevmode
	{\color{3}(Campo hamiltoniano preserva volume.)}
	\[\mathcal{L}_{X_H}\Lambda=0,\qquad \text{onde }\Lambda=\underbrace{\omega\wedge \ldots \wedge \omega}_{n} \]
\end{coro}

\begin{remark}\leavevmode
	Teoría erdógica tem a ver com transformações que preservam medida (como os fluxoshamiltonianos acho).
\end{remark}

\begin{remark}[Nem todo campo simplético é Hamiltoniano] Aquí a gente precisa de uma variedade não simplesmente conexa. Por exemplo o cilindro $S^1\times \mathbb{R}$. Aqui $\omega=d\theta\wedge dh$. Pegue $X=\frac{\partial }{\partial h}$. Então $i_X\omega=d\theta$, i.e. é simplético. Mas não é Hamilatoniano.
\end{remark}

\subsection{Exemplos de campos hamiltonianos}

\begin{enumerate}
	\item \textbf{Função altura em $S^2$.} Pegue  $H=h$. Pegue um valor regular  $k\in \mathbb{R}$. Então a diferencial de  $H$, $dH$ e surjetiva nos conjuntos de nível e por isso o campo hamiltoniano não se anula. Assim a gente tem um fluxo nos círculos da esfera paralelos ao plano $xy$, e nos polos tem singularidades.

	\item \textbf{Levantamentos de campos.} Seja $Q$ uma variedade diferenciável e $X\in\mathfrak{X}(Q)$. {\color{8}Seja $\varphi_t$ o fluxo de $X$. Como ele é um difeomorfismo, a gente pode levantá-lo no cotangente como vimos na lista não sei o que. O levantamento preserva a 1-forma canónica.} Então seja \[\hat{\varphi}_t(x,p)=\Big(\varphi_t(x),(d\varphi_t(x)^*)^{-1}p\Big)\] o levantamento de $\varphi_t$ para $T^*Q$.

		Defina $\hat{X}=\frac{d}{dt}\hat{\varphi}t$. Sabemos que $\hat{X}$ preserva a forma tautológica, i.e. $\mathcal{L}_{\hat{X}}\alpha=0$ onde $\alpha$ é a forma tautológoica. Segue daí que $\hat{X}$ é o campo hamiltoniano de $H:=i_{\hat{X}}\alpha$. De fato:
		\[0=\mathcal{L}_{\hat{X}}\alpha=d(i_{\hat{X}}\alpha)+i_{\hat{X}}d\alpha\implies i_{\hat{X}}\omega=d(i_{\hat{X}}\alpha)\]
{\color{8}Começa com um campo que não tem nada, pode ser cualquer coisa random. Quando você levanta esse cara, ele vai preservar volume. É surpdendente.}
	
\item \textbf{Fibração de Hopf.} Pegue 
	\[H:\mathbb{R}^{2n}\longrightarrow \mathbb{R},\qquad H(x)=\frac{1}{2}\|x\|^2\]
	Então qual o campo hamiltoniano
	\[X_H(x)=-J_0\nabla H(x)=-J_0x\]
	\[\varphi_t(x)=e^{t J_0}x=(\cos t)x-(\sin t)J_0x\]

\item \textbf{?} Pegue $(Q,g)$ Riemanniana.  Pegue a energia cinética
\begin{align*}
	H: T^*Q &\longrightarrow \mathbb{R} \\
	(q,p) &\longmapsto \frac{1}{2}\|p\|^2
\end{align*}
onde (?)
 \[\|p\|=\operatorname{s u p }_{x\in B}\|p(x)\|\]

	dada pela energia cinética

\begin{exercise}\leavevmode
	Mostr que o fluxo Hamiltoniano de $H$ é o fluxo geodésico de $Q$ via a identificação
	\begin{align*}
		: TA &\overset{\cong }{\longrightarrow} T^*Q \\
		(q,v)&\longmapsto (q,g(v,\cdot ))
	\end{align*}
	\textbf{Hint:} ver o livro da Anna. E um trabalhinho ver isso mas vale a pena.
	
\end{exercise}
\end{enumerate}

\section{Aula 23: dinâmica Hamiltoniana}

$(M,\omega)$ variedade simplética.

\begin{defn}\leavevmode
	Seja $S\subset M$ uma hiperfície. O  \textit{\textbf{fibrado de linhas características}} $L_S$ ou \textit{\textbf{distribuição característica}} de $S$ é o subfibrado de $TS$ dado por $\ker(\omega|_{S}$.
\end{defn}

\begin{remark}\leavevmode
	Como $M$ é simplética, $S$ tem dimensão impar e por isso esse kernel tem dimensão (positiva).
\end{remark}

Note que se $S=H^{-1}(k)$ para algum $H:M\to\mathbb{R}$ e  $k$ valor regular, então
\[X_H(x)\in L_S(x)\qquad \forall x\in S\]

\begin{prop}\leavevmode
	Seja $S$ um nível de energia de $H_1$ e $H_2$. Então os fluxos de $H_1$ e $H_2$ em $S$ diferem apenas por uma reparametrização.
\end{prop}

\begin{defn}\leavevmode
	Uma \textit{\textbf{característica fechada}} em $S$ é uma curva regular (derivada não se anula) $\gamma:S^1\to S$ tal que 
	\[ \dot(t)=L_S(\gamma(t)),\qquad \forall t\]
\end{defn}

Portanto a existência de órbitas periódicas é equivalente à existência de características fechadas.

\begin{defn}\leavevmode
	Sejam $S$ e $S'$ variedades diferenciáveis e  $X\in\mathfrak{X}(S)$, $Y\in\mathfrak{X}(S')$. Dizemos que $X$ e $Y$ são topologicamente (resp. diferenciavelmente) equivalentes se existe um homeomorfismo (difeomorfismo) $h:S\to S'$ {\color{5}que leva curvas integrais de uma variedade em outra preservando o tempo?} tal que dado $p\in S$ e $\varepsilon>0$ tais que para $0<t<\delta$, $h_{\varphi_t}(p)=\psi_{\bar{t}}(h(p))$ para algum $0<\bar{t}<\varepsilon$, onde $\varphi_t$ e $\psi_t$ são os fluxos de $X$ e $Y$.

	Quando $\bar{t} =t$ dizemos que são \textit{\textbf{conjugados}}. No caso de difeomorfismos, a  condição acima significa $dh(X)=fY$ para alguma $f$, e são conjugados se $f=1$ (i.e. $h$-relacionados).
\end{defn}

\begin{upshot}\leavevmode
	Tem um Hamiltoniano super bom que pode espalhar numa vizinhança  da subvariedade.
\end{upshot}

\begin{prop}\leavevmode
	Seja $S$ uma hiperfície de contato de  $M$. Existe Hamiltoniana $H$ definida em uma vizinhança $U$ de $S$ tal que $S=H^{-1}(0)$ e $X_H|_{H^{-1}(0)}$ é diferenciablemente conjugado a $X_{H}|_{H^{-1}(k)}$ para todo $k$ suficientmente pequeno.
\end{prop}

\begin{proof}\leavevmode
	Vimos que existe uma vizinhança $U$ de $S$, $\varepsilon>0$ e um simplectomorfismo $f:(U,\omega)\to(S\times (-\varepsilon,\varepsilon), d(e^t\alpha)$ onde $\alpha$ é a forma de contato em $S$. {\color{5}Localmente, a forma de contato é simplectomorfa à sua simplectização, isso foi um exercício a aula,, passada?}

	Tome $H$ tal que $H\circ f^{-1}(x,t)=-e^t-1$. Note que $f_*X_H=X_{-e^t}$. {\color{4}Isso é imediato pelo requerimento de que $H\circ f^{-1}=-e^t-t$, pois o pushforward é Hamiltoniano da composição.}

	Então agora note que $X_{-e^t}=R_{\alpha}$ é o campo de Reeb. De fato:
	\begin{align*}
		i_{R_\alpha}(d(e^t\alpha)&=i_{R_\alpha}(e^tdt\wedge \alpha+e^t d\alpha)\\
&=-e^tdt\wedge \cancelto{1}{\alpha(R\alpha)}\\
&=d(-e^t)
	\end{align*}
\end{proof}

Ahora fixe uma subvariedade $S$ que seja o nível de energia de um Hamiltoniano i.e., $ S=H^{-1}(k)$ para $k$ valor regular. Defina $\varphi_t$ o fluxo Hamiltoniano de $H$.

\begin{prop}\leavevmode
	Então existe uma forma de volume $\Lambda_S$ em $S$ preservada por $\varphi_t$.
\end{prop}
\begin{remark}\leavevmode
	Já sabemos que tem essa forma na variedade simplética: lista não sei o que. Mas não pode só restringir essa à $S$ : porque isso vai dar uma $(2n)$-forma! A gente precisa de uma $2n-1$ forma.
\end{remark}

\begin{proof}\leavevmode
	Como $dH(x)\neq 0$ para todo $x$ em uma vizinhança $U$ de $S$, existe $(2n-1)$-forma $\eta$ em $U$ tal que $\Lambda=dH\wedge \eta$. Seja $i:S\hookrightarrow M$ e a inclusão é
	\[\Lambda_S:=i^*\eta.\]
Vamos ver que essa é boa.	Claramente, $\Lambda_S$ é uma forma de volume em $S$. 
\begin{claim}\leavevmode
	$\Lambda_S$ está únicamente caracterizada pela equação
	\[\Lambda=dH\wedge \eta,\]
	\textbf{ou seja} , se $\xi$ satisfaiz $\Lambda=dH\wedge \xi$, então
	\[\Lambda_S=i^*\eta=i^*\xi.\]
\end{claim}
\begin{proof}[Prova da afirmação]\leavevmode
	De fato, se $dH\wedge (\eta-\xi)=0$ então $\eta- \xi=dH\wedge k$ para alguma $(2n-2)$-forma $\kappa$ em $U$. Como $i^*dH=0$, concluiriamos que $i^*\eta=i^*\xi+i^* (\cancelto{0}{dH}\wedge \kappa=i^*\xi$.
\end{proof}
Como $\varphi^*_t\Lambda=\Lambda$ e $\varphi^*_t=dH$ temos $\varphi^* _t\Lambda_S=\Lambda_S$. De fato, \[\Lambda=\varphi^*_tdH=dH\wedge \varphi^*_t\eta=dH\wedge \varphi_* t\eta\]
portanto, pela unicidade no claim, 
\[i^*\varphi^*_t=i^*\eta\implies \]
\end{proof}

\begin{defn}\leavevmode
	Dizemos que $x\in S$ é \textit{\textbf{recorrente}} para  $\varphi_t$ se existe uma sequência $t_k$ tal que
	 \[\lim_{k \to \pm \infty} t_k=\pm \infty\]
	 (isso é só para realmente pensar no tempo infinito… sem isso a gente vai ficar parada eventualmente,, ou algo assim…)	 e tal que
	 \[\lim_{k \to \pm \infty} \varphi_{t_k}(x)=x\]
	 {\color{6}que a órbita passa arbitrariamente perto do $x$.}
\end{defn}

\begin{thm}[de Recorrência de Poincaré]\leavevmode
	Seja $\varphi_t$ um fluxo que preserva volume em uma variedade comapcta $S$. Então quase todo ponto de  $S$ é recorrente para $\varphi_t$.
	
\end{thm}
\section{Aula 24}

\subsection{Prova do teorema de Reoccrência de Poincaré}

\begin{proof}[Prova do teorema de recorrência]\leavevmode
Usaremos a medida inducida pela forma de volume, i.e. $\mu(A)=\int_{A}\mu_S$.

Seja $\varphi=\varphi_1$. Vamos provar inicalmente que
\begin{equation}\label{eq:medida1}\mu\left(A\cap\left( \underbrace{\bigcap_{k\geq 0} \bigcup_{j\geq k} \varphi^{-j}(S)}_{:=B} \right) \right)=\mu(A)
\end{equation}
para todo borealiano $A$. Observe que $x\in B$ se para todo inteiro $k \geq 0$ existe $j\geq k$ tal que $x\in \varphi^{-j}(A)$, ou seja, $\varphi^j(x)\in A$.

 {\color{4}ou seja, para tempos arbitrários o $\varphi(x)$ já tá em $A$.}

 Para provar \cref{eq:medida1}, defina
 \[A_k=\bigcup_{j\geq k} \varphi^{-j}(A).\]
 Note que 
 \[A_0\supset A_1\supset A_2\supset \ldots\]
 Ademais, como $\varphi(A_k)=\bigcup_{j \geq k}\varphi^{k-j}(A)$, de modo que $\varphi^k(A_k)=A_k$ para toda $k$. Além disso, $A_k\subset A_0$.

 Isso significa, \textit{como a medida é finita}, que $A_k=A_0$ em quase todo ponto.

 Portanto, $\bigcap_{k\geq 0} A_k=A_0$ q.t.p. Como $A\subset A_0$, concluimos que
 \[A\cap\left( \bigcap_{k\geq 0} A_k \right) =A\cap A_0=A \qquad \text{q.t.p.} \]
Isso implica \cref{eq:medida1}.

Agora, use o fato de que existe uma base enumerável de $S$. Portanto, para todo $n\in \mathbb{N}$ existe um conjunto enumerável de bolas abertas $B_j(\frac{1}{n})$ de raio $\frac{1}{n}$ cobrindo $S$. Pela primeira parte do argumento, existe um conjunto de medida nula $N_n \subset S$ tal que todo $x\in S\setminus N_n$ retorna infinitas vezes para a bola que o contém.

Tomando $N=\bigcup_{n} N_n$, temos que $\mu(N)=0$ e todo $x\in S\setminus N$ possui a propriedade de voltar infinitas vezes para  a bola  $B_j(\frac{1}{n})$, seja qual for o valor de $n$.

Isso prova o teorema para tempos futuros. Para temos passados e ó mesmo ussando $\varphi^{-1}$.
\end{proof}

Esse teorema é muito forte: {\color{2}quase todo ponto é recorrente}. Esse teorema é valido em espaços mais gerales que variedades. Na verdade, se definir recorrência só para tempos futuros, pode ussar um simlpes automorfismo que preserva medida num espaço com uma base enumerável.

\subsection{A seguinte questão: órbitas periódicas}

\begin{question}\leavevmode
Dado um nível de energia compaco $S$, existem órbitas periódicas no mesmo?
\end{question}

Este problema possui uma estrutura variacional. Suponha, por simplicidade, que $\omega=-d\alpha$. Dada uma Hamiltoniana $H:M\to\mathbb{R}$ e uma curva fechada $\gamma:[0,T]\to M$ defina sua \textit{\textbf{ação}} como
\[A_H(\gamma)=\int_{\gamma}\alpha-H\]
e esse se chama \textit{\textbf{funcional de ação}}.

\begin{prop}[Princípio de Hamilton]\leavevmode
	Pontos críticos do funcional de ação (no espaço de curvas fechadas de período T fixado) estão em bijeção com as órbitas pers de $H$ de príodo $T$. Mais precisamente, uma curva fechada $\gamma:[0,T]\to M$ é uma órbita de per $X_H$ se e somente se dada uma variação suave $\Gamma:[-\varepsilon,\varepsilon]\times [0,T]\to M$ tal que $\Gamma_0=\gamma$, tem-se
	\[\frac{d A_0(\Gamma_s)}{ds}\Big|_{s=0}=0\]
{\color{5}O $\Gamma_s$ é uma curvinha pertinho de $\gamma$. A maneira de você pensar isso é que $\gamma$ é um ponto no espaço $C^\infty(S^1,M)$ e $\Gamma$ é uma curvinha que passa por $\gamma$ em $s=0$. Deriva o funcional de ação em $\Gamma_s$ e tem um ponto crítico em $s=0$, i.e. em  $\gamma$.}
\end{prop}

\begin{proof}\leavevmode
Seja $X$ o campo vetorial ao longo de $\gamma$ dado por $\frac{d}{ds}\Big|_{s=0}\Gamma_s$. {\color{5}Isso é como a derivada de curva no espaço de lazos, um campo vetorial que poemos no $\gamma$.} Por Stokes
\begin{align*}
	\underbrace{\int_{[0,\varepsilon]\times [0,1]}\Gamma ^* \alpha}_{(*)}&=\int_{\varepsilon\times [0,T]}\Gamma^*\alpha-\int_{0 \times [0,T]}\Gamma ^*\alpha\\
	&=\int_{\Gamma_\varepsilon}\alpha-\int_{\gamma}\alpha
\end{align*}
{\color{5}porque é uma curva fechada no rectángulo, as bordas verticais se cancelam.}

Dividindo ambos lados por $\varepsilon$ e fazendo $\varepsilon \to 0$ temos que o lado direito converge para a derivada de $\int_{\gamma}\alpha$ na direção de $X$ e o lado esquerdo tende a $\int_{\gamma}d \alpha(X,\dot{\gamma})$:
\begin{align*}
	(*)&=\int_{0}^\varepsilon\underbrace{\int_{0}^T \Gamma^*d\alpha_{(s,t)}\left( \frac{\partial }{\partial s},\frac{\partial }{\partial t} \right) dt}_{:=f(s)}ds\\
	&=F(\varepsilon)-F(0),\qquad F:=\int fds
\end{align*}

Mas,
\begin{align*}
\lim_{\varepsilon \to 0} \frac{F(\varepsilon)-F(0)}{\varepsilon}&=F'(0)=f(0)\\
&=\int_{0}^T \Gamma ^*d\alpha_{(0,t)}\left( \frac{\partial }{\partial s},\frac{\partial }{\partial t} \right) dt\\
&=\int_{0}^T d \alpha_{\Gamma(0,t)}\left( \frac{\partial }{\partial s}\Gamma,\frac{\partial }{\partial t}\Gamma \right) dt\\
&=\int_{0}^Td \alpha_{\gamma(t)}(X,\dot{\gamma})dt
\end{align*}
$\mathsf{OK}$. Agora vamos botar o funcional de ação:
\begin{align*}
	\frac{d A_H(\Gamma_A)}{ds}\Big|_{s=0}&=\int_{0}^T-d \alpha(\dot{\gamma},X)-dH(\gamma)Xdt\\
					     &=\int_{0}^T \omega(\dot{\gamma},X)dt\\
					     &=\int_{0}^T\omega( \dot{\gamma}-X_H,X)dt
\end{align*}
que é zero see $\dot{\gamma}=X_H$.
\end{proof}

\begin{remark}\leavevmode
	O argumento anterior se aplica para a prova o seguinte: se $\gamma:[0,T]\to M$ é uma curva não necessariamente fechada e $L_1,L_2\subset M$ são subvariedades tais que $\alpha|_{L_i}=0$ para $i=1,2$ e  $\gamma(0)\in L_1$, $\gamma(1)\in L_2$. Entao $\gamma$ é uma órbita de $X_H$ sse é um ponto crítico de $A_H$ para qualquer variação 
	\[\Gamma_0=\gamma,\qquad \Gamma_s(0)\in L_1\qquad \text{ e} \qquad \Gamma_s(1)\in L_2\qquad \forall s\]

	\begin{example}\leavevmode
		$M=T^*Q$, $L_1=T^*_xQ$, $L_2=T^*_yQ$.
	\end{example}
\end{remark}

\subsection{Extra: a note from Graham's talk just today's evening}

In a chart, $U\subseteq \mathbb{R}^{n}$, $g_{ij}:U \to \operatorname{Sym}^+(\mathbb{R}^{+})$, define the \textit{\textbf{action functional}} as
\[A[\gamma]=\int_{0}^1g_{ij}(\gamma(t))\dot{\gamma}^i \dot{\gamma}^jdt\]
\begin{enumerate}[label=(\roman*)]
\item Cricital points of $A$ are geodesics.
\item This is a standard Lagrangian action (with $V=0$).
	 \begin{align*}
		\mathcal{L}: TU &\overset{g}{\longrightarrow} T^*U \\
		(x,v)&\longmapsto (x,g_{ij}v^j)
	\end{align*}
	So the moral is: Legendre transform is a musical isomorphism.
\end{enumerate}

\section{Aula 25}

\begin{question}\leavevmode
	Dado um nível de energia regular compacto $S$ de uma Ham. $H$ existem órbitas periódicas em $S$?
\end{question}

\begin{thing4}{Resposta}\leavevmode
	Em geral não: existem níveis de energia compactos sem órbitas periódicas.
\end{thing4}

\begin{thm}[Rabinowitz, '78]\leavevmode
	Se $S\subset\mathbb{R}^{2n}$ é star-shaped (o campo radial é transversal a ele) então possui uma orbita periodica.
\end{thm}

\begin{thing4}{Conjetura}[Weinstein $\sim$'79]\leavevmode
Toda hiperfície de contato $S$ possui uma órbita periódica.	
\end{thing4}

\begin{thing3}{Respostas parciais}\leavevmode
\begin{itemize}
\item $S\subset M^4$ ($S$ dim 3) é veradde (Taubes).
\item $\dim S>3$ resultados parciais positivos. 
\end{itemize}
\end{thing3}

\begin{remark}\leavevmode
	É muito difícil demostrar se numa variedade de contato aribtrárie o fluxo de Reeb tem ou não tem uma órbita periodica.
\end{remark}

\begin{thing4}{Exemplos}\leavevmode
\begin{enumerate}
\item \textbf{(Dinâmica em elipsóides)}.\hspace{1em} Pega a Hamiltoniana
	\begin{align*}
		H: \mathbb{C}^n &\longrightarrow \mathbb{R} \\
		(z_1,\ldots,z_n) &\longmapsto \sum_{j=1}^n\frac{|z_j|^2}{r_j}
	\end{align*}
onde $0<r_1\leq r_2\leq \ldots\leq r_n$. Como $H$ é homogênea (as variaveis "saem" do argumento da função elevadas ao grau de homogênidade) de grau 2, as dinâmicas em níveis de energia regulares de $H$ são equivalentes {\color{3}(Exercício) geometricamente é claro: a função $\varphi:\mathbb{C}^{n}\to \mathbb{C}^{n}:z\mapsto cz$ 'e conformemente simplética, ela manda o fibrado caracterísico $L_S$ em  $L_{cS}$ (de um elipsoide pequeno a um elipsoide grande)}.

Portanto, basta considerar $S=H^{-1}(1)$. Como
\[X^{(z)}_H=J\nabla H(z)=J\left( \frac{2}{r_1}z_1,\ldots,\frac{2}{r_n}z_n \right),\]
o fluxo de $X_H$ é
\[\varphi_t(z_1,\ldots,z_n)=\left( e^{t\frac{i}{r_1}}z_1,\ldots,e^{t\frac{i}{r_n}}z_n\right) \]
Consequentemente, temos ao menos $n$ orbitas periodicas em $S$ dadas por
\[\gamma_j(t)=(0,\ldots,0,z_j(t),0,\ldots,0)\]
onde $z_j(t)=e^{t\frac{i}{r_j}}z$. De fato, existem duas possíveis situações:
\begin{itemize}
\item $r_1,\ldots,r_n$ são \textit{\textbf{racionalmente independentes}}, ou seja,
	\[\left<(r_1,\ldots,r_n,v\right>\neq 0\qquad \forall v \in\mathbb{Z}^n\setminus\{0\}\]
	(geometricamente signfica que a linha determinada pelo vetor não toca a lattice $\mathbb{Z}^n$). Neste caso, não existem outras órbias periódicas.	
\item $r_1,\ldots,r_n$ são racionalmente dependentes. Neste caso, existem infinitas órbitas periódicas. 
\end{itemize}
\begin{exercise}\leavevmode
	Prove estas afirmações.
\end{exercise}
\item \textbf{(Fluxos geodésicos)}.\hspace{1em} $(Q,g)$ variedade Riemanniana.
	\begin{align*}
		f: TQ &\longrightarrow T^*Q \\
		(x,v) &\longmapsto (x,i_{v}g)
	\end{align*}
	isomorfismo de fibrados. Então puxamos a forma simplética canônica, $\omega=f^* \omega_0$. Considre
	\begin{align*}
		H: TQ &\longrightarrow \mathbb{R} \\
		H(x,v) &=\frac{1}{2}g(u,v)
	\end{align*}
O fluxo Hamiltoniano $\varphi_t$ de $H$ é o fluxo geodésico de $(Q,g)$, ou seja,
\[\varphi_t(x,v)=\left( \gamma_{(x,v)}(t),\dot{\gamma}_{(x,v)}(t) \right) \]
onde $\gamma_{(x,v)}$ é a única geodésica tal que $\gamma_{(x,v)}(0)=x$ é $\dot{\gamma}_{(x,v)}(0)=v$.
Como $H$ é homogênea (de grau 2) as dinâmicas em $H^{-1}(c)$, $c\neq 0$, são dif. equivalentes. Portanto basta considerar o caso $S=H^{-1}(1/2)=SQ$ (o fibrado esférico ou unitário). É um fato conhecido que se $Q$ for compacta então existe uma geodésica fechada. O caso em que $\pi_{1}(Q)\neq 0$ é fácil de provar e devido a Cartan. O caso $\pi_{1}(Q)=0$ é devido a Lustenik-Fet.

\begin{thing4}{What is going on?}\leavevmode
	Nesses casos tem órbita períodica
\end{thing4}

\begin{enumerate}
\item (\textbf{Fluxo geodésico em $S^n$ curvatura seccional $K\equiv 1$})\hspace{1em} Neste caso, as órbitas periódicas são todas fechadas (dadas pelos grandes círculos) e de mesmo período mínimo.

	\begin{thing4}{Comentario}\leavevmode
		A redução seguinta nos diz que a Grassmaniana orientada $\operatorname{Gr}^+(2,n+1)$ é uma variedade simplética:
		\[\begin{tikzcd}
			S^1\arrow[r]& S S^n\arrow[d]\\
			&\operatorname{Gr}^+(2,n+1)
		\end{tikzcd}\]
	\end{thing4}
\item (\textbf{Curvatura secional $K\equiv 0$, todo}).\hspace{1em} Neste caso, $ST^2=\bigcup_{v\in S^1} v\times T^n$. Esses toros são invariantes pelo fluxo global, e em cada termo $v \times T^2$ temos um campo constante igual a $v$. É fácil ver que se $v=(v_1,\ldots,v_n)$ é tal que $v_1,\ldots,v_n$ são racionalmente independentes então todas as órbitas em $v\times T^n$ são densas. Se $v_1,\ldots, v_n$ são racionalmente dependentes, todas as órbitas são fechadas (em $v \times T^n$). {\color{3}(Exercício).}

\item (\textbf{Caso hierbólico, i.e. $K<-1$} \hspace{1em} Neste caso,pode-se mostrar que existe uma decomposição \[TS=\operatorname{s pan}\{x\}\oplus E^s \oplus  E^u\] onde $\dot X=\dot{\varphi}_t=X_H$, tal que  $d\varphi_t(E^{s,u}=E^{s,u}$ e existem constantes 
	\[\|d\varphi_t(X)\leq a^t\|v\|,\qquad \forall v\in E^1,\qquad \forall t>0,\]
	\[\|d\varphi_t(v)\|\geq b^t\|v\|,\qquad \forall v\in E^u,\qquad \forall t>0.\]
Fluxos com esta propriedade se chaman  \textit{\textbf{Anosov}}.

Utilizando teoría hiperbólica de sistemas dinâmicos, pode se provar que o conjugado das órbitas periódicas de $\varphi_t$ é denso e, ao mesmo tempo, existem órbitas periódicas.
\end{enumerate}
\end{enumerate}
\end{thing4}

\section{Aula 26}

Hoje vamos terminhar a parte de dinâmica Hamiltoniana falando de uma classe de sistemas Hamiltonianos muito interessantes que se chaman circuitos magnéticos.

\subsection{Fluxos magnéticos}

Pegue uma variedade com métrica Riemanniana $(Q,g,\Omega)$ onde $\Omega$ é o \textit{\textbf{campo magnético}}, uma 2-forma fechada {\color{3}(a condição de ser fechada tem a ver com as equações de Maxwell)}. De novo, como $M$ é Riemanniana podemos pegar o pullback da forma canônica no fibrado cotangente, obtemos a forma $\omega=f^*\omega_{\operatorname{can}}$.

\begin{exercise}\leavevmode
	Dada a projeção $\pi:TQ\to Q$, verifique que $\omega+\pi^*Q$ é não degenerada. Chamamos ela de \textit{\textbf{forma simplética twist}}
\end{exercise}

Agora pegue o mesmo Hamiltoniano que na aula passada, a energia cinética
\begin{align*}
	H: TQ &\longrightarrow \mathbb{R} \\
	(x,v) &\longmapsto \frac{1}{2}g(v,v)
\end{align*}
Se pegamos o fluxo Hamiltoniano respeito a $\omega$, a gente já viu, é o fluxo geodésico. Mas…
\begin{defn}\leavevmode
O \textit{\textbf{fluxo magnético}} de $(g,\Omega)$ é o fluxo Hamiltoniano de $H$ com respeito a $\omega+\pi^*\Omega$.	
\end{defn}

\begin{exercise}\leavevmode
	Seja $Y:TQ\to TQ$ a aplicação de fibrados definida pela equação
	\[\Omega(v,w)=g(Y(v),w),\qquad \forall v,w,\]
	ou seja, $Y(v)$ é o vetor que representa a forma $\Omega(v,\cdot)$ (mediante o isomorfismo musical da métrica, Riesz). Se chama \textit{\textbf{força de Lorentz}} de $\Omega$.

	Mostre que $\gamma$ é uma órbita do fluxo magnético se e somente se $\gamma=(c,\dot c)$ onde $c$ satisfaz a segunda equação de Newton
	\[\nabla_{\dot c}\dot c=Y(\dot c\]
	{\color{3}ou seja, sem termo magnético é simplesmente uma geodésica porque $Y=0$ assi que  não temos aceleração, aqui temos aceleração que é a força de Lorenz.}	
\end{exercise}

\begin{thing4}{Algumas diferenças em relação ao fluxo geodésico}\leavevmode
	\begin{itemize}
	\item Não é homogêneo. (No caso do fluxo geodésico, multiplicar um vetor por um número real resulta na mesma geodésica, só que reparametrizada). Aqui a dinâmica varia com o nível de energia.
	\item Existem níveis energia que não são de contato. Isso é porque quando torcemos a forma simplética já não temos a propriedade de ser conformemente simplético.
	\end{itemize}
\end{thing4}

\begin{thing4}{Exemplos}\leavevmode
\begin{enumerate}
\item O caso mais simples: uma  superfície e que o fluxo magnético (que não seja zero poruqe isso seria o fluxo geodésico) seja o pullback da forma de volume. Pegue $(S^2,K\equiv 1, \Omega)$ com $\Omega$ a forma de área.

Como $\Omega$ é a forma de área, $Y=J$  é a estrutura complexa de $S^2$ (lembre que $Y$ é um automorfismo do fibrado tangente, $Y(v)$  é o vetor que representa $\Omega(v,\cdot)$). Portanto, as trajetórias magnéticas são dadas por 
\[\nabla_{\dot c}\dot c=J(\dot c)\]
que são curvas com curvatura geodésica constante porque
\begin{align*}
\dot c\left<\dot c, \dot c\right>&=2 \left<\nabla_{\dot c}\dot c, \dot c\right>=2\left<Y(\dot c),\dot c\right>=2 \Omega(\dot c, \dot c)=0.
\end{align*}
(a derivada da norma da curva é zero), é isso é simplesmente porque $J$ preserva a norma.

Por exmeplo, pegue um paralelo na esfera que seja um círculo de raio $a$ paralelo ao plano $xy$, chama ela de  $\psi$. A curvatura geodésica dele é 
\[\left\| \nabla_{\dot \psi}\dot \psi\right\|=\frac{\sqrt{r^2-a^2} }{ra}\]
Ou seja, ela tem curvatura geodésica constante. Agora considere a ação de $\mathsf{SO}(3)$, que é o grupo de isometrias de $S^2$. Bom $\mathsf{SO}(3)$ age transitivamente também no fibrado unitário, e isso faiz que todas a curvas de curvatura geodésica constante são imagens de $\psi$ por elementos de $\mathsf{SO}(3)$.

Portanto, as trajetorias magnéticas de cada nível de energia são fechadas com o mesmo período mínimo.
\item $(T^2,K\equiv 0,\Omega=\text{área})$.
Todo é como no exemplo anterior, temos a estrutura complexa assim que as curvas são curvas de curvatura geodésica constante. Aqui no toro essas curvas podem ser entendidas no recobrimento universal, onde são círculos no plano. Elas são contrátileis. Isso não acontence no caso geodésico, pois a geodésicas nunca são contratileis.

Agora lembre que o fibrado unitário do toro é trivial
\[\begin{tikzcd}
ST^2=S^1 \times T^2\arrow[d]\\
T^2
\end{tikzcd}\]
é isso mostra que não temos uma superfície de contato, já que a forma de volume não é exata, assim que o seu pullback não pode ser exata---as superfícies de contato precisam formas exatas.

\item $(M^2,K\equiv-1,\Omega=\text{área} )$. Tudo é igual mas o que são as curvas de curvatura geodésica constante no recobrimento universal, que é o disco de Poincaré ou o semiplano superior?

Neste caso, pode-se provar o seguinte:
\begin{itemize}
\item Se $H<1/2$ então todas as trajetórias  são fechadas com projeção contrátil em $M$.
\item Se $H=1/2$ todas as trajetóricas são densas. Trata-se de horociclos: órbitas que tangencian a borda do disco, e perto desse ponto passam por infinitos dominios fundamentais (triangulinhos), tantos que quando projetamos obtemos uma curva densa.
\item Se $H>1/2$ o fluxo  é dif equivalente ao fluxo geodésico, em particular é Anosov (caótico).
\end{itemize}
\begin{remark}\leavevmode
	Isso se reflite no caso de $M$ e $TM$, onde  $H^{-1}(1/2)$ não é de contato, mas $H^{-1}(K)$ sim para $K>1/2$ ou $K<1/2$.
\end{remark}
\end{enumerate}
\end{thing4}

\subsection{Introdução à topologia simplética: invariantes globais}

\begin{question}\leavevmode
Como distinguir estruturas simpléticas? Quais invariantes existem?
\end{question}

Um invariante obvio é a dimensão. Esse é o único invariante local! Pelo teorema de Darboux. Um invariante global é o volume. Mas, será que isso é suficiente? Pode ter duas variedades difeomorfas com o mesmo volume que não são simplectomorfas?

\begin{question}\leavevmode
	Existem outros invariantes?
\end{question}


\section{Aula 27: invariantes globais}

\subsection{Teorema de non-squeezing de Gromov}


\begin{thm}[Non-squeezing de Gromov]\leavevmode
	Sejam $B^{2n}(r)=\{(x,y)\in\mathbb{R}^{2n}:\|x\|^2+\|y\|^2<r^2\}$ a bola de raio $r$ e $Z^{2n}(R)=\{(x,y)\in\mathbb{R}^{2n}:x_1^2+y^2<R^2\}$ o cilindro de raio $R$ sobre o $(x_1,y_1)$ ambos com a forma simplética $\sum_{i} dx_i\wedge dy_i$.

	Existe um mergulho simplético $B^{2n}(n)\hookrightarrow  Z^{2n}(R)$ sse $r\leq R$.
\end{thm}

\begin{remark}\leavevmode
É fundamental tomarmos um plano \textit{simplético} $( x_1,y_1)$; se tomassemos, por exemplo o plano isotrópico $(x_1,y_1)$ então a transformacão simplética $(x,y)\mapsto \left(\delta x, \frac{y}{\delta}\right)$ leva $B^{2n}(1)$ em $Z^{2n}(\delta)$ para qualquer $\delta>0$.
\end{remark}

\begin{defn}\leavevmode
	A \textit{\textbf{capacidade de Gromov}} de uma variedade simplética $(M^{2n},\omega)$ é
	\[c_G(M,\omega)=\operatorname{sup}\{\pi r^2:\exists \text{ mergulho simplético } B^{2n}(r)\hookrightarrow  M\}\]
\end{defn}

\begin{thm}[Non-squeezing de Gromov]\leavevmode
A capacidade de Gromov do cilindro é $\pi R^2$.
\end{thm}

\begin{defn}[Høfer-Zehnder]\leavevmode
	Considere a classe de variedades simpléticas $(M^{2n},\omega)$ c/n finado. Uma \textit{\textbf{capacidade simplética}} é um mapa $(M,\omega)\mapsto c(M,\omega)\in [0,\infty]$ tal que
\begin{itemize}
\item \textbf{Monotonocidade}: se $\exists $ um mergulho simplético $(M\omega)\hookrightarrow (N,\tau)$ então
	 \[c(M,\omega)\leq c(N,\tau)\]
\item \textbf{Conformalidade}: $c(M,a\omega)=|a|c(M,\omega)$ para toda $a \in \mathbb{R}\setminus\{0\}$.
\item \textbf{Não trivialidade}: $c(B^{2n}(r),\omega)=c(Z^{2n}(r),\omega)=\pi r^2$.
\end{itemize}
\end{defn}

\begin{remark}\leavevmode
Monotonocidade implica que a capacidade é um invariante simplético.	
\end{remark}

A capacidade de Gromov $c_G$ é uma capacidade simplética. Monotonicidade é obvia: so meter uma bola dentro de outra. A conformalidade segue do fato que $(B^{2n}(\sqrt{|a|} r),\omega)$ é simplectomofra a $(B^{2n}(r),a\omega)$ (exercício). E claro, a não trivialidade e o teorema de non-squeezing.

	A capacidade de Gromov é a menor capacidade simplética i.e. $c_G\leq c$ para qualquer capacidade $c$. Isso segue da monotonicidade e a não-trivialidade.

\begin{defn}\leavevmode
A \textit{\textbf{capacidade cilíndrica}} de $(M,\omega)$ é \[\bar{c}(M,\omega)=\operatorname{inf}\{\pi r^2:\;\exists \; \text{mergulho simplético } (M,\omega)\hookrightarrow  Z^{2n}(r)\}\]
\end{defn}
Note que $\bar{c}$ é a maior capacidade simplética. Portanto, dada qualquer capacidade simplética $c$ temos
 \[c_G(M,\omega)\leq c(M,\omega)\leq \bar{c}(M,\omega)\]

\begin{remark}\leavevmode
	Existem exemplos $(M,\omega)$ para as quais a desigualdade é estrita.
\end{remark}

\begin{example}[Capacidade de elipsóides]\leavevmode
	Considere o elipsóide
	\[E=\left\{ v \in\mathbb{R}^{2n}:\sum_{j,k=1}^{2n}a_{jk}v_jv_k\leq 1 \right\} \]
com $a_{jk}>0$.
\end{example}

\begin{exercise}\leavevmode
	Mostre que existe um simplectomorfismo levando  $E$ em
	\[E(r_1,\ldots,r_n)=\left\{ (z_1,\ldots,z_n)\in \mathbb{C}^n : \sum_{j=1}^n\left| \frac{z_j^2}{r_j} \right| \leq 1 \right\} \]
para algum $0 \leq  r_1 \leq  r_2 \leq \ldots\leq r_n$ e os números $r_i$ são únicamente determinados por $E$.
\end{exercise}

Como $B^{2n}(r_1)\subset E(r_1,\ldots,r_n)\subset Z^{2n}(r_1)$. Consequentemente, $c(E)=\pi r_1^2$. Note que enquanto o volume depende de todas as $r_i$, a capacidade depende só da menor! Então teorema de non-squeezing the Gromov para elipsoide diz que existe $r$ tal que $B^{2n}(r)\hookrightarrow  E$ sse $r \leq r_1$.

\subsection{Breve esboço da prova do teorema non-squeezing}

\begin{defn}\leavevmode
Seja $(M,J)$ uma variedade quase-complexa e $(\Sigma, j)$ uma superfície de Riemann ($j$ é complexa enquanto $J$ é só quase-complexa). Uma aplicação suave $u:\Sigma\to M$ tal que
\[du\circ j=J\circ du\]
é uma \textit{\textbf{curva $J$-holomorfa}} ou \textit{\textbf{pseudoholomorfa}}.
\end{defn}

\begin{remark}\leavevmode
	Cauchy-Riemann é equivalente a que uma função seja $\mathbb{C}$-linear, i.e.
	\[\frac{\partial f}{\partial x}-i\frac{\partial f}{\partial y}=0\iff df(ix)=idf(x).\]
\end{remark}

Seja agora $\varphi:B^{2n}(R)\longrightarrow Z^{2n}(R)$um mergulho simplético e  $r'<r$. Como $\varphi\left( \overline{B^{2n}(r')} \right) $ é compacto, existe um mergulho simplético
\[\varphi:B^{2n}(r')\longrightarrow\left( B^2(R)\times T^{2n-2},\omega_0 \oplus  \omega_0 \right) \]
onde $T^{2n-2}=\mathbb{R}^{2n-2}/a\mathbb{Z}^{2n-2}$ com $a>0$ suficientemente grande  {\color{3}(como para conter a projeção da bola merulhada no cilindro)}.

Portanto, a teorema de non-squeezing segue do seguinte resultado:

\begin{thm}[non-squeezing para produtos]\leavevmode
Seja $(N^{2n-2},\tau)$ variedade simplética compacta tal que $\pi_{2}(N)=0$. Se existe um mergulho simplético $B^{2n}(r)\hookrightarrow B^2(R)\times N$ então $r<R$.
\end{thm}

Note que o teorema de non-squeezing the Gromov segue do teorema non-squeezing para produtos porque o $\pi_{2}(\text{toro} )$ é zero pelo seu recobrimento universal.

\begin{question}\leavevmode
	Por que foi que o teorema non-squeezing foi reduzido para o teorema non-squeezing para produtos?
\end{question}

\begin{thing4}{Resposta}\leavevmode
Considera o mergulho a bola no cilintro, considera uma projeção a $\mathbb{C}^{n-1}$, obtendo um conjunto compacto em $\mathbb{C}^{n-1}$. Bota isso dentro de um dominío fundamental e quocienta para obter um toro. (Lembre que o $\pi_{2}(\text{toro} )=0$.) Mas lembre que você projetou, então tem que pegar um produto,, i.e. continuar o mergulho da bola que tava mergulada no cilindo mergulhando isso num produto de bola por toro.
\end{thing4}

\section{Aula 28}

\begin{thm}[Gromov]\leavevmode
	Seja $(N,\tau)$ variedade simplética compacta tal que $\pi_{2}(N)=0$ e $J$ estrutura complexa em $S^2 \times N$ compatível com a forma simplética produto $\sigma \oplus  \tau$, onde $\sigma$ é a forma de área de $S^2$. Então dado $p \in S^2 \times N$, existe uma curva $J$-holomorfa $u:S^2\to S^2 \times N$ passando por $p$ e homológa à fibra $S^2 \times \{x\}$.
\end{thm}

\subsection{Sketch da prova do teorema non-squeezing usando o teorema de Gromov}
Se existe um mergulho simplético
\[B^{2n}(r) \hookrightarrow B^2(R)\times N\]
então existe
\[B^{2n}(r) \hookrightarrow  S^2(R/2) \times N.\]
(Por que?) Seja $J_0$ a estrutura quase-complexa canônica em $\mathbb{R}^{2n}$ e considere a equação $\varphi_*J_0$ em $\varphi(B^{2n}(r))$.

\begin{prop}\leavevmode
	Dada $(M,\omega)$ uma subvariedade $S \subset M$ e uma estrutura quase-complexa $J_0$ definida ao longo de $S$ (i.e. um endomorfismo em  $TM|_{S}$ tal que $J_0^2=-\operatorname{Id}$---definido em todo $TM$ mas  só deve funcionar em $S$), compatível com $\omega$. Então existe uma estrtura quase-complexa canônica $J$ em $M$ compatível com $\omega$ que extende a $J_0$, i.e. $J|_{S}=J_0$
\end{prop}

\begin{proof}\leavevmode
Exercício.
\end{proof}

Aplicamos essa proposção à subvariedade \[\varphi:((B^{2n}(r)),J_0)\hookrightarrow  (S^2(R/2)\times N,J),\]
obtendo assim que $\varphi$ não é so um simplectomorfismo mas também é holomorfa, onde $J$ é a extensão de $\varphi_*J_0$.

Agora aplicamos o teorema do Gromov com respeito a esa estrutura quase-complexa: dado qualquer ponto existe uma curva $J$-holomorfa $\psi$ que \textit{é homológa à fibra}  $S^2\times\{x\}$. Pegue a origem como ponto distinguido, assim a curva $J$-holomorfa passa por $\varphi(0) \in \operatorname{img}\psi$. Por Stokes, a integral depende só da classe da homologia, i.e.
\[\int_{S^2}\phi^* (\sigma \oplus \tau)=\int_{S^2\times \{x\}}\sigma \oplus  \tau=\pi \mathbb{R}^2\]
porque, lembre, passamos nosso mergulho para $S^2(R/2)\times N$.

Decrescendo $r$ um pouco se necessário, podemos assumir que $\psi:S^2 \to S^2\times N$ é transversal a $\varphi(\overline{B^{2n}(r))}$, exercício. Portanto, $S':=\psi^{-1}\Big(\overline{\varphi(B^{2n}(r))}\Big)$ é uma 2-variedade com bordo. {\color{3}É só olhara a pre imagem da bola que é dominio da $\varphi$ no domínio da $\psi$.}

\begin{remark}[Por que tem que ter bordo?]\leavevmode
	Note que a imagem de $\psi$ não pode estar contida na imagem de $\varphi$ porque a forma simplética em $\varphi$ é zero porque a segunda cohomologia da esfera é zero, assim toda forma fechada é exata, porém a gente viu que a área da imagem de  $\psi$ não é zero.
\end{remark}

\begin{lemma}[de monotonicidade]\leavevmode
Seja $f: S' \to B^{2n}(r)$ uma curva $J_0$-holomorfa tal que $f(\partial S^1)\subset \partial B^{2n}(r)$ e $0 \in \operatorname{img} f$, então
\[\int_{S'}f^*\omega_0\geq \pi r^2.\]
\end{lemma}

\begin{proof}\leavevmode
Isso usa o fato de que superfícies holomorfas sempre são mínimas. De fato, pela construção a preimagem da $\psi$ no domínio da $\varphi$ (na verdade estamos sim no domínio da $\varphi$: a gente toma $f=\varphi^{-1}\circ\psi|_{S'}$) fica uma curva holomorfa com bordo. E a área da superfície mínima com esse bordo é como mínimo $\pi r^2$.
\end{proof}

Enão temos

\begin{align*}
\int_{S'}f^*\psi&=\int_{S'}(\varphi^{-1} \circ \psi)^*  \omega_0=\int_{S'}\psi^*\circ(\varphi^{-1})^*\omega_0\\
&=\int_{S'}\psi^*(\sigma  \oplus \tau)
\end{align*}
já que $\varphi^{-1}$ é simplética---a forma simplética na esfera é $\omega_0$ é no produto $\sigma \oplus  \tau$.

Consequentemente
 \begin{align*}
	 \pi R^2&=\int_{S^2}\varphi^* (\sigma \oplus  \tau) \geq \int_{S'}\psi^*(\sigma \oplus  \tau) \geq \pi r^2.
	 &=
\end{align*}
Isso conclui o nosso sketch da prova do teorema do non-squeezing. A caixa preta é a existência da curva holomorfa, o teorema devido a Gromov.

\end{document}
